\documentclass[10pt,a4paper,landscape]{article}
\usepackage[landscape,margin=1.5cm]{geometry}
\usepackage[ngerman]{babel}
\usepackage{fontspec}
\usepackage{longtable}
\usepackage{array}
\usepackage{booktabs}
\usepackage[table]{xcolor}
\usepackage{ragged2e}

\setlength{\LTpre}{0pt}
\setlength{\LTpost}{0pt}
\setlength{\tabcolsep}{4pt}
\renewcommand{\arraystretch}{1.2}

\newcolumntype{L}[1]{>{\RaggedRight\arraybackslash}p{#1}}

\begin{document}
\footnotesize

\begin{center}
{\Large\bfseries Synopse IKJHG - Vergleich der Referentenentwürfe 2024 und 2026}
\end{center}
\vspace{0.5cm}

\begin{longtable}{|L{8.6cm}|L{8.6cm}|L{8.6cm}|}
\hline
\multicolumn{3}{|c|}{\cellcolor{gray!20}\textbf{Artikel 1 (2026): Änderung des Achten Buches Sozialgesetzbuch}} \\
\hline
\textbf{Geltendes Recht} & \textbf{Änderungen RefE 2024} & \textbf{Änderungen RefE 2026} \\
\hline
\endhead

\multicolumn{3}{|c|}{\cellcolor{gray!10}\textbf{\S{} 1 -- Recht auf Erziehung, Elternverantwortung, Jugendhilfe}} \\
\hline
(1) Jeder junge Mensch hat ein Recht auf \newline Förderung seiner Entwicklung und auf \newline Erziehung zu einer selbstbestimmten, \newline eigenverantwortlichen und \newline gemeinschaftsfähigen Persönlichkeit & \textit{unverändert} & (1) Jeder junge Mensch hat ein Recht auf \newline Förderung seiner Entwicklung und auf \newline Erziehung zu einer selbstbestimmten, \newline eigenverantwortlichen und \newline gemeinschaftsfähigen Persönlichkeit \textbf{und \newline auf Förderung seiner vollen, wirksamen \newline und gleichberechtigten Teilhabe am \newline Leben in der Gesellschaft } \\
\hline
(2) Pflege und Erziehung der Kinder sind \newline das natürliche Recht der Eltern und die \newline zuvörderst ihnen obliegende Pflicht. Über \newline ihre Betätigung wacht die staatliche \newline Gemeinschaft. & \textit{unverändert} & \textit{(2) unverändert} \\
\hline
(3) Jugendhilfe soll zur Verwirklichung des \newline Rechts nach Absatz 1 insbesondere \newline ihrer individuellen Fähigkeiten in allen sie \newline betreffenden Lebensbereichen \newline selbstbestimmt zu interagieren und damit \newline gleichberechtigt am Leben in der \newline Gesellschaft teilhaben zu können, & \textit{unverändert} & \textit{(3) unverändert} \\
\hline
1. junge Menschen in ihrer individuellen \newline und sozialen Entwicklung fördern und dazu \newline beitragen, Benachteiligungen zu vermeiden \newline oder abzubauen, & \textit{unverändert} & 1. unverändert \\
\hline
2. jungen Menschen ermöglichen oder \newline erleichtern, entsprechend ihrem Alter und & \textit{unverändert} & 2. unverändert \\
\hline
3. Eltern und andere Erziehungsberechtigte \newline bei der Erziehung beraten und \newline unterstützen, & \textit{unverändert} & 3. unverändert \\
\hline
4. Kinder und Jugendliche vor Gefahren für \newline ihr Wohl schützen & \textit{unverändert} & 4. unverändert \\
\hline
5. dazu beitragen, positive \newline Lebensbedingungen für junge Menschen \newline und ihre Familien sowie eine kinder- und \newline familienfreundliche Umwelt zu erhalten \newline oder zu schaffen & \textit{unverändert} & 5. unverändert \\
\hline
\multicolumn{3}{|c|}{\cellcolor{gray!10}\textbf{\S{} 2 -- Aufgaben der Jugendhilfe}} \\
\hline
(1) Die Jugendhilfe umfasst Leistungen und \newline andere Aufgaben zugunsten junger \newline Menschen und Familien. & \textit{(1) unverändert} & \textit{(1) unverändert} \\
\hline
(2) Leistungen der Jugendhilfe sind: & (2) Leistungen der Jugendhilfe sind: \newline \textbf{a) \newline Hilfe zur Erziehung und ergänzende \newline Leistungen (\S{}\S{} 27 Absatz 2, 27a, 28 bis \newline 35, 36 bis 37, 39 bis 40),} \newline \textbf{b) Leistungen der Eingliederungshilfe \newline für Kinder und Jugendliche mit \newline Behinderungen und ergänzende \newline Leistungen (\S{}\S{} 27 Absatz 3 bis 3b, 35a \newline bis 40), } & (2) Leistungen der Jugendhilfe sind: \newline \textbf{a) \newline Hilfe zur Erziehung und ergänzende \newline Leistungen (\S{}\S{} 27 Absatz 2, 27a, 28 bis \newline 35, 36 bis 37, 39 bis 40),} \newline \textbf{b) Leistungen der Eingliederungshilfe \newline für Kinder und Jugendliche mit \newline Behinderungen und ergänzende \newline Leistungen (\S{}\S{} 27 Absatz 3 und 4, 35a \newline bis 40), } \\
\hline
1. \newline Angebote der Jugendarbeit, der \newline Jugendsozialarbeit, der Schulsozialarbeit \newline und des erzieherischen Kinder- und \newline Jugendschutzes (\S{}\S{} 11 bis 14), & 1. \newline Angebote der Jugendarbeit, der \newline Jugendsozialarbeit, der Schulsozialarbeit \newline und des erzieherischen Kinder- und \newline Jugendschutzes (\S{}\S{} 11 bis 14), & \textbf{2}. \newline Angebote der Jugendarbeit, der \newline Jugendsozialarbeit, der Schulsozialarbeit \newline und des erzieherischen Kinder- und \newline Jugendschutzes (\S{}\S{} 11 bis 14), \\
\hline
2. \newline Angebote zur Förderung der Erziehung in \newline der Familie (\S{}\S{} 16 bis 21), & 2. \newline Angebote zur Förderung der Erziehung in \newline der Familie (\S{}\S{} 16 bis 21), & \textbf{3}. \newline Angebote zur Förderung der Erziehung in \newline der Familie (\S{}\S{} 16 bis 21), \\
\hline
3. \newline Angebote zur Förderung von Kindern in \newline Tageseinrichtungen und in \newline Kindertagespflege (\S{}\S{} 22 bis 25), & 3. \newline Angebote zur Förderung von Kindern in \newline Tageseinrichtungen und in \newline Kindertagespflege (\S{}\S{} 22 bis 25), & \textbf{4}. \newline Angebote zur Förderung von Kindern in \newline Tageseinrichtungen und in \newline Kindertagespflege (\S{}\S{} 22 bis 25), \\
\hline
4. \newline Hilfe zur Erziehung und ergänzende \newline Leistungen (\S{}\S{} 27 bis 35, 36, 37, 39, 40), & 4. \newline Hilfe zur Erziehung und ergänzende \newline Leistungen \textbf{(\S{}\S{} 27 bis 40),} & \textbf{5. \newline Leistungen zur Entwicklung, zur \newline Erziehung und zur Teilhabe} \textbf{(\S{}\S{} 27 bis \newline 40),} \\
\hline
5. \newline Hilfe für seelisch behinderte Kinder und \newline Jugendliche und ergänzende Leistungen \newline (\S{}\S{} 35a bis 37, 39, 40), & 5. \newline Hilfe für seelisch behinderte Kinder und \newline Jugendliche und ergänzende Leistungen \newline (\S{}\S{} 35a bis 37, 39, 40), & 5. \newline Hilfe für seelisch behinderte Kinder und \newline Jugendliche und ergänzende Leistungen \newline (\S{}\S{} 35a bis 37, 39, 40), \\
\hline
6. \newline Hilfe für junge Volljährige und \newline Nachbetreuung (den \S{}\S{} 41 und 41a). & \textbf{5}. \newline Hilfe für junge Volljährige und \newline Nachbetreuung (den \S{}\S{} 41 und 41a). & \textbf{6}. \newline Hilfe für junge Volljährige und \newline Nachbetreuung (den \S{}\S{} 41 und 41a). \\
\hline
\textit{--- Fassung 2024 ---} \newline (3) Andere Aufgaben der Jugendhilfe sind \newline \textit{--- Fassung 2026 ---} \newline (3) Andere Aufgaben der Jugendhilfe sind \newline einer Einrichtung sowie die Erteilung \newline nachträglicher Auflagen und die damit \newline verbundenen Aufgaben (\S{}\S{} 45 bis 47, 48a), & (3) Andere Aufgaben der Jugendhilfe sind & (3) Andere Aufgaben der Jugendhilfe sind \\
\hline
1. \newline die Inobhutnahme von Kindern und \newline Jugendlichen (\S{} 42), & 1. unverändert & 1. unverändert \\
\hline
2. \newline die vorläufige Inobhutnahme von \newline ausländischen Kindern und Jugendlichen \newline nach unbegleiteter Einreise (\S{} 42a), & 2. unverändert & 2. unverändert \\
\hline
3. \newline die Erteilung, der Widerruf und die \newline Zurücknahme der Pflegeerlaubnis (\S{}\S{} 43, \newline 44), & 3. unverändert & 3. unverändert \\
\hline
\textit{--- Fassung 2024 ---} \newline 4. \newline die Erteilung, der Widerruf und die \newline Zurücknahme der Erlaubnis für den Betrieb \newline einer Einrichtung sowie die Erteilung \newline nachträglicher Auflagen und die damit \newline verbundenen Aufgaben (\S{}\S{} 45 bis 47, 48a), \newline \textit{--- Fassung 2026 ---} \newline 4. \newline die Erteilung, der Widerruf und die \newline Zurücknahme der Erlaubnis für den Betrieb & 4. unverändert & 4. unverändert \\
\hline
5. \newline die Tätigkeitsuntersagung (\S{}\S{} 48, 48a), & 5. unverändert & 5. unverändert \\
\hline
6. \newline die Mitwirkung in Verfahren vor den \newline Familiengerichten (\S{} 50), & 6. unverändert & 6. unverändert \\
\hline
7. \newline die Beratung und Belehrung in Verfahren \newline zur Annahme als Kind (\S{} 51), & 7. unverändert & 7. unverändert \\
\hline
8. \newline die Mitwirkung in Verfahren nach dem \newline Jugendgerichtsgesetz (\S{} 52), & 8. unverändert & 8. unverändert \\
\hline
9. \newline die Beratung und Unterstützung von \newline Müttern bei Vaterschaftsfeststellung und \newline Geltendmachung von \newline Unterhaltsansprüchen sowie von Pflegern \newline und Vormündern (\S{}\S{} 52a, 53a), & 9. unverändert & 9. unverändert \\
\hline
10. \newline die Erteilung, der Widerruf und die \newline Zurücknahme der Anerkennung als \newline Vormundschaftsverein (\S{} 54) & 10. unverändert & 10. unverändert \\
\hline
11. \newline Beistandschaft, Pflegschaft und \newline Vormundschaft des Jugendamts (\S{}\S{} 55 bis \newline 57), & 11. unverändert & 11. unverändert \\
\hline
12. \newline Beurkundung (\S{} 59), & 12. unverändert & 12. unverändert \\
\hline
13. \newline die Aufnahme von vollstreckbaren \newline Urkunden (\S{} 60) & 13. unverändert & 13. unverändert \\
\hline
\multicolumn{3}{|c|}{\cellcolor{gray!10}\textbf{\S{} 3 -- Freie und öffentliche Jugendhilfe}} \\
\hline
(1) Die Jugendhilfe ist gekennzeichnet \newline durch die Vielfalt von Trägern \newline unterschiedlicher Wertorientierungen und \newline die Vielfalt von Inhalten, Methoden und \newline Arbeitsformen. & \textit{unverändert} & \textit{(1) unverändert} \\
\hline
(2) Leistungen der Jugendhilfe werden von \newline Trägern der freien Jugendhilfe und von \newline Trägern der öffentlichen Jugendhilfe \newline erbracht. Leistungsverpflichtungen, die \newline durch dieses Buch begründet werden, \newline richten sich an die Träger der öffentlichen \newline Jugendhilfe. & \textit{unverändert} & \textit{(2) unverändert} \\
\hline
(3) Andere Aufgaben der Jugendhilfe \newline werden von Trägern der öffentlichen \newline Jugendhilfe wahrgenommen. Soweit dies \newline ausdrücklich bestimmt ist, können Träger \newline der freien Jugendhilfe diese Aufgaben \newline wahrnehmen oder mit ihrer Ausführung \newline betraut werden. & \textit{unverändert} & \textit{(3) unverändert} \\
\hline
 & \textit{unverändert} & \textbf{(4) Die Träger der öffentlichen \newline Jugendhilfe sind für Leistungen nach \S{} \newline 2 Absatz 2 Nummer 4 Buchstabe b \newline Rehabilitationsträger im Sinne von \S{} 6 \newline Absatz 1 Nummer 6 des Neunten \newline Buches. \S{} 7 des Neunten Buches ist zu \newline beachten. } \\
\hline
\multicolumn{3}{|c|}{\cellcolor{gray!10}\textbf{\S{} 5 -- Wunsch und Wahlrecht}} \\
\hline
(1) Die Leistungsberechtigten haben das \newline Recht, zwischen Einrichtungen und \newline Diensten verschiedener Träger zu wählen \newline und Wünsche hinsichtlich der Gestaltung \newline der Hilfe zu äußern. Sie sind auf dieses \newline Recht hinzuweisen. & \textit{(1) unverändert} & \textit{(1) unverändert} \\
\hline
(2) Der Wahl und den Wünschen soll \newline entsprochen werden, sofern dies nicht mit \newline unverhältnismäßigen Mehrkosten \newline verbunden ist. Wünscht der \newline Leistungsberechtigte die Erbringung einer \newline in \S{} 78a genannten Leistung in einer \newline Einrichtung, mit deren Träger keine \newline Vereinbarungen nach \S{} 78b bestehen, so \newline soll der Wahl nur entsprochen werden, \newline wenn die Erbringung der Leistung in dieser \newline Einrichtung im Einzelfall oder nach \newline Maßgabe des Hilfeplans (\S{} 36) geboten ist. & \textit{(2) unverändert} & \textit{(2) unverändert} \\
\hline
 & \textbf{(3) Bei der Entscheidung nach Absatz 2 \newline ist zunächst die Zumutbarkeit einer \newline Abweichung von den Wünschen des \newline Leistungsberechtigten zu prüfen. Dabei \newline sin die persönlichen, familiären und \newline örtlichen Umstände einschließlich der \newline gewünschten Wohnform angemessen \newline zu berücksichtigen. Bei Unzumutbarkeit \newline einer abweichenden \newline Leistungsgestaltung ist ein \newline Kostenvergleich nicht vorzunehmen. \newline Für Leistungsberechtigte nach \S{} 27 \newline Absatz 3 gilt im Übrigen \S{} 104 Absatz 3 \newline Satz 3 und 4 des Neunten Buches \newline entsprechend. } & \textbf{(3) Eine von den Wünschen des \newline Leistungsberechtigten abweichende \newline Leistung kommt nicht in Betracht, wenn \newline diese Abweichung für den \newline Leistungsberechtigten unzumutbar ist. \newline Bei der Prüfung der Zumutbarkeit sind \newline die persönlichen, familiären und \newline örtlichen Umstände einschließlich der \newline gewünschten Wohnform angemessen \newline zu berücksichtigen. Bei Unzumutbarkeit \newline einer abweichenden \newline Leistungsgestaltung ist ein \newline Kostenvergleich nicht vorzunehmen. \newline Für Leistungsberechtigte nach \S{} 27 \newline Absatz 3 gilt im Übrigen \S{} 104 Absatz 4 \newline des Neunten Buches entsprechend. } \\
\hline
\multicolumn{3}{|c|}{\cellcolor{gray!10}\textbf{\S{} 8a -- Schutzauftrag bei Kindeswohlgefährdung}} \\
\hline
(1) Werden dem Jugendamt gewichtige \newline Anhaltspunkte für die Gefährdung des \newline Wohls eines Kindes oder Jugendlichen \newline bekannt, so hat es das Gefährdungsrisiko \newline im Zusammenwirken mehrerer Fachkräfte \newline einzuschätzen. Soweit der wirksame \newline Schutz dieses Kindes oder dieses \newline Jugendlichen nicht in Frage gestellt wird, \newline hat das Jugendamt die \newline Erziehungsberechtigten sowie das Kind \newline oder den Jugendlichen in die \newline Gefährdungseinschätzung einzubeziehen \newline und, sofern dies nach fachlicher \newline Einschätzung erforderlich ist, \newline 1. sich dabei einen unmittelbaren Eindruck \newline von dem Kind und von seiner persönlichen \newline Umgebung zu verschaffen sowie \newline 2. Personen, die gemäß \S{} 4 Absatz 3 des \newline Gesetzes zur Kooperation und Information \newline im Kinderschutz dem Jugendamt Daten \newline übermittelt haben, in geeigneter Weise an & \textit{unverändert} & (1) Werden dem Jugendamt gewichtige \newline Anhaltspunkte für die Gefährdung des \newline Wohls eines Kindes oder Jugendlichen \newline bekannt, so hat es das Gefährdungsrisiko \newline im Zusammenwirken mehrerer Fachkräfte \newline \textbf{und, soweit dies notwendig ist, unter \newline Beteiligung betroffener Einrichtungen, \newline Dienste oder anderer Stellen} \newline einzuschätzen. Soweit der wirksame \newline Schutz dieses Kindes oder dieses \newline Jugendlichen nicht in Frage gestellt wird, \newline hat das Jugendamt die \newline Erziehungsberechtigten sowie das Kind \newline oder den Jugendlichen in die \newline Gefährdungseinschätzung einzubeziehen \newline und, sofern dies nach fachlicher \newline Einschätzung erforderlich ist, \newline 1. sich dabei einen unmittelbaren Eindruck \newline von dem Kind und von seiner persönlichen \newline Umgebung zu verschaffen sowie \newline 2. Personen, die gemäß \S{} 4 Absatz 3 des \newline Gesetzes zur Kooperation und Information \newline im Kinderschutz dem Jugendamt Daten \\
\hline
Absätze 2 bis 6 & \textit{unverändert} & \textit{unverändert} \\
\hline
\multicolumn{3}{|c|}{\cellcolor{gray!10}\textbf{\S{} 10 -- Verhältnis zu anderen Leistungen und Verpflichtungen}} \\
\hline
(1) Verpflichtungen anderer, insbesondere \newline der Träger anderer Sozialleistungen und \newline der Schulen, werden durch dieses Buch \newline nicht berührt. Auf Rechtsvorschriften \newline beruhende Leistungen anderer dürfen nicht \newline deshalb versagt werden, weil nach diesem \newline Buch entsprechende Leistungen \newline vorgesehen sind. & \textit{(1) unverändert} & \textit{(1) unverändert} \\
\hline
(2) Unterhaltspflichtige Personen werden \newline nach Maßgabe der \S{}\S{} 90 bis 97b an den \newline Kosten für Leistungen und vorläufige \newline Maßnahmen nach diesem Buch beteiligt. \newline Soweit die Zahlung des Kostenbeitrags die \newline Leistungsfähigkeit des Unterhaltspflichtigen \newline mindert oder der Bedarf des jungen \newline Menschen durch Leistungen und vorläufige \newline Maßnahmen nach diesem Buch gedeckt \newline ist, ist dies bei der Berechnung des \newline Unterhalts zu berücksichtigen. & \textit{(2) unverändert} & \textit{(2) unverändert} \\
\hline
(3) Die Leistungen nach diesem Buch \newline gehen Leistungen nach dem Zweiten Buch \newline vor. Abweichend von Satz 1 gehen \newline Leistungen nach \S{} 3 Absatz 2, den \S{}\S{} 14 \newline bis 16g, 16k, \S{} 19 Absatz 2 in Verbindung \newline mit \S{} 28 Absatz 6 des Zweiten Buches \newline sowie Leistungen nach \S{} 6b Absatz 2 des \newline Bundeskindergeldgesetzes in Verbindung \newline mit \S{} 28 Absatz 6 des Zweiten Buches den \newline Leistungen nach diesem Buch vor. & \textit{(3) unverändert} & \textit{(3) unverändert} \\
\hline
(4) Die Leistungen nach diesem Buch \newline gehen Leistungen nach dem Neunten und \newline Zwölften Buch vor. Abweichend von Satz 1 \newline gehen Leistungen nach \S{} 27a Absatz 1 in \newline Verbindung mit \S{} 34 Absatz 6 des Zwölften \newline Buches und Leistungen der \newline Eingliederungshilfe nach dem Neunten \newline Buch für junge Menschen, die körperlich \newline oder geistig behindert oder von einer \newline solchen Behinderung bedroht sind, den \newline Leistungen nach diesem Buch vor. \newline Landesrecht kann regeln, dass Leistungen \newline der Frühförderung für Kinder unabhängig \newline von der Art der Behinderung vorrangig von \newline anderen Leistungsträgern gewährt werden. & \textbf{(4) Die Leistungen nach diesem Buch \newline gehen Leistungen nach dem Neunten \newline und Zwölften Buch vor. } & \textbf{(4) Die Leistungen nach diesem Buch \newline gehen Leistungen nach dem Neunten \newline und Zwölften Buch vor. } \\
\hline
(5) Soweit Leistungen zum Lebensunterhalt \newline nach \S{} 39 erbracht werden, gehen sie den \newline Leistungen zum Lebensunterhalt nach \S{} 93 \newline des Vierzehnten Buches vor. & \textbf{(5) Die Leistungen nach diesem Buch \newline gehen Leistungen nach dem Zwölften \newline Buch vor. Abweichend von Satz 1 gehen \newline Leistungen nach \S{} 27a Absatz 1 in \newline Verbindung mit \S{} 34 Absatz 6 des \newline Zwölften Buches den Leistungen nach \newline diesem Buch vor.1 } & \textbf{(5) Die Leistungen nach diesem Buch \newline gehen Leistungen nach dem Zwölften \newline Buch vor. Abweichend von Satz 1 gehen \newline Leistungen für den Lebensunterhalt \newline 1. nach dem Dritten Kapitel des Zwölften \newline Buches zur Deckung der Bedarfe nach \newline dem Ersten, Zweiten, Vierten und \newline Fünften Abschnitt, von Leistungen zur \newline Deckung von Bedarfen des Dritten \newline Abschnittes nur diejenigen nach \S{} 34 \newline Absatz 6 des Zwölften Buches und \newline 2. nach dem Vierten Kapitel des \newline Zwölften Buches \newline den Leistungen nach diesem Buch vor. } \\
\hline
 & \textbf{(6)2 Soweit Leistungen zum \newline Lebensunterhalt nach \S{} 39 erbracht \newline werden, gehen sie den Leistungen zum \newline Lebensunterhalt nach \S{} 93 des \newline Vierzehnten Buches vor. } & \textit{unverändert} \\
\hline
\multicolumn{3}{|c|}{\cellcolor{gray!10}\textbf{\S{} 10a -- Beratung}} \\
\hline
Absatz 3 & \textit{unverändert} & \textit{unverändert} \\
\hline
(2) Die Beratung umfasst insbesondere \newline 1. die Familiensituation oder die \newline persönliche Situation des jungen \newline Menschen, Bedarfe, vorhandene \newline Ressourcen sowie mögliche Hilfen, \newline 2. die Leistungen der Kinder- und \newline Jugendhilfe einschließlich des Zugangs \newline zum Leistungssystem, \newline 3. die Leistungen anderer Leistungsträger, \newline 4. mögliche Auswirkungen und Folgen \newline einer Hilfe, \newline 5. die Verwaltungsabläufe, \newline 6. Hinweise auf Leistungsanbieter und \newline andere Hilfemöglichkeiten im Sozialraum & \textit{unverändert} & (2) Die Beratung umfasst insbesondere \newline 1. die Familiensituation oder die \newline persönliche Situation des jungen \newline Menschen, Bedarfe, vorhandene \newline Ressourcen sowie mögliche Hilfen, \newline 2. die Leistungen der Kinder- und \newline Jugendhilfe einschließlich des Zugangs \newline zum Leistungssystem, \newline 3. die Leistungen anderer Leistungsträger, \newline 4. mögliche Auswirkungen und Folgen \newline einer Hilfe, \newline 5. die Verwaltungsabläufe, \newline 6. Hinweise auf Leistungsanbieter und \newline andere Hilfemöglichkeiten im Sozialraum \\
\hline
\multicolumn{3}{|c|}{\cellcolor{gray!10}\textbf{\S{} 10b -- Verfahrenslotse}} \\
\hline
(1) Junge Menschen, die Leistungen der \newline Eingliederungshilfe wegen einer \newline Behinderung oder wegen einer drohenden \newline Behinderung geltend machen oder bei \newline denen solche Leistungsansprüche in \newline Betracht kommen, sowie ihre Mütter, Väter, \newline Personensorge- und \newline Erziehungsberechtigten haben bei der \newline Antragstellung, Verfolgung und \newline Wahrnehmung dieser Leistungen Anspruch \newline auf Unterstützung und Begleitung durch \newline einen Verfahrenslotsen. Der \newline Verfahrenslotse soll die \newline Leistungsberechtigten bei der \newline Verwirklichung von Ansprüchen auf \newline Leistungen der Eingliederungshilfe \newline unabhängig unterstützen sowie auf die \newline Inanspruchnahme von Rechten hinwirken. \newline Diese Leistung wird durch den örtlichen \newline Träger der öffentlichen Jugendhilfe \newline erbracht. & (1) Junge Menschen, die Leistungen \textbf{zur \newline Teilhabe }wegen einer Behinderung oder \newline wegen einer drohenden Behinderung \newline geltend machen oder bei denen solche \newline Leistungsansprüche in Betracht kommen, \newline sowie ihre Mütter, Väter, Personensorge- \newline und Erziehungsberechtigten haben bei der \newline Antragstellung, Verfolgung und \newline Wahrnehmung dieser Leistungen Anspruch \newline auf Unterstützung und Begleitung durch \newline einen Verfahrenslotsen. Der \newline Verfahrenslotse soll die \newline Leistungsberechtigten bei der \newline Verwirklichung von Ansprüchen auf \newline Leistungen \textbf{zur Teilhabe} unabhängig \newline unterstützen sowie auf die \newline Inanspruchnahme von Rechten hinwirken. \newline Diese Leistung wird durch den örtlichen \newline Träger der öffentlichen Jugendhilfe \newline erbracht. & (1) Junge Menschen, die Leistungen \textbf{zur \newline Teilhabe }wegen einer Behinderung oder \newline wegen einer drohenden Behinderung \newline geltend machen oder bei denen solche \newline Leistungsansprüche in Betracht kommen, \newline sowie ihre Mütter, Väter, Personensorge- \newline und Erziehungsberechtigten haben bei der \newline Antragstellung, Verfolgung und \newline Wahrnehmung dieser Leistungen Anspruch \newline auf Unterstützung und Begleitung durch \newline einen Verfahrenslotsen. Der \newline Verfahrenslotse soll die \newline Leistungsberechtigten bei der \newline Verwirklichung von Ansprüchen auf \newline Leistungen \textbf{zur Teilhabe} unterstützen \newline sowie auf die Inanspruchnahme von \newline Rechten hinwirken. \textbf{Der Verfahrenslotse \newline soll auf Wunsch der in Satz 1 genannten \newline Person auch zu Ansprüchen im Rahmen \newline der Pflegeversicherung und deren \newline Inanspruchnahme beraten und die \newline Anspruchsberechtigten unterstützen. \newline Die Leistung nach Satz 1 und 2 wird \newline durch den örtlichen Träger der \newline öffentlichen Jugendhilfe funktionell, \newline organisatorisch und personell getrennt \newline von seinen übrigen Aufgaben erbracht.} \\
\hline
\textit{--- Fassung 2024 ---} \newline (2) Der Verfahrenslotse unterstützt den \newline örtlichen Träger der öffentlichen \newline Jugendhilfe bei der Zusammenführung der \newline Leistungen der Eingliederungshilfe für \newline junge Menschen in dessen Zuständigkeit. \newline Hierzu berichtet er gegenüber dem \newline örtlichen Träger der öffentlichen \newline Jugendhilfe halbjährlich insbesondere über \newline Erfahrungen der strukturellen \newline Zusammenarbeit mit anderen Stellen und \newline öffentlichen Einrichtungen, insbesondere \newline mit anderen Rehabilitationsträgern. \newline \textit{--- Fassung 2026 ---} \newline (2) Der Verfahrenslotse unterstützt den \newline örtlichen Träger der öffentlichen \newline Jugendhilfe bei der Zusammenführung der \newline Leistungen der Eingliederungshilfe für & (2) Der Verfahrenslotse unterstützt den \newline örtlichen Träger der öffentlichen \newline Jugendhilfe bei der \textbf{inklusiven \newline Wahrnehmung der Aufgaben nach \newline diesem Buch insbesondere im Rahmen \newline der Jugendhilfeplanung nach \S{} 80.} & (2) Der Verfahrenslotse unterstützt den \newline örtlichen Träger der öffentlichen \newline Jugendhilfe bei der \textbf{inklusiven \newline Wahrnehmung der Aufgaben nach } \\
\hline
\textit{--- Fassung 2024 ---} \newline Unterabschnitt 1: \newline Hilfe zur Erziehung \newline \textit{--- Fassung 2026 ---} \newline junge Menschen in dessen Zuständigkeit. \newline Hierzu berichtet er gegenüber dem \newline örtlichen Träger der öffentlichen \newline Jugendhilfe halbjährlich insbesondere über \newline Erfahrungen der strukturellen \newline Zusammenarbeit mit anderen Stellen und \newline öffentlichen Einrichtungen, insbesondere \newline mit anderen Rehabilitationsträgern. & Unterabschnitt 1:\textbf{ \newline Leistungen zur Entwicklung, zur \newline Erziehung und zur Teilhabe } & \textbf{diesem Buch insbesondere im Rahmen \newline der Jugendhilfeplanung nach \S{} 80.} \\
\hline
\multicolumn{3}{|c|}{\cellcolor{gray!10}\textbf{\S{} 13 -- Jugendsozialarbeit}} \\
\hline
Absatz 4 & \textit{unverändert} & \textit{unverändert} \\
\hline
(3) Jungen Menschen kann während der \newline Teilnahme an schulischen oder beruflichen \newline Bildungsmaßnahmen oder bei der \newline beruflichen Eingliederung Unterkunft in \newline sozialpädagogisch begleiteten \newline Wohnformen angeboten werden. In diesen \newline Fällen sollen auch der notwendige \newline Unterhalt des jungen Menschen \newline sichergestellt und Krankenhilfe nach \newline Maßgabe des \S{} 40 geleistet werden & \textit{unverändert} & (3) Jungen Menschen kann \textbf{in Ergänzung \newline oder unabhängig von Hilfen oder \newline Maßnahmen nach Absatz 1 oder 2}während der Teilnahme an schulischen \newline oder beruflichen Bildungsmaßnahmen oder \newline bei der beruflichen Eingliederung \newline Unterkunft in sozialpädagogisch \newline begleiteten Wohnformen angeboten \newline werden. In diesen Fällen sollen auch der \newline notwendige Unterhalt des jungen \newline Menschen sichergestellt und Krankenhilfe \newline nach Maßgabe des \S{} 40 geleistet werden \\
\hline
\multicolumn{3}{|c|}{\cellcolor{gray!10}\textbf{\S{} 18 -- Beratung und Unterstützung bei der Ausübung der Personensorge und des Umgangsrechts}} \\
\hline
Unterabschnitt 1: \newline Hilfe zur Erziehung & \textit{unverändert} & Unterabschnitt 1:\textbf{ \newline Leistungen zur Entwicklung, zur \newline Erziehung und zur Teilhabe } \\
\hline
(3) Kinder und Jugendliche haben \newline Anspruch auf Beratung und Unterstützung \newline bei der Ausübung des Umgangsrechts \newline nach \S{} 1684 Absatz 1 des Bürgerlichen \newline Gesetzbuchs. Sie sollen darin unterstützt \newline werden, dass die Personen, die nach \newline Maßgabe der \S{}\S{} 1684, 1685 und 1686a \newline des Bürgerlichen Gesetzbuchs zum \newline Umgang mit ihnen berechtigt sind, von \newline diesem Recht zu ihrem Wohl Gebrauch \newline machen. Eltern, andere \newline Umgangsberechtigte sowie Personen, in \newline deren Obhut sich das Kind befindet, haben \newline Anspruch auf Beratung und Unterstützung \newline bei der Ausübung des Umgangsrechts. Bei & \textit{unverändert} & (3) Kinder und Jugendliche haben \newline Anspruch auf Beratung und Unterstützung \newline bei der Ausübung des Umgangsrechts \newline nach \S{} 1684 Absatz 1 des Bürgerlichen \newline Gesetzbuchs. Sie sollen darin unterstützt \newline werden, dass die Personen, die nach \newline Maßgabe der \S{}\S{} 1684, 1685 und 1686a \newline des Bürgerlichen Gesetzbuchs zum \newline Umgang mit ihnen berechtigt sind, von \newline diesem Recht zu ihrem Wohl Gebrauch \newline machen. Eltern, andere \newline Umgangsberechtigte sowie Personen, in \newline deren Obhut sich das Kind befindet, haben \newline Anspruch auf Beratung und Unterstützung \newline bei der Ausübung des Umgangsrechts. Bei \\
\hline
\multicolumn{3}{|c|}{\cellcolor{gray!10}\textbf{\S{} 27 -- Hilfe zur Erziehung}} \\
\hline
(1) Ein Personensorgeberechtigter hat bei \newline der Erziehung eines Kindes oder eines \newline Jugendlichen Anspruch auf Hilfe (Hilfe zur \newline Erziehung), wenn eine dem Wohl des \newline Kindes oder des Jugendlichen \newline entsprechende Erziehung nicht \newline gewährleistet ist und die Hilfe für seine \newline Entwicklung geeignet und notwendig ist. & \textbf{(1) Zur Verwirklichung des Rechts eines \newline jeden jungen Menschen auf Förderung \newline seiner Entwicklung und auf Erziehung \newline zu einer selbstbestimmten, \newline eigenverantwortlichen und \newline gemeinschaftsfähigen Persönlichkeit \newline durch Leistungen zur Entwicklung, zur \newline Erziehung und zur Teilhabe haben \newline Kinder, Jugendliche oder \newline Personensorgeberechtigte einen \newline Anspruch auf Hilfe zur Erziehung und \newline auf Leistungen der Eingliederungshilfe \newline nach Maßgabe der Absätze 2 und 3. } & \textbf{(1) Zur Verwirklichung des Rechts eines \newline jeden jungen Menschen auf Förderung \newline seiner Entwicklung, auf Erziehung zu \newline einer selbstbestimmten, \newline eigenverantwortlichen und \newline gemeinschaftsfähigen Persönlichkeit \newline und auf Förderung seiner vollen \newline wirksamen und gleichberechtigten \newline Teilhabe am Leben in der Gesellschaft \newline durch Leistungen zur Entwicklung, zur \newline Erziehung und zur Teilhabe haben \newline Kinder, Jugendliche oder \newline Personensorgeberechtigte einen \newline Anspruch auf Hilfe zur Erziehung oder \newline auf Leistungen der Eingliederungshilfe \newline nach Maßgabe der Absätze 2 und 3. } \\
\hline
\textit{--- Fassung 2024 ---} \newline (2) Hilfe zur Erziehung wird insbesondere \newline nach Maßgabe der \S{}\S{} 28 bis 35 gewährt. \newline Art und Umfang der Hilfe richten sich nach \newline dem erzieherischen Bedarf im Einzelfall; \newline dabei soll das engere soziale Umfeld des \newline Kindes oder des Jugendlichen einbezogen \newline werden. Unterschiedliche Hilfearten \newline können miteinander kombiniert werden, \newline sofern dies dem erzieherischen Bedarf des \newline Kindes oder Jugendlichen im Einzelfall \newline entspricht. \newline \textit{--- Fassung 2026 ---} \newline (2) Hilfe zur Erziehung wird insbesondere \newline nach Maßgabe der \S{}\S{} 28 bis 35 gewährt. \newline Art und Umfang der Hilfe richten sich nach \newline dem erzieherischen Bedarf im Einzelfall; & \textbf{(2) Personensorgeberechtigte haben \newline einen Anspruch auf Hilfe zur Erziehung, \newline wenn und solange eine dem Kindeswohl \newline entsprechende Erziehung nicht \newline gewährleistet ist und die Hilfe für die \newline Entwicklung der jungen Menschen \newline geeignet und notwendig ist. Liegen die \newline Voraussetzungen nach Satz 1 vor, \newline haben auch Jugendliche einen \newline Anspruch auf Hilfe zur Erziehung, die \newline außerhalb des Elternhauses erbracht \newline wird. } & \textbf{(2) Personensorgeberechtigte haben \newline einen Anspruch auf Hilfe zur Erziehung, \newline wenn und solange eine dem Kindeswohl \newline entsprechende Erziehung nicht } \\
\hline
dabei soll das engere soziale Umfeld des \newline Kindes oder des Jugendlichen einbezogen \newline werden. Unterschiedliche Hilfearten \newline können miteinander kombiniert werden, \newline sofern dies dem erzieherischen Bedarf des \newline Kindes oder Jugendlichen im Einzelfall \newline entspricht. & \textit{unverändert} & \textbf{gewährleistet ist und die Hilfe für die \newline Entwicklung des jungen Menschen \newline geeignet und notwendig ist.} \\
\hline
(2a) Ist eine Erziehung des Kindes oder \newline Jugendlichen außerhalb des Elternhauses \newline erforderlich, so entfällt der Anspruch auf \newline Hilfe zur Erziehung nicht dadurch, dass \newline eine andere unterhaltspflichtige Person \newline bereit ist, diese Aufgabe zu übernehmen; \newline die Gewährung von Hilfe zur Erziehung \newline setzt in diesem Fall voraus, dass diese \newline Person bereit und geeignet ist, den \newline Hilfebedarf in Zusammenarbeit mit dem \newline Träger der öffentlichen Jugendhilfe nach \newline Maßgabe der \S{}\S{} 36 und 37 zu decken. & (2a) entfällt & (2a) entfällt \\
\hline
(3) Hilfe zur Erziehung umfasst \newline insbesondere die Gewährung \newline pädagogischer und damit verbundener \newline therapeutischer Leistungen. Bei Bedarf soll \newline sie Ausbildungs- und \newline Beschäftigungsmaßnahmen im Sinne des \newline \S{} 13 Absatz 2 einschließen und kann mit \newline anderen Leistungen nach diesem Buch \newline kombiniert werden. Die in der Schule oder \newline Hochschule wegen des erzieherischen \newline Bedarfs erforderliche Anleitung und \newline Begleitung können als Gruppenangebote \newline an Kinder oder Jugendliche gemeinsam \newline erbracht werden, soweit dies dem Bedarf \newline des Kindes oder Jugendlichen im Einzelfall \newline entspricht. & \textbf{(3) Kinder oder Jugendliche mit \newline Behinderungen im Sinne von \S{} 7 Absatz \newline 2, die in der gleichberechtigten Teilhabe \newline an der Gesellschaft eingeschränkt sind \newline oder von einer solchen Behinderung \newline bedroht sind, haben einen Anspruch auf \newline Leistungen der Eingliederungshilfe, \newline wenn und solange diese Leistungen \newline nach der Besonderheit des Einzelfalles \newline geeignet und notwendig sind, den \newline jungen Menschen eine individuelle \newline Lebensführung zu ermöglichen, die der \newline Würde des Menschen entspricht, ihre \newline volle, wirksame und gleichberechtigte \newline Teilhabe am Leben in der Gesellschaft \newline zu fördern und sie zu befähigen, ihre \newline Lebensplanung und -führung möglichst \newline selbstbestimmt und eigenverantwortlich \newline wahrnehmen zu können. } & \textbf{(3) Kinder oder Jugendliche mit \newline Behinderungen oder von Behinderung \newline bedrohte Kinder oder Jugendliche im \newline Sinne von \S{} 7 Absatz 2 haben einen \newline Anspruch auf Leistungen der \newline Eingliederungshilfe, wenn und solange \newline diese Leistungen nach der Besonderheit \newline des Einzelfalles geeignet und notwendig \newline sind, die Aufgaben der \newline Eingliederungshilfe nach \S{} 90 des \newline Neunten Buches zu erfüllen. } \\
\hline
(4) Wird ein Kind oder eine Jugendliche \newline während ihres Aufenthalts in einer \newline Einrichtung oder einer Pflegefamilie selbst \newline Mutter eines Kindes, so umfasst die Hilfe \newline zur Erziehung auch die Unterstützung bei \newline der Pflege und Erziehung dieses Kindes. & \textbf{(4) Die Bundesregierung kann durch \newline Rechtsverordnung mit Zustimmung des \newline Bundesrates Bestimmungen über die \newline Konkretisierung der \newline Leistungsberechtigung nach dieser \newline Vorschrift erlassen. } & \textbf{(4) Maßgeblich für die Eignung und \newline Notwendigkeit der Leistungen der \newline Eingliederungshilfe sind die \newline Wechselwirkungen der geistigen, \newline seelischen, körperlichen oder \newline Sinnesbeeinträchtigungen mit \newline einstellungs- und umweltbedingten \newline Barrieren im Einzelfall und deren \newline konkrete Auswirkungen auf die Teilhabe \newline der jungen Menschen an der \newline Gesellschaft. Andere Leistungen der \newline Eingliederungshilfe können gewährt \newline werden. } \\
\hline
 & \textbf{(5) Besteht gleichzeitig ein Anspruch auf \newline Hilfe zur Erziehung nach Absatz 2 und \newline ein Anspruch auf Leistungen der \newline Eingliederungshilfe nach Absatz 3, so \newline sollen Einrichtungen, Dienste und \newline Personen die Hilfe und Leistungen \newline erbringen, die geeignet sind, sowohl \newline den erzieherischen Bedarf zu decken als \newline auch die Aufgaben der \newline Eingliederungshilfe zu erfüllen. } & \textbf{(5) Geeignete Leistungen können \newline gewährt werden, wenn die \newline Voraussetzungen der Notwendigkeit der \newline Leistungen nach Absatz 3 und Absatz 4 \newline nicht vorliegen. } \\
\hline
 & \textit{unverändert} & \textbf{(6) Besteht gleichzeitig ein Anspruch auf \newline Hilfe zur Erziehung nach Absatz 2 und \newline ein Anspruch auf Leistungen der \newline Eingliederungshilfe nach Absatz 3, so \newline sollen Einrichtungen, Dienste und \newline Personen die Hilfe und Leistungen \newline erbringen, die geeignet sind, sowohl \newline den erzieherischen Bedarf zu decken als \newline auch die Aufgaben der \newline Eingliederungshilfe zu erfüllen. } \\
\hline
 & \textbf{(3a) Maßgeblich für die Eignung und \newline Notwendigkeit der Leistungen der \newline Eingliederungshilfe sind insbesondere \newline die Wechselwirkungen der geistigen, \newline seelischen, körperlichen oder \newline Sinnesbeeinträchtigungen mit \newline einstellungs- und umweltbedingten \newline Barrieren im Einzelfall und deren \newline konkrete Auswirkungen auf die Teilhabe \newline der jungen Menschen an der \newline Gesellschaft. } & \textit{unverändert} \\
\hline
 & \textbf{(3b) Von einer Behinderung bedroht \newline sind Kinder oder Jugendliche, bei denen \newline der Eintritt einer Behinderung nach \newline fachlicher Erkenntnis mit hoher \newline Wahrscheinlichkeit zu erwarten ist. } & \textit{unverändert} \\
\hline
\multicolumn{3}{|c|}{\cellcolor{gray!10}\textbf{\S{} 27a -- Hilfe zur Erziehung}} \\
\hline
 & \textbf{(1) Besteht ein Anspruch auf Hilfe zur \newline Erziehung nach \S{} 27 Absatz 2 wird diese \newline insbesondere nach Maßgabe der \S{}\S{} 28 \newline bis 35 gewährt. Art und Umfang der Hilfe \newline richten sich nach dem erzieherischen \newline Bedarf im Einzelfall; dabei soll das \newline engere soziale Umfeld des Kindes oder \newline des Jugendlichen einbezogen werden. \newline Unterschiedliche Hilfearten können \newline miteinander kombiniert werden, sofern \newline dies dem erzieherischen Bedarf des \newline Kindes oder Jugendlichen im Einzelfall \newline entspricht. } & \textbf{(1) Besteht ein Anspruch auf Hilfe zur \newline Erziehung nach \S{} 27 Absatz 2 wird diese \newline insbesondere nach Maßgabe der \S{}\S{} 28 \newline bis 35 gewährt. Art und Umfang der Hilfe \newline richten sich nach dem erzieherischen \newline Bedarf im Einzelfall; dabei soll das \newline engere soziale Umfeld des Kindes oder \newline des Jugendlichen einbezogen werden. \newline Unterschiedliche Hilfearten können \newline miteinander kombiniert werden, sofern \newline dies dem erzieherischen Bedarf des \newline Kindes oder Jugendlichen im Einzelfall \newline entspricht. } \\
\hline
 & \textbf{(2) Ist eine Erziehung des Kindes oder \newline Jugendlichen außerhalb des \newline Elternhauses erforderlich, so entfällt der \newline Anspruch auf Hilfe zur Erziehung nicht \newline dadurch, dass eine andere \newline unterhaltspflichtige Person bereit ist, \newline diese Aufgabe zu übernehmen; die \newline Gewährung von Hilfe zur Erziehung \newline setzt in diesem Fall voraus, dass diese \newline Person bereit und geeignet ist, den \newline Hilfebedarf in Zusammenarbeit mit dem \newline Träger der öffentlichen Jugendhilfe \newline nach Maßgabe der \S{}\S{} 36, 36a und 38 zu \newline decken. } & \textbf{(2)} \textbf{Ist eine Erziehung des Kindes oder \newline Jugendlichen außerhalb des \newline Elternhauses erforderlich, so entfällt der \newline Anspruch auf Hilfe zur Erziehung nicht \newline dadurch, dass eine andere \newline unterhaltspflichtige Person bereit ist, \newline diese Aufgabe zu übernehmen. Die \newline Gewährung von Hilfe zur Erziehung \newline setzt in diesem Fall voraus, dass diese \newline Person bereit und geeignet ist, den \newline Hilfebedarf in Zusammenarbeit mit dem \newline Träger der öffentlichen Jugendhilfe \newline nach Maßgabe der \S{}\S{} 36, 36a und 39 zu \newline decken. } \\
\hline
 & \textbf{(3) Hilfe zur Erziehung umfasst \newline insbesondere die Gewährung \newline pädagogischer und damit verbundener \newline therapeutischer Leistungen. Bei Bedarf \newline soll sie Ausbildungs- und \newline Beschäftigungsmaßnahmen im Sinne \newline des \S{} 13 Absatz 2 einschließen und \newline kann mit anderen Leistungen nach \newline diesem Buch kombiniert werden. Die in \newline der Schule oder Hochschule wegen des \newline erzieherischen Bedarfs erforderliche \newline Anleitung und Begleitung können als \newline Gruppenangebote an Kinder oder \newline Jugendliche gemeinsam erbracht \newline werden, soweit dies dem Bedarf des \newline Kindes oder Jugendlichen im Einzelfall \newline entspricht, für diesen nach \S{} 5 Absatz 3 \newline zumutbar ist und mit \newline Leistungserbringern entsprechende \newline Vereinbarungen bestehen. Die \newline Leistungen nach Satz 1 sind auf \newline Wunsch des Leistungsberechtigten \newline gemeinsam zu erbringen. } & \textbf{(3) Hilfe zur Erziehung umfasst \newline insbesondere die Gewährung \newline pädagogischer und damit verbundener \newline therapeutischer Leistungen. Bei Bedarf \newline soll sie Ausbildungs- und \newline Beschäftigungsmaßnahmen im Sinne \newline des \S{} 13 Absatz 2 einschließen und \newline kann mit anderen Leistungen nach \newline diesem Buch kombiniert werden. } \\
\hline
 & \textbf{(4) Wird ein Kind oder eine Jugendliche \newline während ihres Aufenthalts in einer \newline Einrichtung oder einer Pflegefamilie \newline selbst Mutter eines Kindes, so umfasst \newline die Hilfe zur Erziehung auch die \newline Unterstützung bei der Pflege und \newline Erziehung dieses Kindes. } & \textbf{(4) Sofern infrastrukturelle Angebote \newline oder Regelangebote insbesondere nach \newline \S{}\S{} 16 bis 18, \S{}\S{} 22 bis 25 oder \S{} 13 im \newline Hinblick auf den Bedarf des Kindes oder \newline des Jugendlichen im Einzelfall \newline geeigneter oder gleichermaßen geeignet \newline sind, werden diese vorrangig gewährt. \newline Bei Jugendlichen und jungen \newline Volljährigen werden Hilfen oder \newline Maßnahmen nach \S{} 13 vorrangig \newline gewährt, wenn sie gleichermaßen \newline geeignet sind. } \\
\hline
 & \textit{unverändert} & \textbf{(5)} \textbf{Die in einer Tageseinrichtung für \newline Kinder nach \S{} 22 Absatz 1 Satz 1 oder in \newline der Schule oder Hochschule wegen des \newline erzieherischen Bedarfs erforderliche \newline Anleitung und Begleitung werden als \newline infrastrukturelle Angebote nach \S{} 80a \newline gewährt; die Vorschriften zur Hilfe- und \newline Leistungsplanung (\S{}\S{} 36 bis 38d) finden \newline keine Anwendung. Nur wenn dem \newline erzieherischen Bedarf im Einzelfall \newline ausschließlich durch eine an dem \newline jeweiligen Kind oder Jugendlichen \newline erbrachte Anleitung und Begleitung in \newline einer Tageseinrichtung für Kinder nach \newline \S{} 22 Absatz 1 Satz 1 oder in der Schule \newline oder Hochschule entsprochen werden \newline kann, besteht nach \S{} 27 Absatz 2 ein \newline Anspruch auf diese Einzelhilfe. } \\
\hline
 & \textit{unverändert} & \textbf{(6)} \textbf{Wird ein Kind oder eine Jugendliche \newline während des Aufenthalts in einer \newline Einrichtung oder einer Pflegefamilie \newline selbst Mutter eines Kindes, so umfasst \newline die Hilfe zur Erziehung auch die \newline Unterstützung bei der Pflege und \newline Erziehung dieses Kindes. } \\
\hline
\multicolumn{3}{|c|}{\cellcolor{gray!10}\textbf{\S{} 34 -- Heimerziehung, sonstige betreute Wohnform}} \\
\hline
(1) Hilfe zur Erziehung in einer Einrichtung \newline über Tag und Nacht (Heimerziehung) oder \newline in einer sonstigen betreuten Wohnform soll \newline Kinder und Jugendliche durch eine \newline Verbindung von Alltagserleben mit \newline pädagogischen und therapeutischen \newline Angeboten in ihrer Entwicklung fördern. Sie \newline soll entsprechend dem Alter und \newline Entwicklungsstand des Kindes oder des \newline Jugendlichen sowie den Möglichkeiten der \newline Verbesserung der Erziehungsbedingungen \newline in der Herkunftsfamilie & (3) unverändert \newline \textbf{Unterabschnitt 3: Leistungen der \newline Eingliederungshilfe für Kinder und \newline Jugendliche mit Behinderungen } & (1) unverändert \newline \textbf{Unterabschnitt 3: Leistungen der \newline Eingliederungshilfe für Kinder und \newline Jugendliche mit Behinderungen } \\
\hline
1. \newline eine Rückkehr in die Familie zu erreichen \newline versuchen oder & 1. unverändert & 1. unverändert \\
\hline
2. \newline die Erziehung in einer anderen Familie \newline vorbereiten oder & 2. unverändert & 2. unverändert \\
\hline
3. \newline eine auf längere Zeit angelegte \newline Lebensform bieten und auf ein \newline selbständiges Leben vorbereiten. \newline Jugendliche sollen in Fragen der \newline Ausbildung und Beschäftigung sowie der \newline allgemeinen Lebensführung beraten und \newline unterstützt werden. & 3. unverändert & 3. unverändert \\
\hline
\multicolumn{3}{|c|}{\cellcolor{gray!10}\textbf{\S{} 35a -- Eingliederungshilfe für Kinder und Jugendliche mit seelischer Behinderung oder drohender seelischer Behinderung}} \\
\hline
(1) Kinder oder Jugendliche haben \newline Anspruch auf Eingliederungshilfe, wenn & \textbf{(1) Besteht ein Anspruch auf Leistungen \newline der Eingliederungshilfe nach \S{} 27 \newline Absatz 3, werden diese insbesondere \newline nach Maßgabe der Absätze 2 und 3 \newline sowie der \S{}\S{} 35b bis 35i und den \newline Kapiteln 9 bis 13 des Teils 1 des \newline Neunten Buches gewährt; die Kapitel 3 \newline bis 6 des Teils 2 des Neunten Buches \newline gelten im Übrigen entsprechend. Art \newline und Umfang der Leistungen richten sich \newline nach dem Ergebnis der Prüfung gemäß \newline \S{} 27 Absatz 3a und bestimmen sich \newline nach der Besonderheit des Einzelfalles, \newline insbesondere nach dem individuellen \newline Bedarf, den persönlichen Verhältnissen, \newline dem engeren sozialen Umfeld, dem \newline Sozialraum und den eigenen Kräften \newline und Mitteln; dabei ist auch die \newline Wohnform zu würdigen. \newline Unterschiedliche Leistungsarten \newline können miteinander kombiniert werden, \newline sofern dies dem Bedarf des Kindes oder \newline Jugendlichen im Einzelfall entspricht. } & \textbf{(1) Besteht ein Anspruch auf Leistungen \newline der Eingliederungshilfe nach \S{} 27 \newline Absatz 3, werden diese insbesondere \newline nach Maßgabe der Absätze 2 bis 7 \newline sowie der \S{}\S{} 35b bis 35i und der Kapitel \newline 9 bis 13 des Teils 1 des Neunten Buches \newline gewährt; \S{} 107 sowie die Kapitel 3 bis 6 } \newline \textbf{des Teils 2 des Neunten Buches gelten \newline im Übrigen entsprechend. Art und \newline Umfang der Leistungen richten sich \newline nach dem Ergebnis der Prüfung gemäß \newline \S{} 27 Absatz 4 und bestimmen sich nach \newline der Besonderheit des Einzelfalles, \newline insbesondere nach dem individuellen \newline Bedarf, den persönlichen Verhältnissen, \newline dem engeren sozialen Umfeld, dem \newline Sozialraum und den eigenen Kräften \newline und Mitteln; dabei ist auch die \newline Wohnform zu würdigen. \newline Unterschiedliche Leistungsarten der \newline Eingliederungshilfe können miteinander \newline kombiniert werden, sofern dies dem \newline Bedarf des Kindes oder Jugendlichen \newline im Einzelfall entspricht. } \\
\hline
1. \newline ihre seelische Gesundheit mit hoher \newline Wahrscheinlichkeit länger als sechs \newline Monate von dem für ihr Lebensalter \newline typischen Zustand abweicht, und & 1. entfällt & 1. entfällt \\
\hline
2. \newline daher ihre Teilhabe am Leben in der \newline Gesellschaft beeinträchtigt ist oder eine \newline solche Beeinträchtigung zu erwarten ist. \newline Von einer seelischen Behinderung bedroht \newline im Sinne dieser Vorschrift sind Kinder oder \newline Jugendliche, bei denen eine \newline Beeinträchtigung ihrer Teilhabe am Leben \newline in der Gesellschaft nach fachlicher \newline Erkenntnis mit hoher Wahrscheinlichkeit zu \newline erwarten ist. \S{} 27 Absatz 4 gilt \newline entsprechend. & 2. entfällt & 2. entfällt \\
\hline
\textit{--- Fassung 2024 ---} \newline (1a) Hinsichtlich der Abweichung der \newline seelischen Gesundheit nach Absatz 1 Satz \newline 1 Nummer 1 hat der Träger der öffentlichen \newline Jugendhilfe die Stellungnahme \newline \textit{--- Fassung 2026 ---} \newline (1a) Hinsichtlich der Abweichung der \newline seelischen Gesundheit nach Absatz 1 Satz \newline 1 Nummer 1 hat der Träger der öffentlichen \newline Jugendhilfe die Stellungnahme \newline für die Behandlung von Kindern und \newline Jugendlichen oder & (1a) entfällt & (1a) entfällt \\
\hline
1. \newline eines Arztes für Kinder- und \newline Jugendpsychiatrie und -psychotherapie, & 1. entfällt & 1. entfällt \\
\hline
\textit{--- Fassung 2024 ---} \newline 2. \newline eines Kinder- und \newline Jugendlichenpsychotherapeuten, eines \newline Psychotherapeuten mit einer Weiterbildung \newline für die Behandlung von Kindern und \newline Jugendlichen oder \newline \textit{--- Fassung 2026 ---} \newline 2. \newline eines Kinder- und \newline Jugendlichenpsychotherapeuten, eines \newline Psychotherapeuten mit einer Weiterbildung & 2. entfällt & 2. entfällt \\
\hline
3. \newline eines Arztes oder eines psychologischen \newline Psychotherapeuten, der über besondere \newline Erfahrungen auf dem Gebiet seelischer \newline Störungen bei Kindern und Jugendlichen \newline verfügt, einzuholen. Die Stellungnahme ist \newline auf der Grundlage der Internationalen \newline Klassifikation der Krankheiten in der vom \newline Bundesinstitut für Arzneimittel und \newline Medizinprodukte herausgegebenen \newline deutschen Fassung zu erstellen. Dabei ist \newline auch darzulegen, ob die Abweichung \newline Krankheitswert hat oder auf einer Krankheit \newline beruht. Enthält die Stellungnahme auch \newline Ausführungen zu Absatz 1 Satz 1 Nummer \newline 2, so sollen diese vom Träger der \newline öffentlichen Jugendhilfe im Rahmen seiner \newline Entscheidung angemessen berücksichtigt \newline werden. Die Hilfe soll nicht von der Person \newline oder dem Dienst oder der Einrichtung, der \newline die Person angehört, die die \newline Stellungnahme abgibt, erbracht werden. & 3. entfällt & 3. entfällt \\
\hline
(2) Die Hilfe wird nach dem Bedarf im \newline Einzelfall & \textbf{(2) Die Leistungen der \newline Eingliederungshilfe umfassen } & \textbf{(2) Die Leistungen der \newline Eingliederungshilfe umfassen } \\
\hline
1. \newline in ambulanter Form, & \textbf{1. \newline Leistungen zur medizinischen \newline Rehabilitation, } & \textbf{1. \newline Leistungen zur medizinischen \newline Rehabilitation, } \\
\hline
2. \newline in Tageseinrichtungen für Kinder oder in \newline anderen teilstationären Einrichtungen, & \textbf{2. \newline Leistungen zur Teilhabe am \newline Arbeitsleben, } & \textbf{2. \newline Leistungen zur Teilhabe am \newline Arbeitsleben, } \\
\hline
3. \newline durch geeignete Pflegepersonen und & \textbf{3. \newline Leistungen zur Teilhabe an Bildung und } & \textbf{3. \newline Leistungen zur Teilhabe an Bildung und } \\
\hline
4. \newline in Einrichtungen über Tag und Nacht sowie \newline sonstigen Wohnformen geleistet. & \textbf{4. \newline Leistungen zur Sozialen Teilhabe. \newline Leistungen nach Satz 1 Nummer 1 bis 3 \newline gehen den Leistungen nach Satz 1 \newline Nummer 4 vor. } & \textbf{4. \newline Leistungen zur Sozialen Teilhabe. \newline Leistungen nach Satz 1 Nummer 1 bis 3 \newline gehen den Leistungen nach Satz 1 \newline Nummer 4 vor. } \\
\hline
 & \textbf{5. \newline in Einrichtungen über Tag und Nacht \newline sowie sonstigen Wohnformen erbracht. \newline Dabei sollen Einrichtungen, Dienste \newline oder Personen die Leistungen erbringen \newline oder in Anspruch genommen werden, in \newline denen Kinder oder Jugendliche mit \newline Behinderungen und Kinder oder \newline Jugendliche ohne Behinderungen \newline gemeinsam Leistungen erhalten \newline können, wenn die Aufgaben der \newline Eingliederungshilfe erfüllt werden \newline können; die besonderen Bedürfnisse \newline von Kindern oder Jugendlichen mit \newline Behinderungen und von Kindern oder \newline Jugendlichen, die von einer \newline Behinderung bedroht sind, sind zu \newline berücksichtigen. } & \textit{unverändert} \\
\hline
\textit{--- Fassung 2024 ---} \newline (3) Aufgabe und Ziele der Hilfe, die \newline Bestimmung des Personenkreises sowie \newline Art und Form der Leistungen richten sich \newline nach Kapitel 6 des Teils 1 des Neunten \newline Buches sowie \S{} 90 und den Kapiteln 3 bis \newline 6 des Teils 2 des Neunten Buches, soweit \newline diese Bestimmungen auch auf seelisch \newline behinderte oder von einer solchen \newline Behinderung bedrohte Personen \newline Anwendung finden und sich aus diesem \newline Buch nichts anderes ergibt. \newline \textit{--- Fassung 2026 ---} \newline (3) Aufgabe und Ziele der Hilfe, die \newline Bestimmung des Personenkreises sowie \newline Art und Form der Leistungen richten sich \newline nach Kapitel 6 des Teils 1 des Neunten & \textbf{(3) Leistungen der Eingliederungshilfe \newline werden als Sach-, Geld- oder \newline Dienstleistung erbracht. Sie können bei \newline Bedarf mit anderen Leistungen nach \newline diesem Buch kombiniert werden. } & \textbf{(3) Leistungen der Eingliederungshilfe \newline werden als Sach-, Geld- oder \newline Dienstleistung erbracht. Sie können bei } \\
\hline
Buches sowie \S{} 90 und den Kapiteln 3 bis \newline 6 des Teils 2 des Neunten Buches, soweit \newline diese Bestimmungen auch auf seelisch \newline behinderte oder von einer solchen \newline Behinderung bedrohte Personen \newline Anwendung finden und sich aus diesem \newline Buch nichts anderes ergibt. & \textit{unverändert} & \textbf{Bedarf mit anderen Leistungen nach \newline diesem Buch kombiniert werden. } \\
\hline
(4) Ist gleichzeitig Hilfe zur Erziehung zu \newline leisten, so sollen Einrichtungen, Dienste \newline und Personen in Anspruch genommen \newline werden, die geeignet sind, sowohl die \newline Aufgaben der Eingliederungshilfe zu \newline erfüllen als auch den erzieherischen Bedarf \newline zu decken. Sind heilpädagogische \newline Maßnahmen für Kinder, die noch nicht im \newline schulpflichtigen Alter sind, in \newline Tageseinrichtungen für Kinder zu \newline gewähren und lässt der Hilfebedarf es zu, \newline so sollen Einrichtungen in Anspruch \newline genommen werden, in denen behinderte \newline und nicht behinderte Kinder gemeinsam \newline betreut werden. & \textbf{(4) Dienstleistungen werden nach dem \newline Bedarf im Einzelfall } & \textbf{(4) Dienstleistungen werden nach dem \newline Bedarf im Einzelfall } \\
\hline
 & \textbf{1. in ambulanter Form, } & \textbf{1. in ambulanter Form, } \\
\hline
 & \textbf{2. in Tageseinrichtungen für Kinder oder \newline in anderen teilstationären \newline Einrichtungen, } & \textbf{2. in Tageseinrichtungen für Kinder oder \newline in anderen teilstationären \newline Einrichtungen, } \\
\hline
 & \textit{unverändert} & \textbf{3. durch geeignete Pflegepersonen und } \\
\hline
 & \textbf{4. durch geeignete Pflegepersonen und } & \textbf{4. in Einrichtungen über Tag und Nacht \newline sowie sonstigen Wohnformen erbracht. \newline Dabei sollen Einrichtungen, Dienste \newline oder Personen die Leistungen erbringen \newline oder in Anspruch genommen werden, in \newline denen Kinder oder Jugendliche mit \newline Behinderungen und Kinder oder \newline Jugendliche ohne Behinderungen \newline gemeinsam Leistungen erhalten \newline können, wenn die Aufgaben der \newline Eingliederungshilfe erfüllt werden \newline können; die besonderen Bedürfnisse \newline von Kindern oder Jugendlichen mit \newline Behinderungen und von Kindern oder \newline Jugendlichen, die von einer \newline Behinderung bedroht sind, sind zu \newline berücksichtigen. } \\
\hline
 & \textbf{5. in Einrichtungen über Tag und Nacht \newline sowie sonstigen Wohnformen erbracht. \newline Dabei sollen Einrichtungen, Dienste \newline oder Personen die Leistungen erbringen \newline oder in Anspruch genommen werden, in \newline denen Kinder oder Jugendliche mit \newline Behinderungen und Kinder oder \newline Jugendliche ohne Behinderungen \newline gemeinsam Leistungen erhalten \newline können, wenn die Aufgaben der \newline Eingliederungshilfe erfüllt werden \newline können; die besonderen Bedürfnisse \newline von Kindern oder Jugendlichen mit \newline Behinderungen und von Kindern oder \newline Jugendlichen, die von einer \newline Behinderung bedroht sind, sind zu \newline berücksichtigen. } & \textit{unverändert} \\
\hline
 & \textbf{(5) Leistungen zur Sozialen Teilhabe \newline können mit Zustimmung der \newline Leistungsberechtigten auch in Form \newline einer pauschalen Geldleistung erbracht \newline werden, soweit es nach \S{} 35i \newline vorgesehen ist. Die Träger der \newline öffentlichen Jugendhilfe regeln das \newline Nähere zur Höhe und Ausgestaltung der \newline Pauschalen. } & \textbf{(5) Leistungen zur Sozialen Teilhabe \newline können mit Zustimmung der \newline Leistungsberechtigten auch in Form } \newline \textbf{einer pauschalen Geldleistung erbracht \newline werden, soweit es nach \S{} 35i \newline vorgesehen ist. Die Träger der \newline öffentlichen Jugendhilfe regeln das \newline Nähere zur Höhe und Ausgestaltung der \newline Pauschalen. } \\
\hline
 & \textbf{(6) Die Leistungen der \newline Eingliederungshilfe werden auf Antrag \newline auch als Teil eines Persönlichen \newline Budgets ausgeführt. Die Betroffenen \newline sind entsprechend zu beraten. Die \newline Vorschrift zum Persönlichen Budget \newline nach \S{} 29 des Neunten Buches ist \newline insoweit anzuwenden. } & \textbf{(6) Die Leistungen der \newline Eingliederungshilfe werden auf Antrag \newline auch durch ein Persönliches Budget \newline ausgeführt. Der Leistungsberechtigte \newline und der Personensorgeberechtigte sind \newline entsprechend zu beraten. Die Vorschrift \newline zum Persönlichen Budget nach \S{} 29 des \newline Neunten Buches ist insoweit \newline anzuwenden. } \\
\hline
 & \textit{unverändert} & \textbf{(7) \S{} 103 des Neunten Buches gilt \newline entsprechend. } \\
\hline
\multicolumn{3}{|c|}{\cellcolor{gray!10}\textbf{\S{} 35b -- Leistungen zur medizinischen Rehabilitation}} \\
\hline
 & \textbf{(1) Die Leistungen } & \textbf{(1) Die Leistungen } \\
\hline
 & \textbf{1. \newline zur Assistenz zur Übernahme von \newline Handlungen zur Alltagsbewältigung \newline sowie Begleitung der \newline Leistungsberechtigten (\S{} 35f Absatz 2 \newline Nummer 2 in Verbindung mit \S{} 78 \newline Absatz 2 Nummer 1 und Absatz 5 des \newline Neunten Buches), } & \textbf{1. \newline zur Assistenz zur Übernahme von \newline Handlungen zur Alltagsbewältigung \newline sowie Begleitung der \newline Leistungsberechtigten (\S{} 35f Absatz 2 \newline Nummer 2 in Verbindung mit \S{} 78 \newline Absatz 2 Nummer 1 und Absatz 5 des \newline Neunten Buches), } \\
\hline
 & \textbf{2. \newline zur Förderung der Verständigung (\S{} 35f \newline Absatz 2 Nummer 6) und } & \textbf{2. \newline zur Förderung der Verständigung (\S{} 35f \newline Absatz 2 Nummer 6) und } \\
\hline
 & \textbf{3. \newline zur Beförderung im Rahmen der \newline Leistungen zur Mobilität (\S{} 35f Absatz 2 \newline Nummer 7 in Verbindung mit \S{} 83 \newline Absatz 1 Nummer 1 des Neunten \newline Buches) können mit Zustimmung der \newline Leistungsberechtigten als pauschale \newline Geldleistungen nach \S{} 35a Absatz 3 \newline Satz 3 erbracht werden. Die zuständigen \newline Träger der öffentlichen Jugendhilfe \newline regeln das Nähere zur Höhe und \newline Ausgestaltung der pauschalen \newline Geldleistungen sowie zur \newline Leistungserbringung. } & \textbf{3. \newline zur Beförderung im Rahmen der \newline Leistungen zur Mobilität (\S{} 35f Absatz 2 \newline Nummer 7 in Verbindung mit \S{} 83 \newline Absatz 1 Nummer 1 des Neunten \newline Buches) \newline können mit Zustimmung der \newline Leistungsberechtigten als pauschale \newline Geldleistungen nach \S{} 35a Absatz 3 \newline Satz 3 erbracht werden. Die zuständigen \newline Träger der öffentlichen Jugendhilfe \newline regeln das Nähere zur Höhe und \newline Ausgestaltung der pauschalen \newline Geldleistungen sowie zur \newline Leistungserbringung. } \\
\hline
 & \textbf{(2) Die Leistungen } & \textbf{(2) Die Leistungen } \newline \textbf{zur Förderung der Verständigung (\S{} 35f \newline Absatz 2 Nummer 6), } \\
\hline
 & \textbf{1. \newline zur Assistenz (\S{} 35f Absatz 2 Nummer \newline 2), } & \textbf{1. \newline zur Assistenz (\S{} 35f Absatz 2 Nummer \newline 2), } \\
\hline
 & \textbf{2. \newline zur Heilpädagogik (\S{} 35f Absatz 2 \newline Nummer 3), } & \textbf{2. \newline zur Heilpädagogik (\S{} 35f Absatz 2 \newline Nummer 3), } \\
\hline
 & \textbf{3. \newline zum Erwerb und Erhalt praktischer \newline Fähigkeiten und Kenntnisse (\S{} 35f \newline Absatz 2 Nummer 5), } & \textbf{3. \newline zum Erwerb und Erhalt praktischer \newline Fähigkeiten und Kenntnisse (\S{} 35f \newline Absatz 2 Nummer 5), } \\
\hline
 & \textbf{4. \newline zur Förderung der Verständigung (\S{} 35f \newline Absatz 2 Nummer 6), } & \textbf{4. } \\
\hline
 & \textbf{5. \newline zur Beförderung im Rahmen der \newline Leistungen zur Mobilität (\S{} 35f Absatz 2 \newline Nummer 7 in Verbindung mit \S{} 83 \newline Absatz 1 Nummer 1 des Neunten \newline Buches) und } & \textbf{5. \newline zur Beförderung im Rahmen der \newline Leistungen zur Mobilität (\S{} 35f Absatz 2 \newline Nummer 7 in Verbindung mit \S{} 83 \newline Absatz 1 Nummer 1 des Neunten \newline Buches) und } \\
\hline
 & \textbf{6. \newline zur Erreichbarkeit einer \newline Ansprechperson unabhängig von einer \newline konkreten Inanspruchnahme (\S{} 35f \newline Absatz 2 Nummer 2 in Verbindung mit \S{} \newline 78 Absatz 6 des Neunten Buches) \newline können an mehrere \newline Leistungsberechtigte gemeinsam \newline erbracht werden, soweit dies nach \S{} 5 \newline Absatz 3 für die Leistungsberechtigten \newline zumutbar ist und mit \newline Leistungserbringern entsprechende \newline Vereinbarungen bestehen. Maßgeblich \newline sind die Ermittlungen und \newline Feststellungen nach \S{}\S{} 36 bis 38d } & \textbf{6. \newline zur Erreichbarkeit einer \newline Ansprechperson unabhängig von einer \newline konkreten Inanspruchnahme (\S{} 35f \newline Absatz 2 Nummer 2 in Verbindung mit \S{} \newline 78 Absatz 6 des Neunten Buches) \newline können an mehrere \newline Leistungsberechtigte gemeinsam \newline erbracht werden, soweit dies nach \S{} 5 \newline Absatz 3 für die Leistungsberechtigten \newline zumutbar ist und mit \newline Leistungserbringern entsprechende \newline Vereinbarungen bestehen. Maßgeblich \newline sind die Ermittlungen und \newline Feststellungen nach \S{}\S{} 36 bis 38d. \S{} 35f \newline Absatz 4 bleibt unberührt. } \\
\hline
 & \textbf{(3) Die Leistungen nach Absatz 2 sind \newline auf Wunsch der Leistungsberechtigten \newline gemeinsam zu erbringen, soweit die \newline Teilhabeziele erreicht werden können. } & \textbf{(3) Auf Wunsch der \newline Leistungsberechtigten sind die \newline Leistungen nach Absatz 2 sind an \newline mehrere Leistungsberechtigte \newline gemeinsam zu erbringen, soweit die \newline Teilhabeziele erreicht werden können. } \\
\hline
 & \textbf{(4) Zur Ermöglichung der \newline gemeinschaftlichen Mittagsverpflegung \newline in der Verantwortung einer Werkstatt für \newline behinderte Menschen, einem anderen \newline Leistungsanbieter oder dem \newline Leistungserbringer vergleichbarer \newline anderer tagesstrukturierender \newline Maßnahmen werden die erforderliche \newline sächliche Ausstattung, die personelle \newline Ausstattung und die erforderlichen \newline betriebsnotwendigen Anlagen des \newline Leistungserbringers übernommen. } & \textbf{(4) Die in einer Tageseinrichtung für \newline Kinder nach \S{} 22 Absatz 1 Satz 1 wegen \newline der Behinderung erforderliche Anleitung \newline und Begleitung werden als \newline infrastrukturelle Angebote nach \S{} 80a } \newline \textbf{gewährt; die Vorschriften zur Hilfe- und \newline Leistungsplanung (\S{}\S{} 36 bis 38d) finden \newline keine Anwendung. Nur wenn dem \newline Ergebnis der Prüfung gemäß \S{} 27 \newline Absatz 4 und der Besonderheit des \newline Einzelfalles ausschließlich durch eine \newline an dem jeweiligen Kind erbrachte \newline Anleitung und Begleitung in einer \newline Tageseinrichtung für Kinder nach \S{} 22 \newline Absatz 1 Satz 1 entsprochen werden \newline kann, besteht nach \S{} 27 Absatz 2 ein \newline Anspruch auf diese Einzelhilfe. } \\
\hline
 & \textbf{(5) In besonderen Wohnformen des \S{} \newline 42a Absatz 2 Satz 1 Nummer 2 und Satz \newline 3 des Zwölften Buches werden \newline Aufwendungen für Wohnraum oberhalb \newline der Angemessenheitsgrenze nach \S{} 42a \newline Absatz 6 des Zwölften Buches \newline übernommen, sofern dies wegen der \newline besonderen Bedürfnisse des Kindes \newline oder Jugendlichen mit Behinderungen \newline erforderlich ist. \S{}\S{} 78a bis 78g sind \newline anzuwenden. } & \textbf{(5) Zur Ermöglichung der \newline gemeinschaftlichen Mittagsverpflegung \newline in der Verantwortung einer Werkstatt für \newline behinderte Menschen, einem anderen \newline Leistungsanbieter oder dem \newline Leistungserbringer vergleichbarer \newline anderer tagesstrukturierender \newline Maßnahmen werden die erforderliche \newline sächliche Ausstattung, die personelle \newline Ausstattung und die erforderlichen \newline betriebsnotwendigen Anlagen des \newline Leistungserbringers übernommen. } \\
\hline
 & \textbf{(6) Bei einer stationären \newline Krankenhausbehandlung nach \S{} 39 des \newline Fünften Buches werden auch \newline Leistungen für die Begleitung und \newline Befähigung des Kindes oder \newline Jugendlichen durch vertraute \newline Bezugspersonen zur Sicherstellung der \newline Durchführung der Behandlung erbracht, \newline soweit dies aufgrund des \newline Vertrauensverhältnisses des Kindes \newline oder Jugendlichen zur Bezugsperson \newline und aufgrund der \newline behinderungsbedingten besonderen \newline Bedürfnisse erforderlich ist. Vertraute \newline Bezugspersonen im Sinne von Satz 1 \newline sind Eltern, Geschwister oder Personen, \newline die dem Kind oder Jugendlichen \newline gegenüber im Alltag bereits Leistungen \newline der Eingliederungshilfe erbringen. Die \newline Leistungen umfassen Leistungen zur \newline Verständigung und zur Unterstützung \newline im Umgang mit Belastungssituationen \newline als nichtmedizinische Nebenleistungen \newline zur stationären \newline Krankenhausbehandlung. Bei den \newline Leistungen im Sinne von Satz 1 findet \S{} \newline 91 Absatz 1 und 2 des Neunten Buches \newline gegenüber Kostenträgern von \newline Leistungen zur Krankenbehandlung mit \newline Ausnahme der Träger der \newline Unfallversicherung keine Anwendung. \S{} \newline 17 Absatz 2 und 2a des Ersten Buches \newline bleibt unberührt. } \newline \textbf{\S{} 35 g } \newline \textbf{Leistungen zur Mobilität} \newline \textbf{Bei den Leistungen zur Mobilität nach \S{} \newline 35f Absatz 2 Nummer 7 gilt \S{} 83 des \newline Neunten Buches mit der Maßgabe, \newline dass } \newline \textbf{\S{} 35 h } \newline \textbf{Besuchsbeihilfen} \newline \textbf{Werden Leistungen bei einem oder \newline mehreren Anbietern über Tag und Nacht \newline erbracht, können den \newline Leistungsberechtigten oder ihren \newline Angehörigen zum gegenseitigen Besuch \newline Beihilfen geleistet werden, soweit es im \newline Einzelfall erforderlich ist. } \newline \textbf{\S{} 35 i } \newline \textbf{Pauschale Geldleistung, \newline gemeinsame Inanspruchnahme } & \textbf{(6) Bei einer stationären \newline Krankenhausbehandlung nach \S{} 39 des \newline Fünften Buches werden auch \newline Leistungen für die Begleitung und \newline Befähigung des Kindes oder \newline Jugendlichen durch vertraute \newline Bezugspersonen zur Sicherstellung der \newline Durchführung der Behandlung erbracht, \newline soweit dies aufgrund des \newline Vertrauensverhältnisses des Kindes \newline oder Jugendlichen zur Bezugsperson \newline und aufgrund der \newline behinderungsbedingten besonderen \newline Bedürfnisse erforderlich ist. Vertraute \newline Bezugspersonen im Sinne von Satz 1 \newline sind Eltern, Geschwister oder Personen, \newline die dem Kind oder Jugendlichen \newline gegenüber im Alltag bereits Leistungen \newline der Eingliederungshilfe erbringen. Die \newline Leistungen umfassen Leistungen zur \newline Verständigung und zur Unterstützung \newline im Umgang mit Belastungssituationen \newline als nichtmedizinische Nebenleistungen \newline zur stationären \newline Krankenhausbehandlung. Bei den \newline Leistungen im Sinne von Satz 1 findet \S{} \newline 91 Absatz 1 und 2 des Neunten Buches } \newline \textbf{gegenüber Kostenträgern von \newline Leistungen zur Krankenbehandlung mit \newline Ausnahme der Träger der \newline Unfallversicherung keine Anwendung. \S{} \newline 17 Absatz 2 und 2a des Ersten Buches \newline bleibt unberührt. } \newline \textbf{\S{} 35 g } \newline \textbf{Leistungen zur Mobilität} \newline \textbf{Bei den Leistungen zur Mobilität nach \S{} \newline 35f Absatz 2 Nummer 7 gilt \S{} 83 des \newline Neunten Buches mit der Maßgabe, \newline dass } \newline \textbf{\S{} 35 h } \newline \textbf{Besuchsbeihilfen} \newline \textbf{Werden Leistungen bei einem oder \newline mehreren Anbietern über Tag und Nacht \newline erbracht, können den \newline Leistungsberechtigten oder ihren \newline Angehörigen zum gegenseitigen Besuch \newline Beihilfen geleistet werden, soweit es im \newline Einzelfall erforderlich ist. } \newline \textbf{\S{} 35 i } \newline \textbf{Pauschale Geldleistung, \newline gemeinsame Inanspruchnahme } \\
\hline
 & \textbf{1. \newline die Leistungsberechtigten zusätzlich zu \newline den in\S{} 83 Absatz 2 des Neunten \newline Buches genannten Voraussetzungen \newline zur Teilhabe am Leben in der \newline Gemeinschaft ständig auf die Nutzung \newline eines Kraftfahrzeugs angewiesen sind \newline und } & \textbf{1. \newline die Leistungsberechtigten zusätzlich zu \newline den in \S{} 83 Absatz 2 des Neunten \newline Buches genannten Voraussetzungen \newline zur Teilhabe am Leben in der \newline Gemeinschaft ständig auf die Nutzung \newline eines Kraftfahrzeugs angewiesen sind \newline und } \\
\hline
 & \textbf{2. \newline abweichend von \S{} 83 Absatz 3 Satz 2 \newline des Neunten Buches die Vorschriften \newline der \S{}\S{} 6 und 8 der Kraftfahrzeughilfe- \newline Verordnung nicht maßgeblich sind. } & \textbf{2. \newline abweichend von \S{} 83 Absatz 3 Satz 2 \newline des Neunten Buches die Vorschriften \newline der \S{}\S{} 6 und 8 der Kraftfahrzeughilfe- \newline Verordnung nicht maßgeblich sind. } \\
\hline
\textbf{Unterabschnitt 3 } & \textbf{Unterabschnitt 4 \newline Gemeinsame Vorschriften für \newline Leistungen zur Entwicklung, zur \newline Erziehung und zur Teilhabe } & \textbf{Unterabschnitt 4 \newline Gemeinsame Vorschriften für \newline Leistungen zur Entwicklung, zur \newline Erziehung und zur Teilhabe } \\
\hline
\multicolumn{3}{|c|}{\cellcolor{gray!10}\textbf{\S{} 36 -- Mitwirkung, Hilfeplan}} \\
\hline
 & \textbf{(1) Zur Sicherstellung von Kontinuität \newline und Bedarfsgerechtigkeit der \newline Leistungsgewährung sind von den \newline zuständigen öffentlichen Stellen, \newline insbesondere von \newline Sozialleistungsträgern oder \newline Rehabilitationsträgern rechtzeitig im \newline Rahmen des Hilfe- und Leistungsplans \newline Vereinbarungen zur Durchführung des \newline Zuständigkeitsübergangs zu treffen. Im \newline Rahmen der Beratungen zum \newline Zuständigkeitsübergang prüfen der \newline Träger der öffentlichen Jugendhilfe und \newline die andere öffentliche Stelle, \newline insbesondere der andere \newline Sozialleistungsträger oder \newline Rehabilitationsträger gemeinsam, \newline welche Leistung nach dem \newline Zuständigkeitsübergang dem Bedarf \newline des jungen Menschen entspricht. } & \textbf{(1) Zur Sicherstellung von Kontinuität \newline und Bedarfsgerechtigkeit der \newline Leistungsgewährung sind von den \newline zuständigen öffentlichen Stellen, \newline insbesondere von \newline Sozialleistungsträgern oder \newline Rehabilitationsträgern rechtzeitig im \newline Rahmen des Hilfe- und Leistungsplans \newline Vereinbarungen zur Durchführung des \newline Zuständigkeitsübergangs zu treffen. Im \newline Rahmen der Beratungen zum \newline Zuständigkeitsübergang prüfen der \newline Träger der öffentlichen Jugendhilfe und \newline die andere öffentliche Stelle, \newline insbesondere der andere \newline Sozialleistungsträger oder \newline Rehabilitationsträger gemeinsam, \newline welche Leistung nach dem \newline Zuständigkeitsübergang dem Bedarf \newline des jungen Menschen entspricht. } \\
\hline
 & \textbf{(2) Abweichend von Absatz 1 werden \newline bei einem Zuständigkeitsübergang vom \newline Träger der öffentlichen Jugendhilfe auf \newline einen Träger der Eingliederungshilfe \newline rechtzeitig im Rahmen eines \newline Teilhabeplanverfahrens nach \S{} 19 des \newline Neunten Buches die Voraussetzungen \newline für die Sicherstellung einer nahtlosen \newline und bedarfsgerechten \newline Leistungsgewährung nach dem \newline Zuständigkeitsübergang geklärt. Die \newline Teilhabeplanung ist frühzeitig, in der \newline Regel ein Jahr vor dem \newline voraussichtlichen \newline Zuständigkeitswechsel, vom Träger der \newline Jugendhilfe einzuleiten. Mit Zustimmung \newline des Leistungsberechtigten oder seines \newline Personensorgeberechtigten ist eine \newline Teilhabeplankonferenz nach \S{} 20 des \newline Neunten Buches durchzuführen. Stellt \newline der beteiligte Träger der \newline Eingliederungshilfe fest, dass seine \newline Zuständigkeit sowie die \newline Leistungsberechtigung absehbar \newline gegeben sind, soll er entsprechend \S{} 19 \newline Absatz 5 des Neunten Buches die \newline Teilhabeplanung vom Träger der \newline öffentlichen Jugendhilfe übernehmen. \newline Dies beinhaltet gemäß \S{} 21 des Neunten \newline Buches auch die Durchführung des \newline Verfahrens zur Gesamtplanung nach \newline den \S{}\S{} 117 bis 122 des Neunten Buches. } & \textbf{(2) Abweichend von Absatz 1 werden \newline bei einem Zuständigkeitsübergang vom \newline Träger der öffentlichen Jugendhilfe auf \newline einen Träger der Eingliederungshilfe \newline rechtzeitig im Rahmen eines \newline Teilhabeplanverfahrens nach \S{} 19 des \newline Neunten Buches die Voraussetzungen \newline für die Sicherstellung einer nahtlosen \newline und bedarfsgerechten \newline Leistungsgewährung nach dem \newline Zuständigkeitsübergang geklärt. Die \newline Teilhabeplanung ist frühzeitig, in der \newline Regel ein Jahr vor dem \newline voraussichtlichen \newline Zuständigkeitswechsel, vom Träger der \newline Jugendhilfe einzuleiten. Mit Zustimmung \newline des Leistungsberechtigten oder seines \newline Personensorgeberechtigten ist eine \newline Teilhabeplankonferenz nach \S{} 20 des \newline Neunten Buches durchzuführen. Stellt \newline der beteiligte Träger der \newline Eingliederungshilfe fest, dass seine } \newline \textbf{Zuständigkeit sowie die \newline Leistungsberechtigung absehbar \newline gegeben sind, soll er entsprechend \S{} 19 \newline Absatz 5 des Neunten Buches die \newline Teilhabeplanung vom Träger der \newline öffentlichen Jugendhilfe übernehmen. \newline Dies beinhaltet gemäß \S{} 21 des Neunten \newline Buches auch die Durchführung des \newline Verfahrens zur Gesamtplanung nach \newline den \S{}\S{} 117 bis 122 des Neunten Buches. \newline \S{} 41a findet keine Anwendung. } \\
\hline
 & \textbf{(3) Werden Hilfen oder Leistungen \newline abweichend von den Absätzen 1 und 2 \newline vom Leistungsberechtigten selbst \newline beschafft, so ist der Träger der \newline öffentlichen Jugendhilfe zur Übernahme \newline der erforderlichen Aufwendungen nur \newline verpflichtet, wenn } \newline \textbf{a) \newline bis zu einer Entscheidung des Trägers \newline der öffentlichen Jugendhilfe über die \newline Gewährung der Leistung oder } \newline \textbf{b) \newline bis zu einer Entscheidung über ein \newline Rechtsmittel nach einer zu Unrecht \newline abgelehnten Leistung keinen zeitlichen \newline Aufschub geduldet hat. War es dem \newline Leistungsberechtigten unmöglich, den \newline Träger der öffentlichen Jugendhilfe \newline rechtzeitig über den Bedarf in Kenntnis \newline zu setzen, so hat er dies unverzüglich \newline nach Wegfall des Hinderungsgrundes \newline nachzuholen. } \newline \textbf{\S{} 36 d } \newline \textbf{Zusammenarbeit beim \newline Zuständigkeitsübergang } & \textbf{(3) Werden Hilfen oder Leistungen \newline abweichend von den Absätzen 1 und 2 \newline vom Leistungsberechtigten selbst \newline beschafft, so ist der Träger der \newline öffentlichen Jugendhilfe zur Übernahme \newline der erforderlichen Aufwendungen nur \newline verpflichtet, wenn } \newline \textbf{a) \newline bis zu einer Entscheidung des Trägers \newline der öffentlichen Jugendhilfe über die \newline Gewährung der Leistung oder } \newline \textbf{b) \newline bis zu einer Entscheidung über ein \newline Rechtsmittel nach einer zu Unrecht \newline abgelehnten Leistung keinen zeitlichen \newline Aufschub geduldet hat. War es dem \newline Leistungsberechtigten unmöglich, den \newline Träger der öffentlichen Jugendhilfe \newline rechtzeitig über den Bedarf in Kenntnis } \newline \textbf{zu setzen, so hat er dies unverzüglich \newline nach Wegfall des Hinderungsgrundes \newline nachzuholen. } \newline \textbf{\S{} 36 d } \newline \textbf{Zusammenarbeit beim \newline Zuständigkeitsübergang } \\
\hline
 & \textbf{1. \newline der Leistungsberechtigte den Träger der \newline öffentlichen Jugendhilfe vor der \newline Selbstbeschaffung über den Bedarf in \newline Kenntnis gesetzt hat, } & \textbf{1. \newline der Leistungsberechtigte den Träger der \newline öffentlichen Jugendhilfe vor der \newline Selbstbeschaffung über den Bedarf in \newline Kenntnis gesetzt hat, } \\
\hline
 & \textbf{2. \newline die Voraussetzungen für die Gewährung \newline der Hilfe oder Leistung vorlagen und } & \textbf{2. \newline die Voraussetzungen für die Gewährung \newline der Hilfe oder Leistung vorlagen und } \\
\hline
 & \textbf{3. \newline die Deckung des Bedarfs } & \textbf{3. \newline die Deckung des Bedarfs } \\
\hline
 & \textbf{(4) Werden bei der Durchführung der \newline Hilfe oder Leistung andere Personen, \newline Dienste oder Einrichtungen tätig, so \newline sind sie oder deren Mitarbeiterinnen \newline und Mitarbeiter an der Aufstellung des \newline Hilfe- und Leistungsplans und seiner \newline Überprüfung zu beteiligen. Soweit dies \newline zur Feststellung des Bedarfs, der zu \newline gewährenden Art der Hilfe oder Leistung \newline oder von deren notwendigen Leistungen \newline nach Inhalt, Umfang und Dauer \newline erforderlich ist, sollen öffentliche \newline Stellen, insbesondere andere \newline Sozialleistungsträger, \newline Rehabilitationsträger oder die Schule \newline beteiligt werden. } & \textbf{(4) Werden bei der Durchführung der \newline Hilfe oder Leistung andere Personen, \newline Dienste oder Einrichtungen tätig, so \newline sind sie oder deren Mitarbeitende an der \newline Aufstellung des Hilfe- und \newline Leistungsplans und seiner Überprüfung \newline zu beteiligen. Soweit dies zur \newline Feststellung des Bedarfs, der zu \newline gewährenden Art der Hilfe oder Leistung \newline oder von deren notwendigen Leistungen \newline nach Inhalt, Umfang und Dauer \newline erforderlich ist, sollen öffentliche \newline Stellen, insbesondere andere \newline Sozialleistungsträger, \newline Rehabilitationsträger oder die Schule \newline beteiligt werden. } \\
\hline
 & \textbf{(5) Soweit dies zur Feststellung des \newline Bedarfs, der zu gewährenden Art der \newline Hilfe oder Leistung oder von deren \newline notwendigen Leistungen nach Inhalt, \newline Umfang und Dauer erforderlich ist und \newline dadurch der Hilfe- oder Leistungszweck \newline nicht in Frage gestellt wird, sollen \newline Eltern, die nicht \newline personensorgeberechtigt sind, an der \newline Aufstellung des Hilfe- und \newline Leistungsplans und seiner Überprüfung \newline beteiligt werden; die Entscheidung, ob, \newline wie und in welchem Umfang deren \newline Beteiligung erfolgt, soll im \newline Zusammenwirken mehrerer Fachkräfte \newline unter Berücksichtigung der \newline Willensäußerung und der Interessen des \newline Kindes oder Jugendlichen sowie der \newline Willensäußerung des \newline Personensorgeberechtigten getroffen \newline werden. } & \textbf{(5) Soweit dies zur Feststellung des \newline Bedarfs, der zu gewährenden Art der \newline Hilfe oder Leistung oder von deren \newline notwendigen Leistungen nach Inhalt, \newline Umfang und Dauer erforderlich ist und \newline dadurch der Hilfe- oder Leistungszweck \newline nicht in Frage gestellt wird, sollen \newline Eltern, die nicht \newline personensorgeberechtigt sind, an der } \newline \textbf{Aufstellung des Hilfe- und \newline Leistungsplans und seiner Überprüfung \newline beteiligt werden; die Entscheidung, ob, \newline wie und in welchem Umfang deren \newline Beteiligung erfolgt, soll im \newline Zusammenwirken mehrerer Fachkräfte \newline unter Berücksichtigung der \newline Willensäußerung und der Interessen des \newline Kindes oder Jugendlichen sowie der \newline Willensäußerung des \newline Personensorgeberechtigten getroffen \newline werden. } \\
\hline
\textit{--- Fassung 2024 ---} \newline Zusammenarbeit beim \newline Zuständigkeitsübergang \newline \textit{--- Fassung 2026 ---} \newline Teilhabeplanung ist frühzeitig, in der Regel \newline ein Jahr vor dem voraussichtlichen \newline Zuständigkeitswechsel, vom Träger der \newline Jugendhilfe einzuleiten. Mit Zustimmung \newline des Leistungsberechtigten oder seines \newline Personensorgeberechtigten ist eine \newline Teilhabeplankonferenz nach \S{} 20 des \newline Neunten Buches durchzuführen. Stellt der \newline beteiligte Träger der Eingliederungshilfe \newline fest, dass seine Zuständigkeit sowie die \newline Leistungsberechtigung absehbar gegeben \newline sind, soll er entsprechend \S{} 19 Absatz 5 \newline des Neunten Buches die Teilhabeplanung \newline vom Träger der öffentlichen Jugendhilfe \newline übernehmen. Dies beinhaltet gemäß \S{} 21 \newline des Neunten Buches auch die \newline Durchführung des Verfahrens zur \newline Gesamtplanung nach den \S{}\S{} 117 bis 122 \newline des Neunten Buches. & \textbf{Hilfe- und Leistungsplankonferenz} & \textbf{der nach \S{} 36a Absatz 4 und 5 an der \newline Aufstellung des Hilfe- und \newline Leistungsplans Beteiligten. Auf \newline Verlangen des \newline Personensorgeberechtigten oder des \newline Kindes oder Jugendlichen wird eine \newline Person seines Vertrauens an der Hilfe- \newline und Leistungsplankonferenz beteiligt.} \newline \textbf{\S{} 36 c } \newline \textbf{Steuerungsverantwortung, \newline Selbstbeschaffung } \\
\hline
 & \textbf{(6) Der Hilfe- und Leistungsplan bedarf \newline der Schriftform. Der Träger der \newline öffentlichen Jugendhilfe stellt ihn dem \newline Leistungsberechtigten zur Verfügung. } & \textbf{(6) Der Hilfe- und Leistungsplan bedarf \newline der Textform. Der Träger der \newline öffentlichen Jugendhilfe stellt ihn dem \newline Leistungsberechtigten zur Verfügung. } \\
\hline
\multicolumn{3}{|c|}{\cellcolor{gray!10}\textbf{\S{} 37 -- Beratung und Unterstützung der Eltern, Zusammenarbeit bei Hilfen außerhalb der eigenen Familie}} \\
\hline
(1) Bei der Aufstellung und Überprüfung \newline des Hilfeplans nach \S{} 36 Absatz 2 Satz 2 \newline ist bei Hilfen außerhalb der eigenen Familie \newline prozesshaft auch die Perspektive der Hilfe \newline zu klären. Der Stand der Perspektivklärung \newline nach Satz 1 ist im Hilfeplan zu \newline dokumentieren. &  &  \\
\hline
(2) Maßgeblich bei der Perspektivklärung \newline nach Absatz 1 ist, ob durch Leistungen \newline nach diesem Abschnitt die Entwicklungs-, \newline Teilhabe- oder Erziehungsbedingungen in \newline der Herkunftsfamilie innerhalb eines im \newline Hinblick auf die Entwicklung des Kindes \newline oder Jugendlichen vertretbaren Zeitraums \newline so weit verbessert werden, dass die \newline Herkunftsfamilie das Kind oder den \newline Jugendlichen wieder selbst erziehen, \newline betreuen und fördern kann. Ist eine \newline nachhaltige Verbesserung der \newline Entwicklungs-, Teilhabe- oder \newline Erziehungsbedingungen in der \newline Herkunftsfamilie innerhalb eines im \newline Hinblick auf die Entwicklung des Kindes \newline oder Jugendlichen vertretbaren Zeitraums \newline nicht erreichbar, so soll mit den beteiligten \newline Personen eine andere, dem Wohl des \newline Kindes oder Jugendlichen förderliche und \newline auf Dauer angelegte Lebensperspektive \newline erarbeitet werden. In diesem Fall ist vor \newline und während der Gewährung der Hilfe \newline insbesondere zu prüfen, ob die Annahme \newline als Kind in Betracht kommt. &  &  \\
\hline
\textbf{Ergänzende Bestimmungen zur \newline Hilfeplanung bei Hilfen außerhalb der \newline eigenen Familie } & entfällt & entfällt \\
\hline
(3) Bei der Auswahl der Einrichtung oder \newline der Pflegeperson sind der \newline Personensorgeberechtigte und das Kind \newline oder der Jugendliche oder bei Hilfen nach \newline \S{} 41 der junge Volljährige zu beteiligen. \newline Der Wahl und den Wünschen des \newline Leistungsberechtigten ist zu entsprechen, \newline sofern sie nicht mit unverhältnismäßigen \newline Mehrkosten verbunden sind. Wünschen die \newline in Satz 1 genannten Personen die \newline Erbringung einer in \S{} 78a genannten \newline Leistung in einer Einrichtung, mit deren \newline Träger keine Vereinbarungen nach \S{} 78b \newline bestehen, so soll der Wahl nur entsprochen \newline werden, wenn die Erbringung der Leistung \newline in dieser Einrichtung nach Maßgabe des \newline Hilfeplans geboten ist. Bei der Auswahl \newline einer Pflegeperson, die ihren gewöhnlichen \newline Aufenthalt außerhalb des Bereichs des \newline örtlich zuständigen Trägers hat, soll der \newline örtliche Träger der öffentlichen Jugendhilfe \newline beteiligt werden, in dessen Bereich die \newline Pflegeperson ihren gewöhnlichen \newline Aufenthalt hat. &  &  \\
\hline
(4) Die Art und Weise der Zusammenarbeit \newline nach \S{} 37 Absatz 2 sowie die damit im \newline Einzelfall verbundenen Ziele sind im \newline Hilfeplan zu dokumentieren. Bei Hilfen \newline nach den \S{}\S{} 33, 35a Absatz 2 Nummer 3 \newline zählen dazu auch der vereinbarte Umfang \newline der Beratung und Unterstützung der Eltern \newline nach \S{} 37 Absatz 1 und der Pflegeperson \newline nach \S{} 37a Absatz 1 sowie die Höhe der \newline laufenden Leistungen zum Unterhalt des \newline Kindes oder Jugendlichen nach \S{} 39. Bei \newline Hilfen für junge Volljährige nach \S{} 41 gilt \newline dies entsprechend in Bezug auf den \newline vereinbarten Umfang der Beratung und \newline Unterstützung der Pflegeperson sowie die \newline Höhe der laufenden Leistungen zum \newline Unterhalt. Eine Abweichung von den im \newline Hilfeplan gemäß den Sätzen 1 bis 3 \newline getroffenen Feststellungen ist nur bei einer \newline Änderung des Hilfebedarfs und \newline entsprechender Änderung des Hilfeplans \newline auch bei einem Wechsel der örtlichen \newline Zuständigkeit zulässig. &  &  \\
\hline
\multicolumn{3}{|c|}{\cellcolor{gray!10}\textbf{\S{} 38 -- Zulässigkeit von Auslandsmaßnahmen}} \\
\hline
 & \textbf{(1) Bei Leistungen der \newline Eingliederungshilfe nach \S{}\S{} 27 Absatz \newline 3, 35a sind die Ergebnisse der \newline Bedarfsermittlung nach \S{} 38b \newline Grundlage für die Beratungen in einer \newline Hilfe- und Leistungsplankonferenz. } & \textbf{(1) Bei Leistungen der \newline Eingliederungshilfe nach \S{}\S{} 27 Absatz \newline 3, 35a sind die Ergebnisse der \newline Bedarfsermittlung nach \S{} 38b \newline Grundlage für die Beratungen in einer \newline Hilfe- und Leistungsplankonferenz. } \\
\hline
b) \newline des Haager Übereinkommens vom 19. \newline Oktober 1996 über die Zuständigkeit, das \newline anzuwendende Recht, die Anerkennung, \newline Vollstreckung und Zusammenarbeit auf \newline dem Gebiet der elterlichen Verantwortung \newline und der Maßnahmen zum Schutz von \newline Kindern zur Erfüllung der Maßgaben des \newline Artikels 33 zu übermitteln. Die \newline erlaubniserteilende Behörde wirkt auf die \newline unverzügliche Beendigung der \newline Leistungserbringung im Ausland hin, wenn \newline sich aus den Angaben nach Satz 1 ergibt, \newline dass die an die Leistungserbringung im \newline Ausland gestellten gesetzlichen \newline Anforderungen nicht erfüllt sind. & b) entfällt \newline \S{} 38 a \newline \textbf{Bedarfsfeststellung bei Leistungen der \newline Eingliederungshilfe für Kinder und \newline Jugendliche mit Behinderungen } & b) entfällt \newline \S{} 38 a \newline \textbf{Bedarfsfeststellung bei Leistungen der \newline Eingliederungshilfe für Kinder und \newline Jugendliche mit Behinderungen } \\
\hline
 & \textbf{(2) In eine Hilfe- und \newline Leistungsplankonferenz können die \newline nach \S{} 38c Absatz 3 und 4 an der \newline Aufstellung und Überprüfung des Hilfe- \newline und Leistungsplans Beteiligten \newline einbezogen werden. } & \textbf{(2) In eine Hilfe- und \newline Leistungsplankonferenz sollen die nach \newline \S{} 38c Absatz 3 und 4 an der Aufstellung \newline und Überprüfung des Hilfe- und \newline Leistungsplans Beteiligten einbezogen \newline werden. } \\
\hline
 & \textbf{(3) Ist der Träger der öffentlichen \newline Jugendhilfe Leistungsverantwortlicher \newline nach \S{} 15 des Neunten Buches, soll er \newline die Hilfe- und Leistungsplankonferenz \newline mit einer Teilhabeplankonferenz nach \S{} \newline 20 des Neunten Buches verbinden. Ist \newline der Träger der öffentlichen Jugendhilfe \newline beteiligter Rehabilitationsträger nach \S{} \newline 15 des Neunten Buches soll er dem \newline Leistungsberechtigten und den anderen \newline Rehabilitationsträgern anbieten, mit \newline deren Einvernehmen das Verfahren \newline anstelle des leistenden \newline Rehabilitationsträgers durchzuführen; \newline die Vorschriften über die \newline Leistungsverantwortung der \newline Rehabilitationsträger nach den \S{}\S{} 14 \newline und 15 des Neunten Buches bleiben \newline hiervon unberührt. } & \textbf{(3) Ist der Träger der öffentlichen \newline Jugendhilfe Leistungsverantwortlicher \newline nach \S{} 15 des Neunten Buches, soll er \newline die Hilfe- und Leistungsplankonferenz \newline mit einer Teilhabeplankonferenz nach \S{} \newline 20 des Neunten Buches verbinden. Ist \newline der Träger der öffentlichen Jugendhilfe \newline beteiligter Rehabilitationsträger nach \S{} \newline 15 des Neunten Buches, soll er dem \newline Leistungsberechtigten und den anderen \newline Rehabilitationsträgern anbieten, mit \newline deren Einvernehmen das Verfahren \newline anstelle des leistenden \newline Rehabilitationsträgers durchzuführen; \newline die Vorschriften über die \newline Leistungsverantwortung der \newline Rehabilitationsträger nach den \S{}\S{} 14 \newline und 15 des Neunten Buches bleiben \newline hiervon unberührt. } \\
\hline
 & \textbf{(4) Bestehen im Einzelfall Anhaltspunkte \newline für eine Pflegebedürftigkeit nach dem \newline Elften Buch, wird die zuständige \newline Pflegekasse mit Zustimmung des \newline Leistungsberechtigten vom Träger der \newline öffentlichen Jugendhilfe informiert und \newline muss an der Aufstellung und \newline Überprüfung des Hilfe- und \newline Leistungsplans beratend teilnehmen, \newline soweit dies zur Feststellung des \newline Bedarfs, der zu gewährenden Art der \newline Leistung oder von deren not-wendiger \newline Ausgestaltung nach Inhalt, Umfang und \newline Dauer erforderlich ist. Bestehen im \newline Einzelfall Anhaltspunkte, dass \newline Leistungen der Hilfe zur Pflege nach \newline dem Siebten Kapitel des Zwölften \newline Buches erforderlich sind, so soll der \newline Träger dieser Leistungen mit Zu- \newline stimmung der Leistungsberechtigten \newline informiert und an der Aufstellung und \newline Überprüfung des Hilfe- und \newline Leistungsplans beteiligt werden, soweit \newline dies zur Feststellung des Bedarfs, der \newline zu gewährenden Art der Leistung oder \newline von deren notwendiger Ausgestaltung \newline nach Inhalt, Umfang und Dauer \newline erforderlich ist. } & \textbf{(4) Bestehen im Einzelfall Anhaltspunkte \newline für eine Pflegebedürftigkeit nach dem \newline Elften Buch, wird die zuständige \newline Pflegekasse mit Zustimmung des \newline Leistungsberechtigten vom Träger der \newline öffentlichen Jugendhilfe informiert und \newline muss an der Aufstellung und \newline Überprüfung des Hilfe- und \newline Leistungsplans beratend teilnehmen, \newline soweit dies zur Feststellung des \newline Bedarfs, der zu gewährenden Art der \newline Leistung oder von deren notwendiger \newline Ausgestaltung nach Inhalt, Umfang und \newline Dauer erforderlich ist. Bestehen im \newline Einzelfall Anhaltspunkte, dass \newline Leistungen der Hilfe zur Pflege nach \newline dem Siebten Kapitel des Zwölften \newline Buches erforderlich sind, so soll der \newline Träger dieser Leistungen mit Zu- \newline stimmung der Leistungsberechtigten \newline informiert und an der Aufstellung und \newline Überprüfung des Hilfe- und \newline Leistungsplans beteiligt werden, soweit \newline dies zur Feststellung des Bedarfs, der \newline zu gewährenden Art der Leistung oder \newline von deren notwendiger Ausgestaltung \newline nach Inhalt, Umfang und Dauer \newline erforderlich ist. } \\
\hline
 & \textbf{(5) Soweit Leistungen verschiedener \newline Leistungsgruppen nach \S{} 5 des Neunten \newline Buches oder mehrere \newline Rehabilitationsträger nach \S{} 6 Absatz 1 \newline des Neunten Buches erforderlich sind, \newline hat der Träger der öffentlichen \newline Jugendhilfe als leistender \newline Rehabilitationsträger die Regelungen \newline zum Teilhabeplan nach \S{} 19 des \newline Neunten Buches und legt ihn seiner \newline Entscheidung über die Gewährung einer \newline Leistung der Eingliederungshilfe \newline zugrunde. Im Übrigen ist \S{} 19 des \newline Neunten Buches anzuwenden. } \newline \textbf{\S{} 38 d } \newline \textbf{Ergänzende Bestimmungen zur Hilfe- \newline und Leistungsplankonferenz bei \newline Leistungen der Eingliederungshilfe für \newline Kinder und Jugendliche mit \newline Behinderungen } & \textbf{(5) Soweit Leistungen verschiedener \newline Leistungsgruppen nach \S{} 5 des Neunten \newline Buches oder mehrere \newline Rehabilitationsträger nach \S{} 6 Absatz 1 \newline des Neunten Buches erforderlich sind, \newline oder der Leistungsberechtigte oder der \newline Personensorgeberechtigte dies \newline wünscht, hat der Träger der öffentlichen \newline Jugendhilfe als leistender \newline Rehabilitationsträger die Regelungen } \newline \textbf{zum Teilhabeplan nach \S{} 19 des \newline Neunten Buches anzuwenden und legt \newline diesen seiner Entscheidung über die \newline Gewährung einer Leistung der \newline Eingliederungshilfe zugrunde. Im \newline Übrigen gilt \S{} 19 des Neunten Buches. } \newline \textbf{\S{} 38 d } \newline \textbf{Besondere Bestimmungen zur Hilfe- und \newline Leistungsplankonferenz bei Leistungen \newline der Eingliederungshilfe für Kinder und \newline Jugendliche mit Behinderungen } \\
\hline
\multicolumn{3}{|c|}{\cellcolor{gray!10}\textbf{\S{} 39 -- Leistungen zum Unterhalt des Kindes oder des Jugendlichen}} \\
\hline
 & \textbf{(1) Wird Hilfe nach den \S{}\S{} 32 bis 35 oder \newline eine Leistung nach \S{} 35a Absatz 4 \newline Nummer 2 bis 4 gewährt, so ist auch der \newline notwendige Unterhalt des Kindes oder \newline Jugendlichen außerhalb des \newline Elternhauses sicherzustellen. Er \newline umfasst die Kosten für den \newline Sachaufwand sowie für die Pflege und \newline Erziehung des Kindes oder \newline Jugendlichen. } & \textbf{(1) Wird Hilfe nach den \S{}\S{} 32 bis 35 oder \newline eine Leistung nach \S{} 35a Absatz 4 \newline Nummer 2 bis 4 gewährt, so ist auch der \newline notwendige Unterhalt des Kindes oder \newline Jugendlichen außerhalb des \newline Elternhauses sicherzustellen. Er \newline umfasst die Kosten für den \newline Sachaufwand sowie für die Pflege und \newline Erziehung des Kindes oder \newline Jugendlichen. } \\
\hline
 & \textbf{(2) Der gesamte regelmäßig \newline wiederkehrende Bedarf soll durch \newline laufende Leistungen gedeckt werden. \newline Sie umfassen außer im Fall des \S{} 32 und \newline des \S{} 35a Absatz 4 Nummer 2 auch \newline einen angemessenen Barbetrag zur \newline persönlichen Verfügung des Kindes \newline oder des Jugendlichen. Die Höhe des \newline Betrages wird in den Fällen der \S{}\S{} 34, \newline 35, 35a Absatz 4 Nummer 4 von der \newline nach Landesrecht zuständigen Behörde \newline festgesetzt; die Beträge sollen nach \newline Altersgruppen gestaffelt sein. Die \newline laufenden Leistungen im Rahmen der \newline Hilfe in Vollzeitpflege (\S{} 33) oder bei \newline einer geeigneten Pflegeperson (\S{} 35a \newline Absatz 4 Nummer 3) sind nach den \newline Absätzen 4 bis 6 zu bemessen. } & \textbf{(2) Der gesamte regelmäßig \newline wiederkehrende Bedarf soll durch \newline laufende Leistungen gedeckt werden. \newline Sie umfassen außer im Fall des \S{} 32 und \newline des \S{} 35a Absatz 4 Nummer 2 auch \newline einen angemessenen Barbetrag zur \newline persönlichen Verfügung des Kindes \newline oder des Jugendlichen. Die Höhe des \newline Betrages wird in den Fällen der \S{}\S{} 34, \newline 35, 35a Absatz 4 Nummer 4 von der \newline nach Landesrecht zuständigen Behörde \newline festgesetzt; die Beträge sollen nach \newline Altersgruppen gestaffelt sein. Die \newline laufenden Leistungen im Rahmen der \newline Hilfe in Vollzeitpflege (\S{} 33) oder bei \newline einer geeigneten Pflegeperson (\S{} 35a \newline Absatz 4 Nummer 3) sind nach den \newline Absätzen 4 bis 6 zu bemessen. } \\
\hline
 & \textbf{(3) Einmalige Beihilfen oder Zuschüsse \newline können insbesondere zur \newline Erstausstattung einer Pflegestelle, bei \newline wichtigen persönlichen Anlässen sowie \newline für Urlaubs- und Ferienreisen des \newline Kindes oder des Jugendlichen gewährt \newline werden. } & \textbf{(3) Einmalige Beihilfen oder Zuschüsse \newline können insbesondere zur \newline Erstausstattung einer Pflegestelle, bei \newline wichtigen persönlichen Anlässen sowie \newline für Urlaubs- und Ferienreisen des \newline Kindes oder des Jugendlichen gewährt \newline werden. } \\
\hline
persönlichen Anlässen sowie für Urlaubs- \newline und Ferienreisen des Kindes oder des \newline Jugendlichen gewährt werden. & \textit{unverändert} & \textbf{Gesetzbuchs die \newline Entscheidungsbefugnisse der \newline Pflegeperson so weit einschränkt, dass \newline die Einschränkung eine dem Wohl des \newline Kindes oder des Jugendlichen \newline förderliche Entwicklung nicht mehr \newline ermöglicht, sollen die Beteiligten das \newline Jugendamt einschalten. Auch bei \newline sonstigen Meinungsverschiedenheiten \newline zwischen ihnen sollen die Beteiligten \newline das Jugendamt einschalten. } \\
\hline
 & \textbf{(4) Die laufenden Leistungen sollen auf \newline der Grundlage der tatsächlichen Kosten \newline gewährt werden, sofern sie einen \newline angemessenen Umfang nicht \newline übersteigen. Die laufenden Leistungen \newline umfassen auch die Erstattung \newline nachgewiesener Aufwendungen für \newline Beiträge zu einer Unfallversicherung \newline sowie die hälftige Erstattung \newline nachgewiesener Aufwendungen zu \newline einer angemessenen Alterssicherung \newline der Pflegeperson. Sie sollen in einem \newline monatlichen Pauschalbetrag gewährt \newline werden, soweit nicht nach der \newline Besonderheit des Einzelfalls \newline abweichende Leistungen geboten sind. \newline Ist die Pflegeperson in gerader Linie mit \newline dem Kind oder Jugendlichen verwandt \newline und kann sie diesem unter \newline Berücksichtigung ihrer sonstigen \newline Verpflichtungen und ohne Gefährdung \newline ihres angemessenen Unterhalts \newline Unterhalt gewähren, so kann der Teil \newline des monatlichen Pauschalbetrages, der \newline die Kosten für den Sachaufwand des \newline Kindes oder Jugendlichen betrifft, \newline angemessen gekürzt werden. Wird ein \newline Kind oder ein Jugendlicher im Bereich \newline eines anderen Jugendamts \newline untergebracht, so soll sich die Höhe des \newline zu gewährenden Pauschalbetrages nach \newline den Verhältnissen richten, die am Ort \newline der Pflegestelle gelten. } & \textbf{(4) Die laufenden Leistungen sollen auf \newline der Grundlage der tatsächlichen Kosten \newline gewährt werden, sofern sie einen \newline angemessenen Umfang nicht \newline übersteigen. Die laufenden Leistungen \newline umfassen auch die Erstattung \newline nachgewiesener Aufwendungen für \newline Beiträge zu einer Unfallversicherung \newline sowie die hälftige Erstattung \newline nachgewiesener Aufwendungen zu \newline einer angemessenen Alterssicherung \newline der Pflegeperson. Sie sollen in einem \newline monatlichen Pauschalbetrag gewährt \newline werden, soweit nicht nach der \newline Besonderheit des Einzelfalls } \newline \textbf{abweichende Leistungen geboten sind. \newline Ist die Pflegeperson in gerader Linie mit \newline dem Kind oder Jugendlichen verwandt \newline und kann sie diesem unter \newline Berücksichtigung ihrer sonstigen \newline Verpflichtungen und ohne Gefährdung \newline ihres angemessenen Unterhalts \newline Unterhalt gewähren, so kann der Teil \newline des monatlichen Pauschalbetrages, der \newline die Kosten für den Sachaufwand des \newline Kindes oder Jugendlichen betrifft, \newline angemessen gekürzt werden. Wird ein \newline Kind oder ein Jugendlicher im Bereich \newline eines anderen Jugendamts \newline untergebracht, so soll sich die Höhe des \newline zu gewährenden Pauschalbetrages nach \newline den Verhältnissen richten, die am Ort \newline der Pflegestelle gelten. } \\
\hline
 & \textbf{(5) Die Pauschalbeträge für laufende \newline Leistungen zum Unterhalt sollen von \newline den nach Landesrecht zuständigen \newline Behörden festgesetzt werden. Dabei ist \newline dem altersbedingt unterschiedlichen \newline Unterhaltsbedarf von Kindern und \newline Jugendlichen durch eine Staffelung der \newline Beträge nach Altersgruppen Rechnung \newline zu tragen. Das Nähere regelt \newline Landesrecht. } & \textbf{(5) Die Pauschalbeträge für laufende \newline Leistungen zum Unterhalt sollen von \newline den nach Landesrecht zuständigen \newline Behörden festgesetzt werden. Dabei ist \newline dem altersbedingt unterschiedlichen \newline Unterhaltsbedarf von Kindern und \newline Jugendlichen durch eine Staffelung der \newline Beträge nach Altersgruppen Rechnung \newline zu tragen. Das Nähere regelt \newline Landesrecht. } \\
\hline
 & \textbf{(6) Wird das Kind oder der Jugendliche \newline im Rahmen des \newline Familienleistungsausgleichs nach \S{} 31 \newline des Einkommensteuergesetzes bei der \newline Pflegeperson berücksichtigt, so ist ein \newline Betrag in Höhe der Hälfte des Betrages, \newline der nach \S{} 66 des \newline Einkommensteuergesetzes für ein \newline erstes Kind zu zahlen ist, auf die \newline laufenden Leistungen anzurechnen. Ist \newline das Kind oder der Jugendliche nicht das \newline älteste Kind in der Pflegefamilie, so \newline ermäßigt sich der Anrechnungsbetrag \newline für dieses Kind oder diesen \newline Jugendlichen auf ein Viertel des \newline Betrages, der für ein erstes Kind zu \newline zahlen ist. } & \textbf{(6) Wird das Kind oder der Jugendliche \newline im Rahmen des \newline Familienleistungsausgleichs nach \S{} 31 \newline des Einkommensteuergesetzes bei der \newline Pflegeperson berücksichtigt, so ist ein \newline Betrag in Höhe der Hälfte des Betrages, \newline der nach \S{} 66 des \newline Einkommensteuergesetzes für ein \newline erstes Kind zu zahlen ist, auf die \newline laufenden Leistungen anzurechnen. Ist \newline das Kind oder der Jugendliche nicht das \newline älteste Kind in der Pflegefamilie, so \newline ermäßigt sich der Anrechnungsbetrag \newline für dieses Kind oder diesen \newline Jugendlichen auf ein Viertel des \newline Betrages, der für ein erstes Kind zu \newline zahlen ist. } \\
\hline
 & \textbf{(7) Wird ein Kind oder eine Jugendliche \newline während ihres Aufenthalts in einer \newline Einrichtung oder einer Pflegefamilie \newline selbst Mutter eines Kindes, so ist auch \newline der notwendige Unterhalt dieses Kindes \newline sicherzustellen. } \newline \textbf{\S{} 39 d } \newline \textbf{Krankenhilfe } \newline \textbf{Wird Hilfe nach den \S{}\S{} 33 bis 35 oder \newline eine Leistung nach \S{} 35a Absatz 4 \newline Nummer 3 oder 4 gewährt, so ist auch \newline Krankenhilfe zu leisten; für den Umfang \newline der Hilfe gelten die \S{}\S{} 47 bis 52 des \newline Zwölften Buches entsprechend. \newline Krankenhilfe muss den im Einzelfall \newline notwendigen Bedarf in voller Höhe \newline befriedigen. Zuzahlungen und \newline Eigenbeteiligungen sind zu \newline übernehmen. Das Jugendamt kann in \newline geeigneten Fällen die Beiträge für eine \newline freiwillige Krankenversicherung \newline übernehmen, soweit sie angemessen \newline sind. } & \textbf{(7) Wird ein Kind oder eine Jugendliche \newline während ihres Aufenthalts in einer \newline Einrichtung oder einer Pflegefamilie \newline selbst Mutter eines Kindes, so ist auch } \newline \textbf{der notwendige Unterhalt dieses Kindes \newline sicherzustellen. } \newline \textbf{\S{} 39 d } \newline \textbf{Krankenhilfe } \newline \textbf{Wird Hilfe nach den \S{}\S{} 33 bis 35 oder \newline eine Leistung nach \S{} 35a Absatz 4 \newline Nummer 3 oder 4 gewährt, so ist auch \newline Krankenhilfe zu leisten; für den Umfang \newline der Hilfe gelten die \S{}\S{} 47 bis 52 des \newline Zwölften Buches entsprechend. \newline Krankenhilfe muss den im Einzelfall \newline notwendigen Bedarf in voller Höhe \newline befriedigen. Zuzahlungen und \newline Eigenbeteiligungen sind zu \newline übernehmen. Das Jugendamt kann in \newline geeigneten Fällen die Beiträge für eine \newline freiwillige Krankenversicherung \newline übernehmen, soweit sie angemessen \newline sind. } \\
\hline
\multicolumn{3}{|c|}{\cellcolor{gray!10}\textbf{\S{} 40 -- Krankenhilfe}} \\
\hline
Wird Hilfe nach den \S{}\S{} 33 bis 35 oder nach \newline \S{} 35a Absatz 2 Nummer 3 oder 4 gewährt, \newline so ist auch Krankenhilfe zu leisten; für den \newline Umfang der Hilfe gelten die \S{}\S{} 47 bis 52 \newline des Zwölften Buches entsprechend. \newline Krankenhilfe muss den im Einzelfall \newline notwendigen Bedarf in voller Höhe \newline befriedigen. Zuzahlungen und \newline Eigenbeteiligungen sind zu übernehmen. \newline Das Jugendamt kann in geeigneten Fällen \newline die Beiträge für eine freiwillige \newline Krankenversicherung übernehmen, soweit \newline sie angemessen sind. & \textbf{(1) Hilfen oder Leistungen nach diesem \newline Abschnitt sind in der Regel im Inland zu \newline erbringen. Sie dürfen nur dann im \newline Ausland erbracht werden, wenn dies \newline nach Maßgabe der Hilfe- und \newline Leistungsplanung zur Erreichung des \newline Hilfezieles im Einzelfall erforderlich ist \newline und die aufenthaltsrechtlichen \newline Vorschriften des aufnehmenden Staates \newline sowie } & \textbf{(1) Hilfen oder Leistungen nach diesem \newline Abschnitt sind in der Regel im Inland zu \newline erbringen. Sie dürfen nur dann im \newline Ausland erbracht werden, wenn dies \newline nach Maßgabe der Hilfe- und \newline Leistungsplanung zur Erreichung des \newline Hilfezieles im Einzelfall erforderlich ist \newline und die aufenthaltsrechtlichen \newline Vorschriften des aufnehmenden Staates \newline sowie } \newline \textbf{2.7.2019, S. 1) die Voraussetzungen des \newline Artikels 82 oder } \\
\hline
 & \textbf{1. \newline im Anwendungsbereich der Verordnung \newline (EU) 2019/1111 des Rates vom 25. Juni \newline 2019 über die Zuständigkeit, die \newline Anerkennung und Vollstreckung von \newline Entscheidungen in Ehesachen und in \newline Verfahren betreffend die elterliche \newline Verantwortung und über internationale \newline Kindesentführungen (ABl. L 178 vom \newline 2.7.2019, S. 1) die Voraussetzungen des \newline Artikels 82 oder } & \textbf{1. \newline im Anwendungsbereich der Verordnung \newline (EU) 2019/1111 des Rates vom 25. Juni \newline 2019 über die Zuständigkeit, die \newline Anerkennung und Vollstreckung von \newline Entscheidungen in Ehesachen und in \newline Verfahren betreffend die elterliche \newline Verantwortung und über internationale \newline Kindesentführungen (ABl. L 178 vom } \\
\hline
 & \textbf{2. \newline im Anwendungsbereich des Haager \newline Übereinkommens vom 19. Oktober 1996 \newline über die Zuständigkeit, das \newline anzuwendende Recht, die Anerkennung, \newline Vollstreckung und Zusammenarbeit auf \newline dem Gebiet der elterlichen \newline Verantwortung und der Maßnahmen \newline zum Schutz von Kindern die \newline Voraussetzungen des Artikels 33 erfüllt \newline sind. } & \textbf{2. \newline im Anwendungsbereich des Haager \newline Übereinkommens vom 19. Oktober 1996 \newline über die Zuständigkeit, das \newline anzuwendende Recht, die Anerkennung, \newline Vollstreckung und Zusammenarbeit auf \newline dem Gebiet der elterlichen \newline Verantwortung und der Maßnahmen \newline zum Schutz von Kindern die \newline Voraussetzungen des Artikels 33 erfüllt \newline sind. } \\
\hline
 & \textbf{(2) Der Träger der öffentlichen \newline Jugendhilfe soll vor der Entscheidung \newline über die Gewährung einer Hilfe, die \newline ganz oder teilweise im Ausland erbracht \newline wird, } \newline \textbf{a) \newline über eine Betriebserlaubnis nach \S{} 45 \newline für eine Einrichtung im Inland verfügt, in \newline der Hilfe zur Erziehung erbracht wird, } \newline \textbf{b) \newline Gewähr dafür bietet, dass er die \newline Rechtsvorschriften des aufnehmenden \newline Staates einschließlich des \newline Aufenthaltsrechts einhält, insbesondere \newline vor Beginn der Leistungserbringung die \newline in Absatz 1 Satz 2 genannten Maßgaben \newline erfüllt, und mit den Behörden des \newline aufnehmenden Staates sowie den \newline deutschen Vertretungen im Ausland \newline zusammenarbeitet, } \newline \textbf{c) \newline mit der Erbringung der Hilfen nur \newline Fachkräfte nach \S{} 72 Absatz 1 betraut, } \newline \textbf{d) \newline über die Qualität der Maßnahme eine \newline Vereinbarung abschließt; dabei sind die \newline fachlichen Handlungsleitlinien des \newline überörtlichen Trägers anzuwenden, } \newline \textbf{e) \newline Ereignisse oder Entwicklungen, die \newline geeignet sind, das Wohl des Kindes \newline oder Jugendlichen zu beeinträchtigen, \newline dem Träger der öffentlichen Jugendhilfe \newline unverzüglich anzeigt und } & \textbf{(2) Der Träger der öffentlichen \newline Jugendhilfe soll vor der Entscheidung \newline über die Gewährung einer Hilfe, die \newline ganz oder teilweise im Ausland erbracht \newline wird, } \newline \textbf{a) \newline über eine Betriebserlaubnis nach \S{} 45 \newline für eine Einrichtung im Inland verfügt, in \newline der Hilfe zur Erziehung oder Leistungen \newline der Eingliederungshilfe erbracht \newline werden, } \newline \textbf{b) } \newline \textbf{Gewähr dafür bietet, dass er die \newline Rechtsvorschriften des aufnehmenden \newline Staates einschließlich des \newline Aufenthaltsrechts einhält, insbesondere \newline vor Beginn der Leistungserbringung die \newline in Absatz 1 Satz 2 genannten Maßgaben \newline erfüllt, und mit den Behörden des \newline aufnehmenden Staates sowie den \newline deutschen Vertretungen im Ausland \newline zusammenarbeitet, } \newline \textbf{c) \newline mit der Erbringung der Hilfen oder \newline Leistungen nur Fachkräfte nach \S{} 72 \newline Absatz 1 betraut, } \newline \textbf{d) \newline über die Qualität der Maßnahme eine \newline Vereinbarung abschließt; dabei sind die \newline fachlichen Handlungsleitlinien des \newline überörtlichen Trägers anzuwenden, } \newline \textbf{e) \newline Ereignisse oder Entwicklungen, die \newline geeignet sind, das Wohl des Kindes \newline oder Jugendlichen zu beeinträchtigen, \newline dem Träger der öffentlichen Jugendhilfe \newline unverzüglich anzeigt und } \\
\hline
 & \textbf{1. \newline zur Feststellung einer seelischen \newline Störung mit Krankheitswert die \newline Stellungnahme einer in \S{} 35a Absatz 1a \newline Satz 1 genannten Person einholen, } & \textbf{1. \newline zur Feststellung einer seelischen \newline Störung mit Krankheitswert die \newline Stellungnahme eines Arztes für Kinder- \newline und Jugendpsychiatrie und - \newline psychotherapie, eines Kinder- oder \newline Jugendlichenpsychotherapeuten, eines \newline Fachpsychotherapeuten, eines \newline Psychotherapeuten mit einer \newline Weiterbildung für die Behandlung von \newline Kindern und Jugendlichen oder eines \newline Arztes oder eines psychologischen \newline Psychotherapeuten, der über besondere \newline Erfahrung auf dem Gebiet seelischer \newline Störungen bei Kindern und \newline Jugendlichen verfügt, } \\
\hline
 & \textbf{2. \newline sicherstellen, dass der \newline Leistungserbringer } & \textbf{2. \newline sicherstellen, dass der \newline Leistungserbringer } \\
\hline
 & \textbf{3. \newline die Eignung der mit der \newline Leistungserbringung zu betrauenden \newline Einrichtung oder Person an Ort und \newline Stelle überprüfen. } & \textbf{3. \newline die Eignung der mit der \newline Leistungserbringung zu betrauenden \newline Einrichtung oder Person an Ort und \newline Stelle überprüfen. } \\
\hline
 & \textbf{(3) Überprüfung und Fortschreibung des \newline Hilfeplans sollen nach Maßgabe von \S{} \newline 36 Absatz 2 Satz 2 am Ort der \newline Leistungserbringung unter Beteiligung \newline des Kindes oder des Jugendlichen \newline erfolgen. Unabhängig von der \newline Überprüfung und Fortschreibung des \newline Hilfeplans nach Satz 1 soll der Träger \newline der öffentlichen Jugendhilfe nach den \newline Erfordernissen im Einzelfall an Ort und \newline Stelle überprüfen, ob die Anforderungen \newline nach Absatz 2 Nummer 2 Buchstabe b \newline und c sowie Nummer 3 weiter erfüllt \newline sind. } & \textbf{(3) Überprüfung und Fortschreibung des \newline Hilfe- und Leistungsplans sollen nach \newline Maßgabe von \S{} 36a Absatz 2 Satz 2 am \newline Ort der Leistungserbringung unter \newline Beteiligung des Kindes oder des \newline Jugendlichen erfolgen. Unabhängig von \newline der Überprüfung und Fortschreibung \newline des Hilfeplans nach Satz 1 soll der \newline Träger der öffentlichen Jugendhilfe \newline nach den Erfordernissen im Einzelfall an \newline Ort und Stelle überprüfen, ob die \newline Anforderungen nach Absatz 2 Nummer \newline 2 Buchstabe b und c sowie Nummer 3 \newline weiter erfüllt sind. } \\
\hline
 & \textbf{(4) Besteht die Erfüllung der \newline Anforderungen nach Absatz 2 Nummer \newline 2 oder die Eignung der mit der \newline Leistungserbringung betrauten \newline Einrichtung oder Person nicht fort, soll \newline die Leistungserbringung im Ausland \newline unverzüglich beendet werden. } & \textbf{(4) Besteht die Erfüllung der \newline Anforderungen nach Absatz 2 Nummer \newline 2 oder die Eignung der mit der \newline Leistungserbringung betrauten \newline Einrichtung oder Person nicht fort, soll \newline die Leistungserbringung im Ausland \newline unverzüglich beendet werden. } \\
\hline
 & \textbf{(5) Der Träger der öffentlichen \newline Jugendhilfe hat der erlaubniserteilenden \newline Behörde unverzüglich } \newline \textbf{a) \newline der Verordnung (EU) 2019/1111 zur \newline Erfüllung der Maßgaben des Artikels 82, } \newline \textbf{b) \newline des Haager Übereinkommens vom 19. \newline Oktober 1996 über die Zuständigkeit, \newline das anzuwendende Recht, die \newline Anerkennung, Vollstreckung und \newline Zusammenarbeit auf dem Gebiet der \newline elterlichen Verantwortung und der \newline Maßnahmen zum Schutz von Kindern \newline zur Erfüllung der Maßgaben des Artikels \newline 33 zu übermitteln. Die \newline erlaubniserteilende Behörde wirkt auf \newline die unverzügliche Beendigung der \newline Leistungserbringung im Ausland hin, \newline wenn sich aus den Angaben nach Satz 1 \newline ergibt, dass die an die \newline Leistungserbringung im Ausland \newline gestellten gesetzlichen Anforderungen \newline nicht erfüllt sind. } & \textbf{(5) Der Träger der öffentlichen \newline Jugendhilfe hat der erlaubniserteilenden \newline Behörde unverzüglich } \newline \textbf{a) \newline der Verordnung (EU) 2019/1111 zur \newline Erfüllung der Maßgaben des Artikels 82, } \newline \textbf{b) \newline des Haager Übereinkommens vom 19. \newline Oktober 1996 über die Zuständigkeit, \newline das anzuwendende Recht, die \newline Anerkennung, Vollstreckung und \newline Zusammenarbeit auf dem Gebiet der \newline elterlichen Verantwortung und der \newline Maßnahmen zum Schutz von Kindern \newline zur Erfüllung der Maßgaben des Artikels } \newline \textbf{33 zu übermitteln. Die \newline erlaubniserteilende Behörde wirkt auf \newline die unverzügliche Beendigung der \newline Leistungserbringung im Ausland hin, \newline wenn sich aus den Angaben nach Satz 1 \newline ergibt, dass die an die \newline Leistungserbringung im Ausland \newline gestellten gesetzlichen Anforderungen \newline nicht erfüllt sind. } \\
\hline
 & \textbf{1. \newline den Beginn und das geplante Ende der \newline Leistungserbringung im Ausland unter \newline Angabe von Namen und Anschrift des \newline Leistungserbringers, des \newline Aufenthaltsorts des Kindes oder \newline Jugendlichen sowie der Namen der mit \newline der Erbringung der Hilfe betrauten \newline Fachkräfte, } & \textbf{1. \newline den Beginn und das geplante Ende der \newline Leistungserbringung im Ausland unter \newline Angabe von Namen und Anschrift des \newline Leistungserbringers, des \newline Aufenthaltsorts des Kindes oder \newline Jugendlichen sowie der Namen der mit \newline der Erbringung der Hilfe oder Leistung \newline betrauten Fachkräfte, } \\
\hline
 & \textbf{2. \newline Änderungen der in Nummer 1 \newline bezeichneten Angaben sowie } & \textbf{2. \newline Änderungen der in Nummer 1 \newline bezeichneten Angaben sowie } \\
\hline
 & \textbf{3. \newline die bevorstehende Beendigung der \newline Leistungserbringung im Ausland zu \newline melden sowie } & \textbf{3. \newline die bevorstehende Beendigung der \newline Leistungserbringung im Ausland zu \newline melden sowie } \\
\hline
 & \textbf{4. \newline einen Nachweis zur Erfüllung der \newline aufenthaltsrechtlichen Vorschriften des \newline aufnehmenden Staates und im \newline Anwendungsbereich } & \textbf{4. \newline einen Nachweis zur Erfüllung der \newline aufenthaltsrechtlichen Vorschriften des \newline aufnehmenden Staates und im \newline Anwendungsbereich } \\
\hline
\multicolumn{3}{|c|}{\cellcolor{gray!10}\textbf{\S{} 41 -- Hilfe für junge Volljährige}} \\
\hline
(1) Junge Volljährige erhalten geeignete \newline und notwendige Hilfe nach diesem \newline Abschnitt, wenn und solange ihre \newline Persönlichkeitsentwicklung eine \newline selbstbestimmte, eigenverantwortliche und \newline selbständige Lebensführung nicht \newline gewährleistet. Die Hilfe wird in der Regel \newline nur bis zur Vollendung des 21. \newline Lebensjahres gewährt; in begründeten \newline Einzelfällen soll sie für einen begrenzten \newline Zeitraum darüber hinaus fortgesetzt \newline werden. Eine Beendigung der Hilfe schließt \newline die erneute Gewährung oder Fortsetzung \newline einer Hilfe nach Maßgabe der Sätze 1 und \newline 2 nicht aus. & \textit{(1) unverändert} & \textit{(1) unverändert} \\
\hline
(2) Für die Ausgestaltung der Hilfe gelten \S{} \newline 27 Absatz 3 und 4 sowie die \S{}\S{} 28 bis 30, \newline 33 bis 36, 39 und 40 entsprechend mit der \newline Maßgabe, dass an die Stelle des \newline Personensorgeberechtigten oder des \newline Kindes oder des Jugendlichen der junge \newline Volljährige tritt. & (2) Für die Ausgestaltung der Hilfe gelten \textbf{\S{} \newline 27a Absatz 3 und 4 sowie die \S{}\S{} 28 bis \newline 30, 33 bis 38d, 39c und 39d }entsprechend \newline mit der Maßgabe, dass an die Stelle des \newline Personensorgeberechtigten oder des \newline Kindes oder des Jugendlichen der junge \newline Volljährige tritt. & (2) Für die Ausgestaltung der Hilfe gelten \textbf{\S{} \newline 27a Absatz 3 und 4 sowie die \S{}\S{} 28 bis \newline 30, 33 bis 38d, 39c und 39d }entsprechend \newline mit der Maßgabe, dass an die Stelle des \newline Personensorgeberechtigten oder des \newline Kindes oder des Jugendlichen der junge \newline Volljährige tritt. \\
\hline
(3) Soll eine Hilfe nach dieser Vorschrift \newline nicht fortgesetzt oder beendet werden, \newline prüft der Träger der öffentlichen \newline Jugendhilfe ab einem Jahr vor dem hierfür \newline im Hilfeplan vorgesehenen Zeitpunkt, ob im \newline Hinblick auf den Bedarf des jungen \newline Menschen ein Zuständigkeitsübergang auf \newline andere Sozialleistungsträger in Betracht \newline kommt; \S{} 36b gilt entsprechend. & (3) Soll eine Hilfe nach dieser Vorschrift \newline nicht fortgesetzt oder beendet werden, \newline prüft der Träger der öffentlichen \newline Jugendhilfe ab einem Jahr vor dem hierfür \newline im \textbf{Hilfe- und Leistungsplan} \newline vorgesehenen Zeitpunkt, ob im Hinblick auf \newline den Bedarf des jungen Menschen ein \newline Zuständigkeitsübergang auf andere \newline Sozialleistungsträger in Betracht kommt; \S{} \newline 36\textbf{d} gilt entsprechend. & (3) Soll eine Hilfe nach dieser Vorschrift \newline nicht fortgesetzt oder beendet werden, \newline prüft der Träger der öffentlichen \newline Jugendhilfe ab einem Jahr vor dem hierfür \newline im \textbf{Hilfe- und Leistungsplan} \newline vorgesehenen Zeitpunkt, ob im Hinblick auf \newline den Bedarf des jungen Menschen ein \newline Zuständigkeitsübergang auf andere \newline Sozialleistungsträger in Betracht kommt; \S{} \newline 36\textbf{d} gilt entsprechend. \\
\hline
\multicolumn{3}{|c|}{\cellcolor{gray!10}\textbf{\S{} 41a -- Nachbetreuung}} \\
\hline
(1) Junge Volljährige werden innerhalb \newline eines angemessenen Zeitraums nach \newline Beendigung der Hilfe bei der \newline Verselbständigung im notwendigen \newline Umfang und in einer für sie verständlichen, \newline nachvollziehbaren und wahrnehmbaren \newline Form beraten und unterstützt. & (1) Junge Volljährige werden innerhalb \newline eines angemessenen Zeitraums nach \newline Beendigung der Hilfe bei der \newline Verselbständigung im notwendigen \newline Umfang und in einer für sie verständlichen, \newline nachvollziehbaren und wahrnehmbaren \newline Form beraten und unterstützt. & (1) Junge Volljährige werden innerhalb \newline eines angemessenen Zeitraums nach \newline Beendigung der Hilfe bei der \newline Verselbständigung im notwendigen \newline Umfang und in einer für sie verständlichen, \newline nachvollziehbaren und wahrnehmbaren \newline Form beraten und unterstützt. \\
\hline
(2) Der angemessene Zeitraum sowie der \newline notwendige Umfang der Beratung und \newline Unterstützung nach Beendigung der Hilfe \newline sollen in dem Hilfeplan nach \S{} 36 Absatz 2 \newline Satz 2, der die Beendigung der Hilfe nach \newline \S{} 41 feststellt, dokumentiert und \newline regelmäßig überprüft werden. Hierzu soll \newline der Träger der öffentlichen Jugendhilfe in \newline regelmäßigen Abständen Kontakt zu dem \newline jungen Volljährigen aufnehmen. & (2) Der angemessene Zeitraum sowie der \newline notwendige Umfang der Beratung und \newline Unterstützung nach Beendigung der Hilfe \newline sollen in dem \textbf{Hilfe- und Leistungsplan \newline nach \S{} 36a}, der die Beendigung der Hilfe \newline nach \S{} 41 feststellt, dokumentiert und \newline regelmäßig überprüft werden. Hierzu soll \newline der Träger der öffentlichen Jugendhilfe in \newline regelmäßigen Abständen Kontakt zu dem \newline jungen Volljährigen aufnehmen. & (2) Der angemessene Zeitraum sowie der \newline notwendige Umfang der Beratung und \newline Unterstützung nach Beendigung der Hilfe \newline sollen in dem \textbf{Hilfe- und Leistungsplan \newline nach \S{} 36a}, der die Beendigung der Hilfe \newline nach \S{} 41 feststellt, dokumentiert und \newline regelmäßig überprüft werden. Hierzu soll \newline der Träger der öffentlichen Jugendhilfe in \newline regelmäßigen Abständen Kontakt zu dem \newline jungen Volljährigen aufnehmen. \\
\hline
\multicolumn{3}{|c|}{\cellcolor{gray!10}\textbf{\S{} 42 -- Absatz 1}} \\
\hline
\textit{--- Fassung 2024 ---} \newline (2) Das Jugendamt hat während der \newline Inobhutnahme unverzüglich das Kind oder \newline den Jugendlichen umfassend und in einer \newline verständlichen, nachvollziehbaren und \newline wahrnehmbaren Form über diese \newline Maßnahme aufzuklären, die Situation, die \newline zur Inobhutnahme geführt hat, zusammen \newline mit dem Kind oder dem Jugendlichen zu \newline klären und Möglichkeiten der Hilfe und \newline Unterstützung aufzuzeigen. Dem Kind oder \newline dem Jugendlichen ist unverzüglich \newline Gelegenheit zu geben, eine Person seines \newline Vertrauens zu benachrichtigen. Das \newline Jugendamt hat während der Inobhutnahme \newline für das Wohl des Kindes oder des \newline Jugendlichen zu sorgen und dabei den \newline notwendigen Unterhalt und die \newline Krankenhilfe sicherzustellen; \S{} 39 Absatz 4 \newline Satz 2 gilt entsprechend. Das Jugendamt \newline ist während der Inobhutnahme berechtigt, \newline alle Rechtshandlungen vorzunehmen, die \newline zum Wohl des Kindes oder Jugendlichen \newline notwendig sind; der mutmaßliche Wille der \newline Personensorge- oder der \newline Erziehungsberechtigten ist dabei \newline angemessen zu berücksichtigen. Im Fall \newline des Absatzes 1 Satz 1 Nummer 3 gehört \newline zu den Rechtshandlungen nach Satz 4, zu \newline denen das Jugendamt verpflichtet ist, \newline insbesondere die unverzügliche Stellung \newline eines Asylantrags für das Kind oder den \newline Jugendlichen in Fällen, in denen Tatsachen \newline die Annahme rechtfertigen, dass das Kind \newline oder der Jugendliche internationalen \newline Schutz im Sinne des \S{} 1 Absatz 1 Nummer \newline 2 des Asylgesetzes benötigt; dabei ist das \newline Kind oder der Jugendliche zu beteiligen. \newline \textit{--- Fassung 2026 ---} \newline (2) Das Jugendamt hat während der \newline Inobhutnahme unverzüglich das Kind oder \newline den Jugendlichen umfassend und in einer \newline verständlichen, nachvollziehbaren und \newline wahrnehmbaren Form über diese \newline Maßnahme aufzuklären, die Situation, die \newline zur Inobhutnahme geführt hat, zusammen \newline mit dem Kind oder dem Jugendlichen zu \newline klären und Möglichkeiten der Hilfe und \newline Unterstützung aufzuzeigen. Dem Kind oder \newline dem Jugendlichen ist unverzüglich \newline Gelegenheit zu geben, eine Person seines \newline Vertrauens zu benachrichtigen. Das \newline Jugendamt hat während der Inobhutnahme & (2) Das Jugendamt hat während der \newline Inobhutnahme unverzüglich das Kind oder \newline den Jugendlichen umfassend und in einer \newline verständlichen, nachvollziehbaren und \newline wahrnehmbaren Form über diese \newline Maßnahme aufzuklären, die Situation, die \newline zur Inobhutnahme geführt hat, zusammen \newline mit dem Kind oder dem Jugendlichen zu \newline klären und Möglichkeiten der Hilfe und \newline Unterstützung aufzuzeigen. Dem Kind oder \newline dem Jugendlichen ist unverzüglich \newline Gelegenheit zu geben, eine Person seines \newline Vertrauens zu benachrichtigen. Das \newline Jugendamt hat während der Inobhutnahme \newline für das Wohl des Kindes oder des \newline Jugendlichen zu sorgen und dabei den \newline notwendigen Unterhalt und die \newline Krankenhilfe sicherzustellen; \textbf{\S{} 39c Absatz \newline 4 Satz 2} gilt entsprechend. Das Jugendamt \newline ist während der Inobhutnahme berechtigt, \newline alle Rechtshandlungen vorzunehmen, die \newline zum Wohl des Kindes oder Jugendlichen \newline notwendig sind; der mutmaßliche Wille der \newline Personensorge- oder der \newline Erziehungsberechtigten ist dabei \newline angemessen zu berücksichtigen. Im Fall \newline des Absatzes 1 Satz 1 Nummer 3 gehört \newline zu den Rechtshandlungen nach Satz 4, zu \newline denen das Jugendamt verpflichtet ist, \newline insbesondere die unverzügliche Stellung \newline eines Asylantrags für das Kind oder den \newline Jugendlichen in Fällen, in denen Tatsachen \newline die Annahme rechtfertigen, dass das Kind \newline oder der Jugendliche internationalen \newline Schutz im Sinne des \S{} 1 Absatz 1 Nummer \newline 2 des Asylgesetzes benötigt; dabei ist das \newline Kind oder der Jugendliche zu beteiligen. & (2) Das Jugendamt hat während der \newline Inobhutnahme unverzüglich das Kind oder \newline den Jugendlichen umfassend und in einer \newline verständlichen, nachvollziehbaren und \newline wahrnehmbaren Form über diese \newline Maßnahme aufzuklären, die Situation, die \newline zur Inobhutnahme geführt hat, zusammen \newline mit dem Kind oder dem Jugendlichen zu \newline klären und Möglichkeiten der Hilfe und \newline Unterstützung aufzuzeigen. Dem Kind oder \newline dem Jugendlichen ist unverzüglich \newline Gelegenheit zu geben, eine Person seines \newline Vertrauens zu benachrichtigen. Das \newline Jugendamt hat während der Inobhutnahme \\
\hline
\textit{--- Fassung 2024 ---} \newline Die Inobhutnahme umfasst die Befugnis, ein \newline Kind oder einen Jugendlichen bei einer \newline geeigneten Person, in einer geeigneten \newline Einrichtung oder in einer sonstigen \newline Wohnform vorläufig unterzubringen; im Fall \newline von Satz 1 Nummer 2 auch ein Kind oder \newline einen Jugendlichen von einer anderen \newline Person wegzunehmen. \newline \textit{--- Fassung 2026 ---} \newline für das Wohl des Kindes oder des \newline Jugendlichen zu sorgen und dabei den \newline notwendigen Unterhalt und die \newline Krankenhilfe sicherzustellen; \S{} 39 Absatz 4 \newline Satz 2 gilt entsprechend. Das Jugendamt \newline ist während der Inobhutnahme berechtigt, \newline alle Rechtshandlungen vorzunehmen, die \newline zum Wohl des Kindes oder Jugendlichen \newline notwendig sind; der mutmaßliche Wille der \newline Personensorge- oder der \newline Erziehungsberechtigten ist dabei \newline angemessen zu berücksichtigen. Im Fall \newline des Absatzes 1 Satz 1 Nummer 3 gehört \newline zu den Rechtshandlungen nach Satz 4, zu \newline denen das Jugendamt verpflichtet ist, \newline insbesondere die unverzügliche Stellung \newline eines Asylantrags für das Kind oder den \newline Jugendlichen in Fällen, in denen Tatsachen \newline die Annahme rechtfertigen, dass das Kind \newline oder der Jugendliche internationalen \newline Schutz im Sinne des \S{} 1 Absatz 1 Nummer \newline 2 des Asylgesetzes benötigt; dabei ist das \newline Kind oder der Jugendliche zu beteiligen. & \textit{unverändert} & für das Wohl des Kindes oder des \newline Jugendlichen zu sorgen und dabei den \newline notwendigen Unterhalt und die \newline Krankenhilfe sicherzustellen; \S{} 39\textbf{c} Absatz \newline 4 Satz 2 gilt entsprechend. Das Jugendamt \newline ist während der Inobhutnahme berechtigt, \newline alle Rechtshandlungen vorzunehmen, die \newline zum Wohl des Kindes oder Jugendlichen \newline notwendig sind; der mutmaßliche Wille der \newline Personensorge- oder der \newline Erziehungsberechtigten ist dabei \newline angemessen zu berücksichtigen. Im Fall \newline des Absatzes 1 Satz 1 Nummer 3 gehört \newline zu den Rechtshandlungen nach Satz 4, zu \newline denen das Jugendamt verpflichtet ist, \newline insbesondere die unverzügliche Stellung \newline eines Asylantrags für das Kind oder den \newline Jugendlichen in Fällen, in denen Tatsachen \newline die Annahme rechtfertigen, dass das Kind \newline oder der Jugendliche internationalen \newline Schutz im Sinne des \S{} 1 Absatz 1 Nummer \newline 2 des Asylgesetzes benötigt; dabei ist das \newline Kind oder der Jugendliche zu beteiligen. \\
\hline
(1) Das Jugendamt ist berechtigt und \newline verpflichtet, ein Kind oder einen \newline Jugendlichen in seine Obhut zu nehmen, & \textit{(1) unverändert} & \textit{unverändert} \\
\hline
1. \newline das Kind oder der Jugendliche um Obhut \newline bittet oder & 1. unverändert & \textit{unverändert} \\
\hline
2. \newline eine dringende Gefahr für das Wohl des \newline Kindes oder des Jugendlichen die \newline Inobhutnahme erfordert und \newline a) \newline die Personensorgeberechtigten nicht \newline widersprechen oder \newline b) \newline eine familiengerichtliche Entscheidung \newline nicht rechtzeitig eingeholt werden kann \newline oder & 2. unverändert & \textit{unverändert} \\
\hline
3. \newline ein ausländisches Kind oder ein \newline ausländischer Jugendlicher unbegleitet \newline nach Deutschland kommt und sich weder \newline Personensorge- noch \newline Erziehungsberechtigte im Inland aufhalten. & 3. unverändert & \textit{unverändert} \\
\hline
(3) Das Jugendamt hat im Fall des Absatzes \newline 1 Satz 1 Nummer 1 und 2 die \newline Personensorge- oder \newline Erziehungsberechtigten unverzüglich von \newline der Inobhutnahme zu unterrichten, sie in \newline einer verständlichen, nachvollziehbaren und \newline wahrnehmbaren Form umfassend über \newline diese Maßnahme aufzuklären und mit ihnen \newline das Gefährdungsrisiko abzuschätzen. \newline Widersprechen die Personensorge- oder \newline Erziehungsberechtigten der Inobhutnahme, \newline so hat das Jugendamt unverzüglich & \textit{(3) unverändert} & \textit{unverändert} \\
\hline
1. \newline das Kind oder den Jugendlichen den \newline Personensorge- oder \newline Erziehungsberechtigten zu übergeben, \newline sofern nach der Einschätzung des \newline Jugendamts eine Gefährdung des \newline Kindeswohls nicht besteht oder die \newline Personensorge- oder \newline Erziehungsberechtigten bereit und in der \newline Lage sind, die Gefährdung abzuwenden \newline oder & 1. unverändert & \textit{unverändert} \\
\hline
2. \newline eine Entscheidung des Familiengerichts \newline über die erforderlichen Maßnahmen zum \newline Wohl des Kindes oder des Jugendlichen \newline herbeizuführen. \newline Sind die Personensorge- oder \newline Erziehungsberechtigten nicht erreichbar, so \newline gilt Satz 2 Nummer 2 entsprechend. Im Fall \newline des Absatzes 1 Satz 1 Nummer 3 ist \newline unverzüglich die Bestellung eines \newline Vormunds oder Pflegers zu veranlassen. \newline Widersprechen die \newline Personensorgeberechtigten der \newline Inobhutnahme nicht, so ist unverzüglich ein \newline Hilfeplanverfahren zur Gewährung einer \newline Hilfe einzuleiten. & 2. unverändert & \textit{unverändert} \\
\hline
(4) Die Inobhutnahme endet mit & \textit{(4) unverändert} & \textit{unverändert} \\
\hline
1. \newline der Übergabe des Kindes oder \newline Jugendlichen an die Personensorge- oder \newline Erziehungsberechtigten, & 1. unverändert & \textit{unverändert} \\
\hline
2. \newline der Entscheidung über die Gewährung von \newline Hilfen nach dem Sozialgesetzbuch. & 2. unverändert & \textit{unverändert} \\
\hline
(5) Freiheitsentziehende Maßnahmen im \newline Rahmen der Inobhutnahme sind nur \newline zulässig, wenn und soweit sie erforderlich \newline sind, um eine Gefahr für Leib oder Leben \newline des Kindes oder des Jugendlichen oder \newline eine Gefahr für Leib oder Leben Dritter \newline abzuwenden. Die Freiheitsentziehung ist \newline ohne gerichtliche Entscheidung spätestens \newline mit Ablauf des Tages nach ihrem Beginn \newline zu beenden. & \textit{(5) unverändert} & \textit{unverändert} \\
\hline
(6) Ist bei der Inobhutnahme die \newline Anwendung unmittelbaren Zwangs \newline erforderlich, so sind die dazu befugten \newline Stellen hinzuzuziehen. & \textit{(6) unverändert} & \textit{unverändert} \\
\hline
\multicolumn{3}{|c|}{\cellcolor{gray!10}\textbf{\S{} 42a -- Vorläufige Inobhutnahme von ausländischen Kindern und Jugendlichen nach unbegleiteter Einreise}} \\
\hline
Absätze 5 bis 6 & \textit{unverändert} & \textit{unverändert} \\
\hline
(4) Das Jugendamt hat der nach \newline Landesrecht für die Verteilung von \newline unbegleiteten ausländischen Kindern und \newline Jugendlichen zuständigen Stelle die \newline vorläufige Inobhutnahme des Kindes oder \newline des Jugendlichen innerhalb von sieben \newline Werktagen nach Beginn der Maßnahme \newline zur Erfüllung der in \S{} 42b genannten \newline Aufgaben mitzuteilen. Zu diesem Zweck \newline sind auch die Ergebnisse der Einschätzung \newline nach Absatz 2 Satz 1 mitzuteilen. Die nach \newline Landesrecht zuständige Stelle hat \newline gegenüber dem Bundesverwaltungsamt \newline innerhalb von drei Werktagen das Kind \newline oder den Jugendlichen zur Verteilung \newline anzumelden oder den Ausschluss der \newline Verteilung anzuzeigen. & \textit{unverändert} & (4) Das Jugendamt hat der nach \newline Landesrecht für die Verteilung von \newline unbegleiteten ausländischen Kindern und \newline Jugendlichen zuständigen Stelle die \newline vorläufige Inobhutnahme des Kindes oder \newline des Jugendlichen innerhalb von einem \newline Monat nach Beginn der Maßnahme zur \newline Erfüllung der in \S{} 42b genannten Aufgaben \newline mitzuteilen. Zu diesem Zweck sind auch \newline die Ergebnisse der Einschätzung nach \newline Absatz 2 Satz 1 mitzuteilen. Die nach \newline Landesrecht zuständige Stelle hat \newline gegenüber dem Bundesverwaltungsamt \newline innerhalb von drei Werktagen das Kind \newline oder den Jugendlichen zur Verteilung \newline anzumelden oder den Ausschluss der \newline Verteilung anzuzeigen. \\
\hline
\multicolumn{3}{|c|}{\cellcolor{gray!10}\textbf{\S{} 42b -- Verfahren zur Verteilung unbegleiteter ausländischer Kinder und Jugendlicher}} \\
\hline
Absätze 5 bis 8 & \textit{unverändert} & \textit{unverändert} \\
\hline
(4) Die Durchführung eines \newline Verteilungsverfahrens ist bei einem \newline unbegleiteten ausländischen Kind oder \newline Jugendlichen ausgeschlossen, wenn \newline 1. dadurch dessen Wohl gefährdet würde, \newline 2. dessen Gesundheitszustand die \newline Durchführung eines Verteilungsverfahrens \newline innerhalb von 14 Werktagen nach Beginn \newline der vorläufigen Inobhutnahme gemäß \S{} \newline 42a nicht zulässt, \newline 3. dessen Zusammenführung mit einer \newline verwandten Person kurzfristig erfolgen \newline kann, zum Beispiel aufgrund der \newline Verordnung (EU) Nr. 604/2013 des \newline Europäischen Parlaments und des Rates \newline vom 26. Juni 2013 zur Festlegung der \newline Kriterien und Verfahren zur Bestimmung \newline des Mitgliedstaats, der für die Prüfung \newline eines von einem Drittstaatsangehörigen \newline oder Staatenlosen in einem Mitgliedstaat \newline gestellten Antrags auf internationalen \newline Schutz zuständig ist (ABl. L 180 vom \newline 29.6.2013, S. 31), und dies dem Wohl des \newline Kindes entspricht oder \newline 4. die Durchführung des \newline Verteilungsverfahrens nicht innerhalb von \newline einem Monat nach Beginn der vorläufigen \newline Inobhutnahme erfolgt. & \textit{unverändert} & (4) Die Durchführung eines \newline Verteilungsverfahrens ist bei einem \newline unbegleiteten ausländischen Kind oder \newline Jugendlichen ausgeschlossen, wenn \newline 1. dadurch dessen Wohl gefährdet würde, \newline 2. dessen Gesundheitszustand die \newline Durchführung eines Verteilungsverfahrens \newline innerhalb von 14 Werktagen nach Beginn \newline der vorläufigen Inobhutnahme gemäß \S{} \newline 42a nicht zulässt, \newline 3. dessen Zusammenführung mit einer \newline verwandten Person kurzfristig erfolgen \newline kann, zum Beispiel aufgrund der \newline Verordnung (EU) Nr. 604/2013 des \newline Europäischen Parlaments und des Rates \newline vom 26. Juni 2013 zur Festlegung der \newline Kriterien und Verfahren zur Bestimmung \newline des Mitgliedstaats, der für die Prüfung \newline eines von einem Drittstaatsangehörigen \newline oder Staatenlosen in einem Mitgliedstaat \newline gestellten Antrags auf internationalen \newline Schutz zuständig ist (ABl. L 180 vom \newline 29.6.2013, S. 31), und dies dem Wohl des \newline Kindes entspricht oder \newline 4. die Durchführung des \newline Verteilungsverfahrens nicht innerhalb von \newline zwei Monaten nach Beginn der vorläufigen \newline Inobhutnahme erfolgt. \\
\hline
\multicolumn{3}{|c|}{\cellcolor{gray!10}\textbf{\S{} 42e -- Berichtspflicht}} \\
\hline
Die Bundesregierung hat dem Deutschen \newline Bundestag jährlich einen Bericht über die \newline Situation unbegleiteter ausländischer \newline Minderjähriger in Deutschland vorzulegen. & \textit{unverändert} & Unbegleitete ausländische Jugendliche \newline sind verpflichtet, ihren gewöhnlichen \newline Aufenthalt zu nehmen in dem Bereich des \newline aufgrund der Zuweisungsentscheidung der \newline zuständigen Landesbehörde nach \S{} 88a \newline Absatz 2 Satz 1 zuständigen Jugendamts. \newline Über Folgen eines Verstoßes sind \newline unbegleitete ausländische Jugendliche \newline aufzuklären. \\
\hline
\multicolumn{3}{|c|}{\cellcolor{gray!10}\textbf{\S{} 42f -- Behördliches Verfahren zur Altersfeststellung}} \\
\hline
(1) Das Jugendamt hat im Rahmen der \newline vorläufigen Inobhutnahme der \newline ausländischen Person gemäß \S{} 42a deren \newline Minderjährigkeit durch Einsichtnahme in \newline deren Ausweispapiere festzustellen oder \newline hilfsweise mittels einer qualifizierten \newline Inaugenscheinnahme einzuschätzen und \newline festzustellen. \S{} 8 Absatz 1 und \S{} 42 Absatz \newline 2 Satz 2 sind entsprechend anzuwenden. & \textit{unverändert} & (1) Das Jugendamt hat im Rahmen der \newline vorläufigen Inobhutnahme der \newline ausländischen Person gemäß \S{} 42a deren \newline Minderjährigkeit durch Einsichtnahme in \newline deren Ausweispapiere festzustellen oder \newline hilfsweise mittels einer qualifizierten \newline Inaugenscheinnahme einzuschätzen und \newline festzustellen. \textbf{Die betroffene Person ist \newline über die Folgen der Altersbestimmung \newline und einer Verweigerung einer \newline Mitwirkung aufzuklären.} \S{} 8 Absatz 1 und \newline \S{} 42 Absatz 2 Satz 2 sind entsprechend \newline anzuwenden. \\
\hline
(2) Auf Antrag des Betroffenen oder seines \newline Vertreters oder von Amts wegen hat das \newline Jugendamt in Zweifelsfällen eine ärztliche \newline Untersuchung zur Altersbestimmung zu \newline veranlassen. Ist eine ärztliche \newline Untersuchung durchzuführen, ist die \newline betroffene Person durch das Jugendamt \newline umfassend über die \newline Untersuchungsmethode und über die \newline möglichen Folgen der Altersbestimmung \newline aufzuklären. Ist die ärztliche Untersuchung \newline von Amts wegen durchzuführen, ist die \newline betroffene Person zusätzlich über die \newline Folgen einer Weigerung, sich der ärztlichen \newline Untersuchung zu unterziehen, aufzuklären; \newline die Untersuchung darf nur mit Einwilligung \newline der betroffenen Person und ihres Vertreters \newline durchgeführt werden. Die \S{}\S{} 60, 62 und 65 \newline bis 67 des Ersten Buches sind \newline entsprechend anzuwenden. & \textit{unverändert} & (2) \textbf{Auf Antrag der betroffenen Person \newline oder ihres Vertreters oder von Amts \newline wegen hat das Jugendamt eine ärztliche \newline Untersuchung zur Altersbestimmung zu \newline veranlassen, wenn nicht mit an \newline Sicherheit grenzender \newline Wahrscheinlichkeit nach Absatz 1 die \newline Volljährigkeit der ausländischen Person \newline festgestellt wurde. Die betroffene \newline Person ist durch das Jugendamt \newline umfassend über die \newline Untersuchungsmethode aufzuklären. \newline Die Untersuchung darf nur mit \newline Einwilligung der betroffenen Person und \newline ihres Vertreters durchgeführt werden.} \newline Die \S{}\S{} 60, 62 und 65 bis 67 des Ersten \newline Buches sind entsprechend anzuwenden. \\
\hline
(3) Widerspruch und Klage gegen die \newline Entscheidung des Jugendamts, aufgrund \newline der Altersfeststellung nach dieser Vorschrift \newline die vorläufige Inobhutnahme nach \S{} 42a \newline oder die Inobhutnahme nach \S{} 42 Absatz 1 \newline Satz 1 Nummer 3 abzulehnen oder zu \newline beenden, haben keine aufschiebende & \textit{unverändert} & \textbf{(3) Eine Klage gegen die Entscheidung \newline des Jugendamts, aufgrund der \newline Altersfeststellung nach dieser Vorschrift \newline die vorläufige Inobhutnahme nach \S{} 42a \newline oder die Inobhutnahme nach \S{} 42 \newline Absatz 1 Satz 1 Nummer 3 abzulehnen \newline oder zu beenden, hat keine } \\
\hline
Wirkung. Landesrecht kann bestimmen, \newline dass gegen diese Entscheidung Klage \newline ohne Nachprüfung in einem Vorverfahren \newline nach \S{} 68 der Verwaltungsgerichtsordnung \newline erhoben werden kann. & \textit{unverändert} & \textbf{aufschiebende Wirkung. Das \newline Vorverfahren nach \S{} 68 der \newline Verwaltungsgerichtsordnung entfällt. } \\
\hline
\multicolumn{3}{|c|}{\cellcolor{gray!10}\textbf{\S{} 44 -- Erlaubnis zur Vollzeitpflege}} \\
\hline
(1) Wer ein Kind oder einen Jugendlichen \newline über Tag und Nacht in seinem Haushalt \newline aufnehmen will (Pflegeperson), bedarf der \newline Erlaubnis. Einer Erlaubnis bedarf nicht, wer \newline ein Kind oder einen Jugendlichen & (1) Wer ein Kind oder einen Jugendlichen \newline über Tag und Nacht in seinem Haushalt \newline aufnehmen will (Pflegeperson), bedarf der \newline Erlaubnis. Einer Erlaubnis bedarf nicht, wer \newline ein Kind oder einen Jugendlichen & (1) Wer ein Kind oder einen Jugendlichen \newline über Tag und Nacht in seinem Haushalt \newline aufnehmen will (Pflegeperson), bedarf der \newline Erlaubnis. Einer Erlaubnis bedarf nicht, wer \newline ein Kind oder einen Jugendlichen \\
\hline
1. \newline im Rahmen von Hilfe zur Erziehung oder \newline von Eingliederungshilfe für seelisch \newline behinderte Kinder und Jugendliche auf \newline Grund einer Vermittlung durch das \newline Jugendamt, & 1. \newline im Rahmen von Hilfe zur Erziehung oder \newline von \textbf{Leistungen der Eingliederungshilfe \newline für Kinder und Jugendliche mit \newline Behinderungen} auf Grund einer \newline Vermittlung durch das Jugendamt, & 1. \newline im Rahmen von Hilfe zur Erziehung oder \newline von \textbf{Leistungen der Eingliederungshilfe \newline für Kinder und Jugendliche mit \newline Behinderungen} auf Grund einer \newline Vermittlung durch das Jugendamt, \\
\hline
2. \newline als Vormund oder Pfleger im Rahmen \newline seines Wirkungskreises, & 2. unverändert & 2. unverändert \\
\hline
3. \newline als Verwandter oder Verschwägerter bis \newline zum dritten Grad, & 3. unverändert & 3. unverändert \\
\hline
4. \newline bis zur Dauer von acht Wochen, & 4. unverändert & 4. unverändert \\
\hline
5. \newline im Rahmen eines Schüler- oder \newline Jugendaustausches, & 5. unverändert & 5. unverändert \\
\hline
6. \newline in Adoptionspflege (\S{} 1744 des \newline Bürgerlichen Gesetzbuchs) \newline über Tag und Nacht aufnimmt. & 6. unverändert & 6. unverändert \\
\hline
(2) Die Erlaubnis ist zu versagen, wenn das \newline Wohl des Kindes oder des Jugendlichen in \newline der Pflegestelle nicht gewährleistet ist. \S{} \newline 72a Absatz 1 und 5 gilt entsprechend. & \textit{(2) unverändert} & \textit{(2) unverändert} \\
\hline
(3) Das Jugendamt soll den Erfordernissen \newline des Einzelfalls entsprechend an Ort und \newline Stelle überprüfen, ob die Voraussetzungen \newline für die Erteilung der Erlaubnis weiter \newline bestehen. Ist das Wohl des Kindes oder \newline des Jugendlichen in der Pflegestelle \newline gefährdet und ist die Pflegeperson nicht \newline bereit oder in der Lage, die Gefährdung \newline abzuwenden, so ist die Erlaubnis \newline zurückzunehmen oder zu widerrufen. & \textit{(3) unverändert} & \textit{(3) unverändert} \\
\hline
(4) Wer ein Kind oder einen Jugendlichen \newline in erlaubnispflichtige Familienpflege \newline aufgenommen hat, hat das Jugendamt \newline über wichtige Ereignisse zu unterrichten, \newline die das Wohl des Kindes oder des \newline Jugendlichen betreffen. & \textit{(4) unverändert} & \textit{(4) unverändert} \\
\hline
\multicolumn{3}{|c|}{\cellcolor{gray!10}\textbf{\S{} 45 -- Erlaubnis für den Betrieb einer Einrichtung}} \\
\hline
\textit{--- Fassung 2024 ---} \newline (1) Der Träger einer Einrichtung, nach \S{} \newline 45a bedarf für den Betrieb der Einrichtung \newline der Erlaubnis. Einer Erlaubnis bedarf nicht, \newline wer \newline \textit{--- Fassung 2026 ---} \newline (1) Der Träger einer Einrichtung, nach \S{} \newline 45a bedarf für den Betrieb der Einrichtung \newline der Erlaubnis. Einer Erlaubnis bedarf nicht, \newline wer \newline Gaststättengewerbes der Aufnahme von \newline Kindern oder Jugendlichen dient. & \textit{(1) unverändert} & \textit{(1) unverändert} \\
\hline
1. \newline eine Jugendfreizeiteinrichtung, eine \newline Jugendbildungseinrichtung, eine \newline Jugendherberge oder ein Schullandheim \newline betreibt, & 1. unverändert & 1. unverändert \\
\hline
2. \newline ein Schülerheim betreibt, das \newline landesgesetzlich der Schulaufsicht \newline untersteht, & 2. unverändert & 2. unverändert \\
\hline
\textit{--- Fassung 2024 ---} \newline 3. \newline eine Einrichtung betreibt, die außerhalb der \newline Jugendhilfe liegende Aufgaben für Kinder \newline oder Jugendliche wahrnimmt, wenn für sie \newline eine entsprechende gesetzliche Aufsicht \newline besteht oder im Rahmen des Hotel- und \newline Gaststättengewerbes der Aufnahme von \newline Kindern oder Jugendlichen dient. \newline \textit{--- Fassung 2026 ---} \newline 3. \newline eine Einrichtung betreibt, die außerhalb der \newline Jugendhilfe liegende Aufgaben für Kinder \newline oder Jugendliche wahrnimmt, wenn für sie \newline eine entsprechende gesetzliche Aufsicht \newline besteht oder im Rahmen des Hotel- und & 3. unverändert & 3. unverändert \\
\hline
(2) Die Erlaubnis ist zu erteilen, wenn das \newline Wohl der Kinder und Jugendlichen in der \newline Einrichtung gewährleistet ist. Dies ist in der \newline Regel anzunehmen, wenn & \textit{(2) unverändert} & \textit{(2) unverändert} \\
\hline
1. \newline der Träger die für den Betrieb der \newline Einrichtung erforderliche Zuverlässigkeit \newline besitzt, & 1. unverändert & 1. \newline die Träger die für den Betrieb der \newline Einrichtung erforderliche Zuverlässigkeit \newline besitzt \textbf{und die Gewähr für eine den \newline Zielen des Grundgesetzes förderliche \newline Arbeit bietet,} \\
\hline
2. \newline die dem Zweck und der Konzeption der \newline Einrichtung entsprechenden räumlichen, \newline fachlichen, wirtschaftlichen und personellen \newline Voraussetzungen für den Betrieb erfüllt \newline sind und durch den Träger gewährleistet \newline werden, & 2. unverändert & 2. \newline die dem Zweck und der Konzeption der \newline Einrichtung entsprechenden räumlichen, \newline fachlichen, wirtschaftlichen und personellen \newline Voraussetzungen für den Betrieb erfüllt \newline sind und durch den Träger gewährleistet \newline werden, wobei sich die personellen \newline Voraussetzungen auch nach den jeweils in \newline der Einrichtung wahrzunehmenden \newline Funktionen richten, \\
\hline
3. \newline die gesellschaftliche und sprachliche \newline Integration und ein gesundheitsförderliches \newline Lebensumfeld in der Einrichtung unterstützt \newline werden sowie die gesundheitliche Vorsorge \newline und die medizinische Betreuung der Kinder \newline und Jugendlichen nicht erschwert werden \newline sowie & 3. unverändert & 3. unverändert \\
\hline
4. \newline zur Sicherung der Rechte und des Wohls \newline von Kindern und Jugendlichen in der \newline Einrichtung die Entwicklung, Anwendung \newline und Überprüfung eines Konzepts zum \newline Schutz vor Gewalt, geeignete Verfahren \newline der Selbstvertretung und Beteiligung sowie \newline der Möglichkeit der Beschwerde in \newline persönlichen Angelegenheiten innerhalb \newline und außerhalb der Einrichtung \newline gewährleistet werden. & 4. unverändert & 4. unverändert \\
\hline
Die nach Satz 2 Nummer 1 erforderliche \newline Zuverlässigkeit besitzt ein Träger \newline insbesondere dann nicht, wenn er & \textit{unverändert} & \textit{unverändert} \\
\hline
1. \newline in der Vergangenheit nachhaltig gegen \newline seine Mitwirkungs- und Meldepflichten \newline nach den \S{}\S{} 46 und 47 verstoßen hat, & 1. unverändert & 1. unverändert \\
\hline
2. \newline Personen entgegen eines behördlichen \newline Beschäftigungsverbotes nach \S{} 48 \newline beschäftigt oder & 2. unverändert & 2. unverändert \\
\hline
3. \newline wiederholt gegen behördliche Auflagen \newline verstoßen hat. & 3. unverändert & 3. unverändert \\
\hline
(3) Zur Prüfung der Voraussetzungen hat \newline der Träger der Einrichtung mit dem Antrag & \textit{(3) unverändert} & \textit{(3) unverändert} \\
\hline
1. \newline die Konzeption der Einrichtung vorzulegen, \newline die auch Auskunft über Maßnahmen zur \newline Qualitätsentwicklung und -sicherung sowie \newline zur ordnungsgemäßen Buch- und \newline Aktenführung in Bezug auf den Betrieb der \newline Einrichtung gibt, sowie & 1. unverändert & 1. unverändert \\
\hline
2. \newline im Hinblick auf die Eignung des Personals \newline nachzuweisen, dass die Vorlage und \newline Prüfung von aufgabenspezifischen \newline Ausbildungsnachweisen sowie von \newline Führungszeugnissen nach \S{} 30 Absatz 5 \newline und \S{} 30a Absatz 1 des \newline Bundeszentralregistergesetzes \newline sichergestellt sind; Führungszeugnisse \newline sind von dem Träger der Einrichtung in \newline regelmäßigen Abständen erneut \newline anzufordern und zu prüfen. & 2. unverändert & 2. unverändert \\
\hline
(4) Die Erlaubnis kann mit \newline Nebenbestimmungen versehen werden. \newline Zur Gewährleistung des Wohls der Kinder \newline und der Jugendlichen können \newline nachträgliche Auflagen erteilt werden. & \textit{(4) unverändert} & \textit{(4) unverändert} \\
\hline
(5) Besteht für eine erlaubnispflichtige \newline Einrichtung eine Aufsicht nach anderen \newline Rechtsvorschriften, so hat die zuständige \newline Behörde ihr Tätigwerden zuvor mit der \newline anderen Behörde abzustimmen. Sie hat \newline den Träger der Einrichtung rechtzeitig auf \newline weitergehende Anforderungen nach \newline anderen Rechtsvorschriften hinzuweisen. & \textit{(5) unverändert} & \textit{(5) unverändert} \\
\hline
(6) Sind in einer Einrichtung Mängel \newline festgestellt worden, so soll die zuständige \newline Behörde zunächst den Träger der \newline Einrichtung über die Möglichkeiten zur \newline Beseitigung der Mängel beraten. Wenn \newline sich die Beseitigung der Mängel auf \newline Entgelte oder Vergütungen nach \S{} 134 des \newline Neunten Buches oder nach \S{} 76 des \newline Zwölften Buches auswirken kann, so ist der \newline Träger der Eingliederungshilfe oder der \newline Sozialhilfe, mit dem Vereinbarungen nach \newline diesen Vorschriften bestehen, an der \newline Beratung zu beteiligen. Werden \newline festgestellte Mängel nicht behoben, so \newline können dem Träger der Einrichtung \newline Auflagen nach Absatz 4 Satz 2 erteilt \newline werden. Wenn sich eine Auflage auf \newline Entgelte oder Vergütungen nach \S{} 134 des \newline Neunten Buches oder nach \S{} 76 des \newline Zwölften Buches auswirkt, so entscheidet \newline die zuständige Behörde nach Anhörung \newline des Trägers der Eingliederungshilfe oder \newline der Sozialhilfe, mit dem Vereinbarungen \newline nach diesen Vorschriften bestehen, über \newline die Erteilung der Auflage. Die Auflage ist \newline nach Möglichkeit in Übereinstimmung mit \newline den nach \S{} 134 des Neunten Buches oder \newline nach den \S{}\S{} 75 bis 80 des Zwölften \newline Buches getroffenen Vereinbarungen \newline auszugestalten. & (6) Sind in einer Einrichtung Mängel \newline festgestellt worden, so soll die zuständige \newline Behörde zunächst den Träger der \newline Einrichtung über die Möglichkeiten zur \newline Beseitigung der Mängel beraten. Wenn \newline sich die Beseitigung der Mängel auf \newline Entgelte oder Vergütungen \textbf{nach \S{} 134 \newline des Neunten Buches oder }nach \S{} 76 des \newline Zwölften Buches auswirken kann, so ist der \newline Träger der Eingliederungshilfe oder der \newline Sozialhilfe, mit dem Vereinbarungen nach \newline diesen Vorschriften bestehen, an der \newline Beratung zu beteiligen. Werden \newline festgestellte Mängel nicht behoben, so \newline können dem Träger der Einrichtung \newline Auflagen nach Absatz 4 Satz 2 erteilt \newline werden. Wenn sich eine Auflage auf \newline Entgelte oder Vergütungen \textbf{nach \S{} 134 \newline des Neunten Buches oder }nach \S{} 76 des \newline Zwölften Buches auswirkt, so entscheidet \newline die zuständige Behörde nach Anhörung \newline des Trägers der Eingliederungshilfe oder \newline der Sozialhilfe, mit dem Vereinbarungen \newline nach diesen Vorschriften bestehen, über \newline die Erteilung der Auflage. Die Auflage ist \newline nach Möglichkeit in Übereinstimmung mit \newline den \textbf{nach \S{} 134 des Neunten Buches \newline oder} nach den \S{}\S{} 75 bis 80 des Zwölften \newline Buches getroffenen Vereinbarungen \newline auszugestalten. & (6) Sind in einer Einrichtung Mängel \newline festgestellt worden, so soll die zuständige \newline Behörde zunächst den Träger der \newline Einrichtung über die Möglichkeiten zur \newline Beseitigung der Mängel beraten. Wenn \newline sich die Beseitigung der Mängel auf \newline Entgelte oder Vergütungen \textbf{nach \S{} 134 \newline des Neunten Buches oder }nach \S{} 76 des \newline Zwölften Buches auswirken kann, so ist der \newline Träger der Eingliederungshilfe oder der \newline Sozialhilfe, mit dem Vereinbarungen nach \newline diesen Vorschriften bestehen, an der \newline Beratung zu beteiligen. Werden \newline festgestellte Mängel nicht behoben, so \newline können dem Träger der Einrichtung \newline Auflagen nach Absatz 4 Satz 2 erteilt \newline werden. Wenn sich eine Auflage auf \newline Entgelte oder Vergütungen \textbf{nach \S{} 134 \newline des Neunten Buches oder }nach \S{} 76 des \newline Zwölften Buches auswirkt, so entscheidet \newline die zuständige Behörde nach Anhörung \newline des Trägers der Eingliederungshilfe oder \newline der Sozialhilfe, mit dem Vereinbarungen \newline nach diesen Vorschriften bestehen, über \newline die Erteilung der Auflage. Die Auflage ist \newline nach Möglichkeit in Übereinstimmung mit \newline den \textbf{nach \S{} 134 des Neunten Buches \newline oder} nach den \S{}\S{} 75 bis 80 des Zwölften \newline Buches getroffenen Vereinbarungen \newline auszugestalten. \\
\hline
\multicolumn{3}{|c|}{\cellcolor{gray!10}\textbf{\S{} 50 -- Mitwirkung in Verfahren vor den Familiengerichten}} \\
\hline
(1) Das Jugendamt unterstützt das \newline Familiengericht bei allen Maßnahmen, die \newline die Sorge für die Person von Kindern und \newline Jugendlichen betreffen. Es hat in folgenden \newline Verfahren nach dem Gesetz über das \newline Verfahren in Familiensachen und in den \newline Angelegenheiten der freiwilligen \newline Gerichtsbarkeit mitzuwirken: & \textit{(1) unverändert} & \textit{(1) unverändert} \\
\hline
1. \newline Kindschaftssachen (\S{} 162 des Gesetzes \newline über das Verfahren in Familiensachen und \newline in den Angelegenheiten der freiwilligen \newline Gerichtsbarkeit), & 1. unverändert & 1. unverändert \\
\hline
2. \newline Abstammungssachen (\S{} 176 des Gesetzes \newline über das Verfahren in Familiensachen und \newline in den Angelegenheiten der freiwilligen \newline Gerichtsbarkeit), & 2. unverändert & 2. unverändert \\
\hline
3. \newline Adoptionssachen (\S{} 188 Absatz 2, \S{}\S{} 189, \newline 194, 195 des Gesetzes über das Verfahren \newline in Familiensachen und in den \newline Angelegenheiten der freiwilligen \newline Gerichtsbarkeit), & 3. unverändert & 3. unverändert \\
\hline
4. \newline Ehewohnungssachen (\S{} 204 Absatz 2, \S{} \newline 205 des Gesetzes über das Verfahren in \newline Familiensachen und in den \newline Angelegenheiten der freiwilligen \newline Gerichtsbarkeit) und & 4. unverändert & 4. unverändert \\
\hline
5. \newline Gewaltschutzsachen (\S{}\S{} 212, 213 des \newline Gesetzes über das Verfahren in \newline Familiensachen und in den \newline Angelegenheiten der freiwilligen \newline Gerichtsbarkeit). & 5. unverändert & 5. unverändert \\
\hline
\textit{--- Fassung 2024 ---} \newline (2) Das Jugendamt unterrichtet \newline insbesondere über angebotene und \newline erbrachte Leistungen, bringt erzieherische \newline und soziale Gesichtspunkte zur \newline Entwicklung des Kindes oder des \newline Jugendlichen ein und weist auf weitere \newline Möglichkeiten der Hilfe hin. In Verfahren \newline nach den \S{}\S{} 1631b, 1632 Absatz 4, den \S{}\S{} \newline 1666, 1666a und 1682 des Bürgerlichen \newline Gesetzbuchs sowie in Verfahren, die die \newline Abänderung, Verlängerung oder \newline Aufhebung von nach diesen Vorschriften \newline getroffenen Maßnahmen betreffen, legt das \newline Jugendamt dem Familiengericht den \newline Hilfeplan nach \S{} 36 Absatz 2 Satz 2 vor. \newline Dieses Dokument beinhaltet ausschließlich \newline das Ergebnis der Bedarfsfeststellung, die \newline vereinbarte Art der Hilfegewährung \newline einschließlich der hiervon umfassten \newline Leistungen sowie das Ergebnis etwaiger \newline Überprüfungen dieser Feststellungen. In \newline anderen die Person des Kindes \newline betreffenden Kindschaftssachen legt das \newline Jugendamt den Hilfeplan auf Anforderung \newline des Familiengerichts vor. Das Jugendamt \newline informiert das Familiengericht in dem \newline Termin nach \S{} 155 Absatz 2 des Gesetzes \newline über das Verfahren in Familiensachen und \newline in den Angelegenheiten der freiwilligen \newline Gerichtsbarkeit über den Stand des \newline Beratungsprozesses. \S{} 64 Absatz 2 und \S{} \newline 65 Absatz 1 Satz 1 Nummer 1 und 2 \newline bleiben unberührt. \newline \textit{--- Fassung 2026 ---} \newline (2) Das Jugendamt unterrichtet \newline insbesondere über angebotene und \newline erbrachte Leistungen, bringt erzieherische \newline und soziale Gesichtspunkte zur \newline Entwicklung des Kindes oder des \newline Jugendlichen ein und weist auf weitere \newline Möglichkeiten der Hilfe hin. In Verfahren \newline nach den \S{}\S{} 1631b, 1632 Absatz 4, den \S{}\S{} \newline 1666, 1666a und 1682 des Bürgerlichen \newline Gesetzbuchs sowie in Verfahren, die die \newline Abänderung, Verlängerung oder \newline Aufhebung von nach diesen Vorschriften \newline getroffenen Maßnahmen betreffen, legt das & (2) Das Jugendamt unterrichtet \newline insbesondere über angebotene und \newline erbrachte Leistungen, bringt erzieherische \newline und soziale Gesichtspunkte zur \newline Entwicklung des Kindes oder des \newline Jugendlichen ein und weist auf weitere \newline Möglichkeiten der Hilfe hin. In Verfahren \newline nach den \S{}\S{} 1631b, 1632 Absatz 4, den \S{}\S{} \newline 1666, 1666a und 1682 des Bürgerlichen \newline Gesetzbuchs sowie in Verfahren, die die \newline Abänderung, Verlängerung oder \newline Aufhebung von nach diesen Vorschriften \newline getroffenen Maßnahmen betreffen, legt das \newline Jugendamt dem Familiengericht den \textbf{Hilfe- \newline und Leistungsplan nach \S{} 36a Absatz 1 \newline Satz 2} vor. Dieses Dokument beinhaltet \newline ausschließlich das Ergebnis der \newline Bedarfsfeststellung, die vereinbarte Art der \newline Hilfegewährung einschließlich der hiervon \newline umfassten Leistungen sowie das Ergebnis \newline etwaiger Überprüfungen dieser \newline Feststellungen. In anderen die Person des \newline Kindes betreffenden Kindschaftssachen \newline legt das Jugendamt den Hilfeplan auf \newline Anforderung des Familiengerichts vor. Das \newline Jugendamt informiert das Familiengericht \newline in dem Termin nach \S{} 155 Absatz 2 des \newline Gesetzes über das Verfahren in \newline Familiensachen und in den \newline Angelegenheiten der freiwilligen \newline Gerichtsbarkeit über den Stand des \newline Beratungsprozesses. \S{} 64 Absatz 2 und \S{} \newline 65 Absatz 1 Satz 1 Nummer 1 und 2 \newline bleiben unberührt. & (2) Das Jugendamt unterrichtet \newline insbesondere über angebotene und \newline erbrachte Leistungen, bringt erzieherische \newline und soziale Gesichtspunkte zur \newline Entwicklung des Kindes oder des \newline Jugendlichen ein und weist auf weitere \newline Möglichkeiten der Hilfe hin. In Verfahren \newline nach den \S{}\S{} 1631b, 1632 Absatz 4, den \S{}\S{} \newline 1666, 1666a und 1682 des Bürgerlichen \newline Gesetzbuchs sowie in Verfahren, die die \newline Abänderung, Verlängerung oder \newline Aufhebung von nach diesen Vorschriften \newline getroffenen Maßnahmen betreffen, legt das \\
\hline
dem nach \S{} 87c Absatz 6 Satz 2 \newline zuständigen Jugendamt zu den in \S{} 58 \newline genannten Zwecken unverzüglich mit. \newline Mitzuteilen sind auch das Geburtsdatum \newline und der Geburtsort des Kindes oder des \newline Jugendlichen sowie der Name, den das \newline Kind oder der Jugendliche zur Zeit der \newline Beurkundung seiner Geburt geführt hat. & \textit{unverändert} & \textit{unverändert} \\
\hline
(3) Das Jugendamt, das in Verfahren zur \newline Übertragung der gemeinsamen Sorge nach \newline \S{} 155a Absatz 4 Satz 1 und \S{} 162 des \newline Gesetzes über das Verfahren in \newline Familiensachen und in den \newline Angelegenheiten der freiwilligen \newline Gerichtsbarkeit angehört wird, teilt & \textit{(3) unverändert} & \textit{(3) unverändert} \\
\hline
1. \newline rechtskräftige gerichtliche Entscheidungen, \newline aufgrund derer die Sorge gemäß \S{} 1626a \newline Absatz 2 Satz 1 des Bürgerlichen \newline Gesetzbuchs den Eltern ganz oder zum \newline Teil gemeinsam übertragen wird oder & 1. unverändert & 1. unverändert \\
\hline
2. \newline rechtskräftige gerichtliche Entscheidungen, \newline die die elterliche Sorge ganz oder zum Teil \newline der Mutter entziehen oder auf den Vater \newline allein übertragen, & 2. unverändert & 2. unverändert \\
\hline
\multicolumn{3}{|c|}{\cellcolor{gray!10}\textbf{\S{} 58a -- Bestätigung eines Todes eines sorgeberechtigten Elternteils}} \\
\hline
 & \textit{unverändert} & \textbf{Ein Elternteil, dem die elterliche Sorge \newline oder ein Teil der elterlichen Sorge \newline bislang gemeinsam mit dem anderen \newline Elternteil zustand und dem infolge des \newline Todes des anderen Elternteils die \newline elterliche Sorge oder ein Teil der \newline elterlichen Sorge allein zusteht, kann \newline bei dem nach \S{} 87c Absatz 6 Satz 2 \newline zuständigen Jugendamt hierüber eine \newline Bestätigung beantragen. Das \newline Jugendamt kann die Bestätigung \newline ablehnen, wenn berechtigte Zweifel an \newline der alleinigen Sorge des Elternteils oder \newline der Anwendbarkeit deutschen Rechts \newline bestehen. } \\
\hline
\multicolumn{3}{|c|}{\cellcolor{gray!10}\textbf{\S{} 74 -- Förderung der freien Jugendhilfe}} \\
\hline
(1) Die Träger der öffentlichen Jugendhilfe \newline sollen die freiwillige Tätigkeit auf dem \newline Gebiet der Jugendhilfe anregen; sie sollen \newline sie fördern, wenn der jeweilige Träger & \textit{(1) unverändert} & \textit{(1) unverändert} \\
\hline
1. \newline die fachlichen Voraussetzungen für die \newline geplante Maßnahme erfüllt und die \newline Beachtung der Grundsätze und Maßstäbe \newline der Qualitätsentwicklung und \newline Qualitätssicherung nach \S{} 79a \newline gewährleistet, & 1. unverändert & 1. unverändert \\
\hline
2. \newline die Gewähr für eine zweckentsprechende \newline und wirtschaftliche Verwendung der Mittel \newline bietet, & 2. unverändert & 2. unverändert \\
\hline
3. \newline gemeinnützige Ziele verfolgt, & 3. unverändert & 3. unverändert \\
\hline
4. \newline eine angemessene Eigenleistung erbringt \newline und & 4. unverändert & 4. unverändert \\
\hline
5. \newline die Gewähr für eine den Zielen des \newline Grundgesetzes förderliche Arbeit bietet. & 5. unverändert & 5. unverändert \\
\hline
\textit{--- Fassung 2024 ---} \newline Eine auf Dauer angelegte Förderung setzt \newline in der Regel die Anerkennung als Träger \newline der freien Jugendhilfe nach \S{} 75 voraus. \newline \textit{--- Fassung 2026 ---} \newline Ausgestaltung der Maßnahme \newline gewährleisten. & \textit{unverändert} & Einflussnahme auf die Ausgestaltung der \newline Maßnahme gewährleisten. \\
\hline
(2) Soweit von der freien Jugendhilfe \newline Einrichtungen, Dienste und \newline Veranstaltungen geschaffen werden, um \newline die Gewährung von Leistungen nach \newline diesem Buch zu ermöglichen, kann die \newline Förderung von der Bereitschaft abhängig \newline gemacht werden, diese Einrichtungen, \newline Dienste und Veranstaltungen nach \newline Maßgabe der Jugendhilfeplanung und \newline unter Beachtung der in \S{} 9 genannten \newline Grundsätze anzubieten. \S{} 4 Absatz 1 bleibt \newline unberührt. & \textit{(2) unverändert} & \textit{(2) unverändert} \\
\hline
(3) Über die Art und Höhe der Förderung \newline entscheidet der Träger der öffentlichen \newline Jugendhilfe im Rahmen der verfügbaren \newline Haushaltsmittel nach pflichtgemäßem \newline Ermessen. Entsprechendes gilt, wenn \newline mehrere Antragsteller die \newline Förderungsvoraussetzungen erfüllen und \newline die von ihnen vorgesehenen Maßnahmen \newline gleich geeignet sind, zur Befriedigung des \newline Bedarfs jedoch nur eine Maßnahme \newline notwendig ist. Bei der Bemessung der \newline Eigenleistung sind die unterschiedliche \newline Finanzkraft und die sonstigen Verhältnisse \newline zu berücksichtigen. & \textit{(3) unverändert} & \textit{(3) unverändert} \\
\hline
\textit{--- Fassung 2024 ---} \newline (4) Bei sonst gleich geeigneten \newline Maßnahmen soll solchen der Vorzug \newline gegeben werden, die stärker an den \newline Interessen der Betroffenen orientiert sind \newline und ihre Einflussnahme auf die \newline Ausgestaltung der Maßnahme \newline gewährleisten. \newline \textit{--- Fassung 2026 ---} \newline (4) Bei sonst gleich geeigneten \newline Maßnahmen soll solchen der Vorzug \newline gegeben werden, die stärker an den \newline Interessen der Betroffenen orientiert sind \newline und ihre Einflussnahme auf die & (4) Bei sonst gleich geeigneten \newline Maßnahmen soll solchen der Vorzug \newline gegeben werden, die stärker \textbf{inklusiv \newline ausgerichtet oder }an den Interessen der \newline Betroffenen orientiert sind und ihre \newline Einflussnahme auf die Ausgestaltung der \newline Maßnahme gewährleisten. & (4) Bei sonst gleich geeigneten \newline Maßnahmen soll solchen der Vorzug \newline gegeben werden, die stärker \textbf{inklusiv \newline ausgerichtet oder }an den Interessen der \newline Betroffenen orientiert sind und ihre \\
\hline
(5) Bei der Förderung gleichartiger \newline Maßnahmen mehrerer Träger sind unter \newline Berücksichtigung ihrer Eigenleistungen \newline gleiche Grundsätze und Maßstäbe \newline anzulegen. Werden gleichartige \newline Maßnahmen von der freien und der \newline öffentlichen Jugendhilfe durchgeführt, so \newline sind bei der Förderung die Grundsätze und \newline Maßstäbe anzuwenden, die für die \newline Finanzierung der Maßnahmen der \newline öffentlichen Jugendhilfe gelten. & \textit{(5) unverändert} & \textit{(5) unverändert} \\
\hline
(6) Die Förderung von anerkannten \newline Trägern der Jugendhilfe soll auch Mittel für \newline die Fortbildung der haupt-, neben- und \newline ehrenamtlichen Mitarbeiter sowie im \newline Bereich der Jugendarbeit Mittel für die \newline Errichtung und Unterhaltung von \newline Jugendfreizeit- und Jugendbildungsstätten \newline einschließen. & \textit{(6) unverändert} & \textit{(6) unverändert} \\
\hline
\multicolumn{3}{|c|}{\cellcolor{gray!10}\textbf{\S{} 75 -- Anerkennung als Träger der freien Jugendhilfe}} \\
\hline
\textit{--- Fassung 2024 ---} \newline (1) Als Träger der freien Jugendhilfe \newline können juristische Personen und \newline Personenvereinigungen anerkannt werden, \newline \textit{--- Fassung 2026 ---} \newline (1) Als Träger der freien Jugendhilfe \newline können juristische Personen und \newline Personenvereinigungen anerkannt werden, \newline wenn sie \newline die Gewähr für eine den Zielen des \newline Grundgesetzes förderliche Arbeit bieten. & \textit{(1) unverändert} & \textit{(1) unverändert} \\
\hline
1. \newline auf dem Gebiet der Jugendhilfe im Sinne \newline des \S{} 1 tätig sind, & 1. unverändert & 1. unverändert \\
\hline
\textit{--- Fassung 2024 ---} \newline 2. \newline \textit{--- Fassung 2026 ---} \newline 2. \newline gemeinnützige Ziele verfolgen, & 2. unverändert & 2. unverändert \\
\hline
3. \newline auf Grund der fachlichen und personellen \newline Voraussetzungen erwarten lassen, dass \newline sie einen nicht unwesentlichen Beitrag zur \newline Erfüllung der Aufgaben der Jugendhilfe zu \newline leisten imstande sind, und & 3. unverändert & 3. unverändert \\
\hline
\textit{--- Fassung 2024 ---} \newline 4. \newline die Gewähr für eine den Zielen des \newline \textit{--- Fassung 2026 ---} \newline 4. & 4. unverändert & 4. unverändert \\
\hline
(2) Einen Anspruch auf Anerkennung als \newline Träger der freien Jugendhilfe hat unter den \newline Voraussetzungen des Absatzes 1, wer auf \newline dem Gebiet der Jugendhilfe mindestens \newline drei Jahre tätig gewesen ist. & (2) Einen Anspruch auf Anerkennung als \newline Träger der freien Jugendhilfe hat unter den \newline Voraussetzungen des Absatzes 1, wer auf \newline dem Gebiet der Jugendhilfe \textbf{oder der \newline Eingliederungshilfe für Kinder und \newline Jugendliche mit Behinderungen}mindestens drei Jahre tätig gewesen ist. & (2) Einen Anspruch auf Anerkennung als \newline Träger der freien Jugendhilfe hat unter den \newline Voraussetzungen des Absatzes 1, wer auf \newline dem Gebiet der Jugendhilfe \textbf{oder der \newline Eingliederungshilfe für Kinder und \newline Jugendliche mit Behinderungen}mindestens drei Jahre tätig gewesen ist. \\
\hline
(3) Die Kirchen und \newline Religionsgemeinschaften des öffentlichen \newline Rechts sowie die auf Bundesebene \newline zusammengeschlossenen Verbände der \newline freien Wohlfahrtspflege sind anerkannte \newline Träger der freien Jugendhilfe. & \textit{(3) unverändert} & \textit{(3) unverändert} \\
\hline
\multicolumn{3}{|c|}{\cellcolor{gray!10}\textbf{\S{} 77 -- Vereinbarungen über Kostenübernahme und Qualitätsentwicklung bei ambulanten Leistungen}} \\
\hline
(1) Werden Einrichtungen und Dienste der \newline Träger der freien Jugendhilfe in Anspruch \newline genommen, so sind Vereinbarungen über \newline die Höhe der Kosten der Inanspruchnahme \newline sowie über Inhalt, Umfang und Qualität der \newline Leistung, über Grundsätze und Maßstäbe \newline für die Bewertung der Qualität der Leistung \newline und über geeignete Maßnahmen zu ihrer \newline Gewährleistung zwischen der öffentlichen \newline und der freien Jugendhilfe anzustreben. Zu \newline den Grundsätzen und Maßstäben für die \newline Bewertung der Qualität der Leistung nach \newline Satz 1 zählen auch Qualitätsmerkmale für \newline die inklusive Ausrichtung der \newline Aufgabenwahrnehmung und die \newline Berücksichtigung der spezifischen \newline Bedürfnisse von jungen Menschen mit \newline Behinderungen. Das Nähere regelt das \newline Landesrecht. Die \S{}\S{} 78a bis 78g bleiben \newline unberührt. & \textit{(1) unverändert} & (1) Werden Einrichtungen und Dienste der \newline Träger der freien Jugendhilfe in Anspruch \newline genommen, so sind Vereinbarungen über \newline die Höhe der Kosten der Inanspruchnahme \newline sowie über Inhalt, Umfang und Qualität der \newline Leistung, über Grundsätze und Maßstäbe \newline für die Bewertung der Qualität der Leistung \newline und über geeignete Maßnahmen zu ihrer \newline Gewährleistung zwischen \textbf{dem Träger der \newline öffentlichen Jugendhilfe und unter \newline Berücksichtigung der \newline Leistungsfähigkeit, Wirtschaftlichkeit \newline und Sparsamkeit geeigneten freien \newline Trägern }anzustreben. Zu den Grundsätzen \newline und Maßstäben für die Bewertung der \newline Qualität der Leistung nach Satz 1 zählen \newline auch Qualitätsmerkmale für die inklusive \newline Ausrichtung der Aufgabenwahrnehmung \newline und die Berücksichtigung der spezifischen \newline Bedürfnisse von jungen Menschen mit \newline Behinderungen. Das Nähere regelt das \newline Landesrecht. Die \S{}\S{} 78a bis 78g bleiben \newline unberührt. \\
\hline
\textit{--- Fassung 2024 ---} \newline (2) Wird eine Leistung nach \S{} 37 Absatz 1 \newline oder \S{} 37a erbracht, so ist der Träger der \newline öffentlichen Jugendhilfe zur Übernahme \newline der Kosten der Inanspruchnahme nur \newline verpflichtet, wenn mit den \newline Leistungserbringern Vereinbarungen über \newline Inhalt, Umfang und Qualität der Leistung, \newline über Grundsätze und Maßstäbe für die \newline Bewertung der Qualität der Leistung sowie \newline über geeignete Maßnahmen zu ihrer \newline Gewährleistung geschlossen worden sind; \newline \S{} 78e gilt entsprechend. \newline \textit{--- Fassung 2026 ---} \newline (2) Wird eine Leistung nach \S{} 37 Absatz 1 \newline oder \S{} 37a erbracht, so ist der Träger der \newline öffentlichen Jugendhilfe zur Übernahme \newline der Kosten der Inanspruchnahme nur & (2) Wird eine Leistung nach \S{} 3\textbf{9} Absatz 1 \newline oder \S{} 3\textbf{9a} erbracht, so ist der Träger der \newline öffentlichen Jugendhilfe zur Übernahme \newline der Kosten der Inanspruchnahme nur \newline verpflichtet, wenn mit den \newline Leistungserbringern Vereinbarungen über \newline Inhalt, Umfang und Qualität der Leistung, \newline über Grundsätze und Maßstäbe für die \newline Bewertung der Qualität der Leistung sowie \newline über geeignete Maßnahmen zu ihrer \newline Gewährleistung geschlossen worden sind; \newline \S{} 78e gilt entsprechend. & (2) Wird eine Leistung nach \S{} 3\textbf{9} Absatz 1 \newline oder \S{} 3\textbf{9a} erbracht, so ist der Träger der \newline öffentlichen Jugendhilfe zur Übernahme \newline der Kosten der Inanspruchnahme nur \\
\hline
verpflichtet, wenn mit den \newline Leistungserbringern Vereinbarungen über \newline Inhalt, Umfang und Qualität der Leistung, \newline über Grundsätze und Maßstäbe für die \newline Bewertung der Qualität der Leistung sowie \newline über geeignete Maßnahmen zu ihrer \newline Gewährleistung geschlossen worden sind; \newline \S{} 78e gilt entsprechend. & \textit{unverändert} & verpflichtet, wenn mit den \newline Leistungserbringern Vereinbarungen über \newline Inhalt, Umfang und Qualität der Leistung, \newline über Grundsätze und Maßstäbe für die \newline Bewertung der Qualität der Leistung sowie \newline über geeignete Maßnahmen zu ihrer \newline Gewährleistung geschlossen worden sind; \newline \S{} 78e gilt entsprechend. \\
\hline
\multicolumn{3}{|c|}{\cellcolor{gray!10}\textbf{\S{} 78a -- Anwendungsbereich}} \\
\hline
(1) Die Regelungen der \S{}\S{} 78b bis 78g \newline gelten für die Erbringung von & \textit{(1) unverändert} & \textit{(1) unverändert} \\
\hline
1. \newline Leistungen für Betreuung und Unterkunft in \newline einer sozialpädagogisch begleiteten \newline Wohnform (\S{} 13 Absatz 3), & 1. unverändert & 1. unverändert \\
\hline
2. \newline Leistungen in gemeinsamen Wohnformen \newline für Mütter/Väter und Kinder (\S{} 19), & 2. unverändert & 2. unverändert \\
\hline
3. \newline Leistungen zur Unterstützung bei \newline notwendiger Unterbringung des Kindes \newline oder Jugendlichen zur Erfüllung der \newline Schulpflicht (\S{} 21 Satz 2), & 3. unverändert & 3. unverändert \\
\hline
4. \newline Hilfe zur Erziehung & 4. unverändert & 4. \newline Hilfe zur Erziehung \\
\hline
b) \newline Einrichtungen über Tag und Nacht sowie \newline sonstigen Wohnformen (\S{} 35a Absatz 2 \newline Nummer 4), & \textbf{b) \newline Einrichtungen über Tag und Nacht \newline sowie sonstigen Wohnformen (\S{} 35a \newline Absatz 4 Nummer 4) } & \textbf{b) \newline Einrichtungen über Tag und Nacht \newline sowie sonstigen Wohnformen (\S{} 35a \newline Absatz 4 Nummer 4) } \\
\hline
6. \newline Hilfe für junge Volljährige (\S{} 41), sofern \newline diese den in den Nummern 4 und 5 \newline genannten Leistungen entspricht, sowie & 6. unverändert & 6. unverändert \\
\hline
7. \newline Leistungen zum Unterhalt (\S{} 39), sofern \newline diese im Zusammenhang mit Leistungen \newline nach den Nummern 4 bis 6 gewährt \newline werden; \S{} 39 Absatz 2 Satz 3 bleibt \newline unberührt. & 7. \newline Leistungen zum Unterhalt (\S{} 39\textbf{c}), sofern \newline diese im Zusammenhang mit Leistungen \newline nach den Nummern 4 bis 6 gewährt \newline werden; \S{} 39\textbf{c} Absatz 2 Satz 3 bleibt \newline unberührt. & 7. \newline Leistungen zum Unterhalt (\S{} 39\textbf{c}), sofern \newline diese im Zusammenhang mit Leistungen \newline nach den Nummern 4 bis 6 gewährt \newline werden; \S{} 39\textbf{c} Absatz 2 Satz 3 bleibt \newline unberührt. \\
\hline
(2) Landesrecht kann bestimmen, dass die \newline \S{}\S{} 78b bis 78g auch für andere Leistungen \newline nach diesem Buch sowie für vorläufige \newline Maßnahmen zum Schutz von Kindern und \newline Jugendlichen (\S{}\S{} 42, 42a) gelten. & \textit{(2) unverändert} & \textit{(2) unverändert} \\
\hline
\multicolumn{3}{|c|}{\cellcolor{gray!10}\textbf{\S{} 78b -- Voraussetzungen für die Übernahme des Leistungsentgelts}} \\
\hline
(1) Wird die Leistung ganz oder teilweise in \newline einer Einrichtung erbracht, so ist der Träger \newline der öffentlichen Jugendhilfe zur Übernahme \newline des Entgelts gegenüber dem \newline Leistungsberechtigten verpflichtet, wenn mit \newline dem Träger der Einrichtung oder seinem \newline Verband Vereinbarungen über & \textit{(1) unverändert} & \textit{(1) unverändert} \\
\hline
\textit{--- Fassung 2024 ---} \newline 1. \newline Inhalt, Umfang und Qualität der \newline Leistungsangebote \newline \textit{--- Fassung 2026 ---} \newline 1. \newline Inhalt, Umfang und Qualität der \newline Leistungsangebote \newline (Leistungsvereinbarung), & 1. unverändert & 1. \newline Inhalt, Umfang und Qualität\textbf{ einschließlich \newline der Wirksamkeit} der Leistungsangebote \newline (Leistungsvereinbarung), \\
\hline
2. \newline differenzierte Entgelte für die \newline Leistungsangebote und die \newline betriebsnotwendigen Investitionen \newline (Entgeltvereinbarung) und & 2. unverändert & 2. unverändert \\
\hline
3. \newline Grundsätze und Maßstäbe für die \newline Bewertung der Qualität der \newline Leistungsangebote sowie über geeignete \newline Maßnahmen zu ihrer Gewährleistung \newline (Qualitätsentwicklungsvereinbarung) \newline abgeschlossen worden sind; dazu zählen \newline auch die Qualitätsmerkmale nach \S{} 79a \newline Satz 2. & 3. unverändert & 3. unverändert \\
\hline
(2) Die Vereinbarungen sind mit den \newline Trägern abzuschließen, die unter \newline Berücksichtigung der Grundsätze der \newline Leistungsfähigkeit, Wirtschaftlichkeit und \newline Sparsamkeit zur Erbringung der Leistung \newline geeignet sind. Vereinbarungen über die \newline Erbringung von Auslandsmaßnahmen \newline dürfen nur mit solchen Trägern \newline abgeschlossen werden, die die Maßgaben \newline nach \S{} 38 Absatz 2 Nummer 2 Buchstabe \newline a bis d erfüllen. & (2) Die Vereinbarungen sind mit den \newline Trägern abzuschließen, die unter \newline Berücksichtigung der Grundsätze der \newline Leistungsfähigkeit, Wirtschaftlichkeit und \newline Sparsamkeit zur Erbringung der Leistung \newline geeignet sind. Vereinbarungen über die \newline Erbringung von Auslandsmaßnahmen \newline dürfen nur mit solchen Trägern \newline abgeschlossen werden, die die Maßgaben \newline nach \textbf{\S{} 40 Absatz 2 Nummer 2 \newline Buchstabe a bis d} erfüllen. & (2) Die Vereinbarungen sind mit unter \newline Berücksichtigung der Leistungsfähigkeit, \newline Wirtschaftlichkeit und Sparsamkeit \newline geeigneten Trägern abzuschließen, die \newline eine bedarfsdeckende Leistungserbringung \newline nach den Besonderheiten des Einzelfalls \newline unter Berücksichtigung des Wunsch- und \newline Wahlrechts des Leistungsberechtigten \newline nach \S{} 5 sicherstellen. Vereinbarungen \newline über die Erbringung von \newline Auslandsmaßnahmen dürfen nur mit \newline solchen Trägern abgeschlossen werden, \newline die die Maßgaben nach \textbf{\S{} 40 Absatz 2 \newline Nummer 2 Buchstabe a bis d} erfüllen. \\
\hline
 & \textit{unverändert} & \textbf{(2a) Die Ergebnisse der Vereinbarungen \newline sind den Leistungsberechtigten in einer \newline für sie verständlichen, \newline nachvollziehbaren und wahrnehmbaren \newline Form zugänglich zu machen. } \\
\hline
(3) Ist eine der Vereinbarungen nach \newline Absatz 1 nicht abgeschlossen, so ist der \newline Träger der öffentlichen Jugendhilfe zur \newline Übernahme des Leistungsentgelts nur \newline verpflichtet, wenn dies insbesondere nach \newline Maßgabe der Hilfeplanung (\S{} 36) im \newline Einzelfall geboten ist. & (3) Ist eine der Vereinbarungen nach \newline Absatz 1 nicht abgeschlossen, so ist der \newline Träger der öffentlichen Jugendhilfe zur \newline Übernahme des Leistungsentgelts nur \newline verpflichtet, wenn dies insbesondere nach \newline Maßgabe \textbf{des Hilfe- und Leistungsplans \newline (\S{}\S{} 36a, 37, 38c)} im Einzelfall geboten ist. & (3) Ist eine der Vereinbarungen nach \newline Absatz 1 nicht abgeschlossen, so ist der \newline Träger der öffentlichen Jugendhilfe zur \newline Übernahme des Leistungsentgelts nur \newline verpflichtet, wenn dies insbesondere nach \newline Maßgabe \textbf{des Hilfe- und Leistungsplans \newline (\S{}\S{} 36a, 37, 38c)} im Einzelfall geboten ist. \\
\hline
 & \textit{unverändert} & \textbf{(4) Liegen die Voraussetzungen für die \newline Übernahme des Leistungsentgelts nach \newline Absatz 1 und 3 vor, hat der \newline Leistungserbringer, der eine bewilligte \newline Leistung gegenüber dem \newline Leistungsberechtigten erbringt, \newline Anspruch auf Vergütung dieser \newline Leistung gegenüber dem Träger der \newline öffentlichen Jugendhilfe. } \\
\hline
\multicolumn{3}{|c|}{\cellcolor{gray!10}\textbf{\S{} 78c -- Inhalt der Leistungs- und Entgeltvereinbarungen}} \\
\hline
(1) Die Leistungsvereinbarung muss die \newline wesentlichen Leistungsmerkmale, \newline insbesondere \newline 1. Art, Ziel und Qualität des \newline Leistungsangebots, \newline 2. den in der Einrichtung zu betreuenden \newline Personenkreis, \newline 3. die erforderliche sächliche und \newline personelle Ausstattung, \newline 4. die Qualifikation des Personals sowie \newline 5. die betriebsnotwendigen Anlagen der \newline Einrichtung \newline festlegen. In die Vereinbarung ist \newline aufzunehmen, unter welchen \newline Voraussetzungen der Träger der \newline Einrichtung sich zur Erbringung von \newline Leistungen verpflichtet. Der Träger muss \newline gewährleisten, dass die Leistungsangebote \newline zur Erbringung von Leistungen nach \S{} 78a \newline Absatz 1 geeignet sowie ausreichend, \newline zweckmäßig und wirtschaftlich sind. & \textit{unverändert} & (1) Die Leistungsvereinbarung muss die \newline wesentlichen Leistungsmerkmale, \newline insbesondere \newline 1. Art, Ziel und Qualität des \newline Leistungsangebots, \newline 2. den in der Einrichtung zu betreuenden \newline Personenkreis, \newline 3. die erforderliche sächliche und \newline personelle Ausstattung, \newline 4. die Qualifikation des Personals, \textbf{die sich \newline nach der Zweckbestimmung der \newline Einrichtung und den jeweils in ihr \newline wahrzunehmenden Funktionen richtet,} \newline sowie \newline 5. die betriebsnotwendigen Anlagen der \newline Einrichtung \newline festlegen. In die Vereinbarung ist \newline aufzunehmen, unter welchen \newline Voraussetzungen der Träger der \newline Einrichtung sich zur Erbringung von \newline Leistungen verpflichtet. Der Träger muss \newline gewährleisten, dass die Leistungsangebote \newline zur Erbringung von Leistungen nach \S{} 78a \newline Absatz 1 geeignet sowie ausreichend, \newline zweckmäßig und wirtschaftlich sind. \\
\hline
Absatz 2 & \textit{unverändert} & \textit{unverändert} \\
\hline
\multicolumn{3}{|c|}{\cellcolor{gray!10}\textbf{\S{} 78g -- Absatz 1}} \\
\hline
(2) Kommt eine Vereinbarung nach \S{} 78b \newline Absatz 1 innerhalb von sechs Wochen \newline nicht zustande, nachdem eine Partei \newline schriftlich zu Verhandlungen aufgefordert \newline hat, so entscheidet die Schiedsstelle auf \newline Antrag einer Partei unverzüglich über die \newline Gegenstände, über die keine Einigung \newline erreicht werden konnte. Gegen die \newline Entscheidung ist der Rechtsweg zu den \newline Verwaltungsgerichten gegeben. Die Klage \newline richtet sich gegen eine der beiden \newline Vertragsparteien, nicht gegen die \newline Schiedsstelle. Einer Nachprüfung der \newline Entscheidung in einem Vorverfahren bedarf \newline es nicht. & \textit{unverändert} & (2) Kommt eine Vereinbarung nach \S{} 78b \newline Absatz 1 innerhalb von sechs Wochen \newline nicht zustande, nachdem eine Partei \newline schriftlich zu Verhandlungen aufgefordert \newline hat, so entscheidet die Schiedsstelle auf \newline Antrag einer Partei unverzüglich über die \newline Gegenstände, über die keine Einigung \newline erreicht werden konnte. Gegen die \newline Entscheidung ist der Rechtsweg zu den \newline \textbf{Sozialgerichten }gegeben. Die Klage \newline richtet sich gegen eine der beiden \newline Vertragsparteien, nicht gegen die \newline Schiedsstelle. Einer Nachprüfung der \newline Entscheidung in einem Vorverfahren bedarf \newline es nicht. \\
\hline
Absätze 3 und 4 & \textit{unverändert} & \textit{unverändert} \\
\hline
\multicolumn{3}{|c|}{\cellcolor{gray!10}\textbf{\S{} 79 -- Gesamtverantwortung}} \\
\hline
Absätze 1 und 2 & \textit{unverändert} & \textit{unverändert} \\
\hline
(3) Die Träger der öffentlichen Jugendhilfe \newline haben für eine ausreichende Ausstattung \newline der Jugendämter und der \newline Landesjugendämter einschließlich der \newline Möglichkeit der Nutzung digitaler Geräte zu \newline sorgen; hierzu gehört auch eine dem \newline Bedarf entsprechende Zahl von \newline Fachkräften. Zur Planung und \newline Bereitstellung einer bedarfsgerechten \newline Personalausstattung ist ein Verfahren zur \newline Personalbemessung zu nutzen. & \textit{unverändert} & (3) Die Träger der öffentlichen Jugendhilfe \newline haben für eine ausreichende Ausstattung \newline der Jugendämter und der \newline Landesjugendämter einschließlich der \newline Möglichkeit der Nutzung digitaler Geräte zu \newline sorgen; hierzu gehört auch eine dem \textbf{den \newline jeweiligen Aufgabenbereichen und den \newline darin jeweils wahrzunehmenden \newline Funktionen} entsprechende Zahl von \newline Fachkräften. Zur Planung und \newline Bereitstellung einer bedarfsgerechten \newline Personalausstattung ist ein Verfahren zur \newline Personalbemessung zu nutzen. \\
\hline
 & \textit{unverändert} & \textbf{(4) Für die Bestimmung der örtlichen \newline Zuständigkeit nach \S{}\S{} 86 bis 88a soll ein \newline von der fachlich zuständigen obersten \newline Bundesbehörde bereitgestelltes \newline automatisiertes Prüfungssystem zur \newline Anwendung kommen. } \\
\hline
\multicolumn{3}{|c|}{\cellcolor{gray!10}\textbf{\S{} 80 -- Jugendhilfeplanung}} \\
\hline
(1) Die Träger der öffentlichen Jugendhilfe \newline haben im Rahmen ihrer \newline Planungsverantwortung & \textit{(1) unverändert} & \textit{(1) unverändert} \\
\hline
1. \newline den Bestand an Einrichtungen und \newline Diensten festzustellen, & 1. unverändert & 1. unverändert \\
\hline
2. \newline den Bedarf unter Berücksichtigung der \newline Wünsche, Bedürfnisse und Interessen der \newline jungen Menschen und der \newline Erziehungsberechtigten für einen \newline mittelfristigen Zeitraum zu ermitteln und & 2. unverändert & 2. unverändert \\
\hline
\textit{--- Fassung 2024 ---} \newline 3. \newline die zur Befriedigung des Bedarfs \newline notwendigen Vorhaben rechtzeitig und \newline ausreichend zu planen; dabei ist Vorsorge \newline zu treffen, dass auch ein \newline unvorhergesehener Bedarf befriedigt \newline \textit{--- Fassung 2026 ---} \newline 3. \newline die zur Befriedigung des Bedarfs \newline notwendigen Vorhaben rechtzeitig und \newline ausreichend zu planen; dabei ist Vorsorge \newline zu treffen, dass auch ein \newline unvorhergesehener Bedarf befriedigt \newline werden kann. & 3. unverändert & 3. unverändert \\
\hline
\textit{--- Fassung 2024 ---} \newline (2) Einrichtungen und Dienste sollen so \newline geplant werden, dass insbesondere \newline \textit{--- Fassung 2026 ---} \newline (2) Einrichtungen und Dienste sollen so \newline geplant werden, dass insbesondere \newline Behinderung gemeinsam unter \newline Berücksichtigung spezifischer \newline Bedarfslagen gefördert werden können, & \textit{(2) unverändert} & \textit{(2) unverändert} \\
\hline
1. \newline Kontakte in der Familie und im sozialen \newline Umfeld erhalten und gepflegt werden \newline können, & 1. unverändert & 1. unverändert \\
\hline
2. \newline ein möglichst wirksames, vielfältiges, \newline inklusives und aufeinander abgestimmtes \newline Angebot von Jugendhilfeleistungen \newline gewährleistet ist, & 2. unverändert & 2. unverändert \\
\hline
3. \newline ein dem nach Absatz 1 Nummer 2 \newline ermittelten Bedarf entsprechendes \newline Zusammenwirken der Angebote von \newline Jugendhilfeleistungen in den Lebens- und \newline Wohnbereichen von jungen Menschen und \newline Familien sichergestellt ist, & 3. unverändert & 3. unverändert \\
\hline
\textit{--- Fassung 2024 ---} \newline 4. \newline junge Menschen mit Behinderungen oder \newline von Behinderung bedrohte junge \newline Menschen mit jungen Menschen ohne \newline Behinderung gemeinsam unter \newline Berücksichtigung spezifischer \newline Bedarfslagen gefördert werden können, \newline \textit{--- Fassung 2026 ---} \newline 4. \newline junge Menschen mit Behinderungen oder \newline von Behinderung bedrohte junge \newline Menschen mit jungen Menschen ohne & 4. unverändert & 4. unverändert \\
\hline
5. \newline junge Menschen und Familien in \newline gefährdeten Lebens- und Wohnbereichen \newline besonders gefördert werden, & 5. unverändert & 5. unverändert \\
\hline
6. \newline Mütter und Väter Aufgaben in der Familie \newline und Erwerbstätigkeit besser miteinander \newline vereinbaren können. & 6. unverändert & 6. unverändert \\
\hline
(3) Die Planung insbesondere von \newline Diensten zur Gewährung niedrigschwelliger \newline ambulanter Hilfen nach Maßgabe von \S{} \newline 36a Absatz 2 umfasst auch Maßnahmen \newline zur Qualitätsgewährleistung der \newline Leistungserbringung. & (3) Die Planung insbesondere von \newline Diensten zur Gewährung niedrigschwelliger \newline ambulanter Hilfen nach Maßgabe von \S{} \newline \textbf{36c Absatz 2} umfasst auch Maßnahmen \newline zur Qualitätsgewährleistung der \newline Leistungserbringung. & (3) Die Planung insbesondere von \newline Diensten zur Gewährung niedrigschwelliger \newline ambulanter Hilfen nach Maßgabe von \S{} \newline \textbf{36c Absatz 2} umfasst auch Maßnahmen \newline zur Qualitätsgewährleistung der \newline Leistungserbringung. \\
\hline
(4) Die Träger der öffentlichen Jugendhilfe \newline haben die anerkannten Träger der freien \newline Jugendhilfe in allen Phasen ihrer Planung \newline frühzeitig zu beteiligen. Zu diesem Zwecke \newline sind sie vom Jugendhilfeausschuss, soweit \newline sie überörtlich tätig sind, im Rahmen der \newline Jugendhilfeplanung des überörtlichen \newline Trägers vom Landesjugendhilfeausschuss \newline zu hören. Das Nähere regelt das \newline Landesrecht. & \textit{(4) unverändert} & \textit{(4) unverändert} \\
\hline
(5) Die Träger der öffentlichen Jugendhilfe \newline sollen darauf hinwirken, dass die \newline Jugendhilfeplanung und andere örtliche \newline und überörtliche Planungen aufeinander \newline abgestimmt werden und die Planungen \newline insgesamt den Bedürfnissen und \newline Interessen der jungen Menschen und ihrer \newline Familien Rechnung tragen. & \textit{(5) unverändert} & \textit{(5) unverändert} \\
\hline
\multicolumn{3}{|c|}{\cellcolor{gray!10}\textbf{\S{} 80a -- Planung infrastruktureller Bildungsassistenz}} \\
\hline
 & \textit{unverändert} & \textbf{Infrastrukturelle Angebote zu einer \newline wegen erzieherischen Bedarfs oder \newline wegen einer Behinderung erforderlichen \newline Anleitung und Begleitung von Kindern \newline und Jugendlichen in \newline Tageseinrichtungen nach \S{} 22 Absatz 1 \newline Satz 1 oder in Schulen oder \newline Hochschulen sind Angebote, die nach \newline Maßgabe von \S{} 80 unter Einbeziehung \newline der nach Landesrecht für die Schulen \newline und Hochschulen zuständigen \newline Behörden geplant werden. Landesrecht \newline regelt das Nähere, insbesondere auch \newline die Beteiligung der in den Ländern für \newline die Finanzierung der Schulen oder \newline Hochschulen verantwortlichen Stellen \newline an den Kosten der Angebote nach Satz \newline 1. } \\
\hline
\multicolumn{3}{|c|}{\cellcolor{gray!10}\textbf{\S{} 84 -- Jugendbericht}} \\
\hline
(1) Die Bundesregierung legt dem \newline Deutschen Bundestag und dem Bundesrat \newline in jeder Legislaturperiode einen Bericht \newline über die Lage junger Menschen und die \newline Bestrebungen und Leistungen der \newline Jugendhilfe vor. Neben der \newline Bestandsaufnahme und Analyse sollen die \newline Berichte Vorschläge zur Weiterentwicklung \newline der Jugendhilfe enthalten; jeder dritte \newline Bericht soll einen Überblick über die \newline Gesamtsituation der Jugendhilfe vermitteln. & \textit{unverändert} & (1) Die Bundesregierung legt dem \newline Deutschen Bundestag und dem Bundesrat \newline in jeder Legislaturperiode einen Bericht \newline über die Lage junger Menschen und die \newline Bestrebungen und Leistungen der \newline Jugendhilfe vor. \textbf{Neben der \newline Bestandsaufnahme und Analyse sollen \newline die Berichte Vorschläge zur \newline Weiterentwicklung der Jugendhilfe \newline enthalten; jeder dritte Bericht soll einen \newline Überblick über die Gesamtsituation der \newline Jugendhilfe vermitteln und auch die \newline Situation unbegleiteter ausländischer \newline Minderjähriger darstellen.} \\
\hline
Absatz 2 & \textit{unverändert} & \textit{unverändert} \\
\hline
\multicolumn{3}{|c|}{\cellcolor{gray!10}\textbf{\S{} 85 -- Sachliche Zuständigkeit}} \\
\hline
(1) Für die Gewährung von Leistungen und \newline die Erfüllung anderer Aufgaben nach \newline diesem Buch ist der örtliche Träger \newline sachlich zuständig, soweit nicht der \newline überörtliche Träger sachlich zuständig ist. & \textit{(1) unverändert} & \textit{(1) unverändert} \\
\hline
\textit{--- Fassung 2024 ---} \newline (2) Der überörtliche Träger ist sachlich \newline zuständig für \newline \textit{--- Fassung 2026 ---} \newline (2) Der überörtliche Träger ist sachlich \newline zuständig für \newline die Wahrnehmung der Aufgaben zum \newline Schutz von Kindern und Jugendlichen in \newline Einrichtungen (\S{}\S{} 45 bis 48a), & \textit{(2) unverändert} & \textit{(2) unverändert} \\
\hline
1. \newline die Beratung der örtlichen Träger und die \newline Entwicklung von Empfehlungen zur \newline Erfüllung der Aufgaben nach diesem Buch, & 1. unverändert & 1. unverändert \\
\hline
2. \newline die Förderung der Zusammenarbeit \newline zwischen den örtlichen Trägern und den \newline anerkannten Trägern der freien \newline Jugendhilfe, insbesondere bei der Planung \newline und Sicherstellung eines bedarfsgerechten \newline Angebots an Hilfen zur Erziehung, \newline Eingliederungshilfen für seelisch \newline behinderte Kinder und Jugendliche und \newline Hilfen für junge Volljährige, & 2. unverändert & 2. unverändert \\
\hline
3. \newline die Anregung und Förderung von \newline Einrichtungen, Diensten und \newline Veranstaltungen sowie deren Schaffung \newline und Betrieb, soweit sie den örtlichen Bedarf \newline übersteigen; dazu gehören insbesondere \newline Einrichtungen, die eine Schul- oder \newline Berufsausbildung anbieten, sowie \newline Jugendbildungsstätten, & 3. unverändert & 3. unverändert \\
\hline
4. \newline die Planung, Anregung, Förderung und \newline Durchführung von Modellvorhaben zur \newline Weiterentwicklung der Jugendhilfe, & 4. unverändert & 4. unverändert \\
\hline
5. \newline die Beratung der örtlichen Träger bei der \newline Gewährung von Hilfe nach den \S{}\S{} 32 bis \newline 35a, insbesondere bei der Auswahl einer \newline Einrichtung oder der Vermittlung einer \newline Pflegeperson in schwierigen Einzelfällen, & 5. \newline die Beratung der örtlichen Träger bei der \newline Gewährung von Hilfe nach den \S{}\S{} 32 bis \newline \textbf{35i}, insbesondere bei der Auswahl einer \newline Einrichtung oder der Vermittlung einer \newline Pflegeperson in schwierigen Einzelfällen, & 5. \newline die Beratung der örtlichen Träger bei der \newline Gewährung von Hilfe nach den \S{}\S{} 32 bis \newline \textbf{35i}, insbesondere bei der Auswahl einer \newline Einrichtung oder der Vermittlung einer \newline Pflegeperson in schwierigen Einzelfällen, \\
\hline
\textit{--- Fassung 2024 ---} \newline 6. \newline die Wahrnehmung der Aufgaben zum \newline Schutz von Kindern und Jugendlichen in \newline Einrichtungen (\S{}\S{} 45 bis 48a), \newline \textit{--- Fassung 2026 ---} \newline 6. & 6. unverändert & 6. unverändert \\
\hline
7. \newline die Beratung der Träger von Einrichtungen \newline während der Planung und Betriebsführung, & 7. unverändert & 7. unverändert \\
\hline
8. \newline die Fortbildung von Mitarbeitern in der \newline Jugendhilfe, & 8. unverändert & 8. unverändert \\
\hline
9. \newline die Gewährung von Leistungen an \newline Deutsche im Ausland (\S{} 6 Absatz 3), \newline soweit es sich nicht um die Fortsetzung \newline einer bereits im Inland gewährten Leistung \newline handelt, & 9. unverändert & 9. unverändert \\
\hline
10. \newline die Anerkennung als \newline Vormundschaftsverein (\S{} 54). & 10. unverändert & 10. unverändert \\
\hline
(3) Für den örtlichen Bereich können die \newline Aufgaben nach Absatz 2 Nummer 3, 4, 7 \newline und 8 auch vom örtlichen Träger \newline wahrgenommen werden. & \textit{(3) unverändert} & \textit{(3) unverändert} \\
\hline
(4) Unberührt bleiben die am Tage des \newline Inkrafttretens dieses Gesetzes geltenden \newline landesrechtlichen Regelungen, die die in \newline den \S{}\S{} 45 bis 48a bestimmten Aufgaben \newline einschließlich der damit verbundenen \newline Aufgaben nach Absatz 2 Nummer 2 bis 5 \newline und 7 mittleren Landesbehörden oder, \newline soweit sie sich auf Kindergärten und \newline andere Tageseinrichtungen für Kinder \newline beziehen, unteren Landesbehörden \newline zuweisen. & \textit{(4) unverändert} & \textit{(4) unverändert} \\
\hline
(5) Ist das Land überörtlicher Träger, so \newline können durch Landesrecht bis zum 30. \newline Juni 1993 einzelne seiner Aufgaben auf \newline andere Körperschaften des öffentlichen \newline Rechts, die nicht Träger der öffentlichen \newline Jugendhilfe sind, übertragen werden. & (5) \textbf{Landesrecht kann bis zum 1.1.2030 \newline bestimmen, dass die Gewährung von \newline Leistungen der Eingliederungshilfe für \newline Kinder und Jugendliche mit \newline Behinderungen im Sinne des \S{} 2 Absatz \newline 2 Nummer 4 Buchstabe b auf den \newline überörtlichen Träger der öffentlichen \newline Jugendhilfe oder auf eine andere \newline Körperschaft des öffentlichen Rechts \newline übertragen wird. Im Falle einer \newline Übertragung nach Satz 1 ist eine \newline ortsnahe Wahrnehmung der Aufgaben \newline nach \S{}\S{} 36 bis 39b unter Einbeziehung \newline des örtlichen Trägers der öffentlichen \newline Jugendhilfe sicherzustellen; \S{} 27 Absatz \newline 5 bleibt unberührt.} & \textbf{(5) Aufgabenübertragungen, die durch \newline ein Land als überörtlichem Träger bis \newline zum 30. Juni 1993 auf andere \newline Körperschaften des öffentlichen Rechts \newline vorgenommen worden sind, bleiben \newline wirksam. } \\
\hline
\multicolumn{3}{|c|}{\cellcolor{gray!10}\textbf{\S{} 86 -- Absätze 1 bis 5}} \\
\hline
(6) Lebt ein Kind oder ein Jugendlicher \newline zwei Jahre bei einer Pflegeperson und ist \newline sein Verbleib bei dieser Pflegeperson auf \newline Dauer zu erwarten, so ist oder wird \newline abweichend von den Absätzen 1 bis 5 der \newline örtliche Träger zuständig, in dessen \newline Bereich die Pflegeperson ihren \newline gewöhnlichen Aufenthalt hat. Er hat die \newline Eltern und, falls den Eltern die \newline Personensorge nicht oder nur teilweise \newline zusteht, den Personensorgeberechtigten \newline über den Wechsel der Zuständigkeit zu \newline unterrichten. Endet der Aufenthalt bei der \newline Pflegeperson, so endet die Zuständigkeit \newline nach Satz 1. & \textit{unverändert} & (6) Lebt ein Kind oder ein Jugendlicher \newline zwei Jahre bei einer Pflegeperson und ist \newline sein Verbleib bei dieser Pflegeperson auf \newline Dauer zu erwarten, so ist oder wird \newline abweichend von den Absätzen 1 bis 5 der \newline örtliche Träger zuständig, in dessen \newline Bereich die Pflegeperson ihren \newline gewöhnlichen Aufenthalt hat. Er hat die \newline Eltern und, falls den Eltern die \newline Personensorge nicht oder nur teilweise \newline zusteht, den Personensorgeberechtigten \newline über den Wechsel der Zuständigkeit zu \newline unterrichten. Endet der Aufenthalt bei der \newline Pflegeperson, so endet die Zuständigkeit \newline nach Satz 1. \textbf{Wurde bei der Auswahl der \newline Pflegeperson der örtliche Träger, in \newline dessen Bereich diese ihren \newline gewöhnlichen Aufenthalt hat, nicht nach \newline \S{} 37 Absatz 3 Satz 5 einbezogen, ist der \newline Wechsel der Zuständigkeit nach Satz 1 \newline ausgeschlossen.} \\
\hline
Absatz 7 & \textit{unverändert} & \textit{unverändert} \\
\hline
\multicolumn{3}{|c|}{\cellcolor{gray!10}\textbf{\S{} 86a -- Örtliche Zuständigkeit für Leistungen an junge Volljährige}} \\
\hline
(1) Für Leistungen an junge Volljährige ist \newline der örtliche Träger zuständig, in dessen \newline Bereich der junge Volljährige vor Beginn \newline der Leistung seinen gewöhnlichen \newline Aufenthalt hat. & \textit{(1) unverändert} & \textit{(1) unverändert} \\
\hline
(2) Hält sich der junge Volljährige in einer \newline Einrichtung oder sonstigen Wohnform auf, \newline die der Erziehung, Pflege, Betreuung, \newline Behandlung oder dem Strafvollzug dient, \newline so richtet sich die örtliche Zuständigkeit \newline nach dem gewöhnlichen Aufenthalt vor der \newline Aufnahme in eine Einrichtung oder \newline sonstige Wohnform. & \textit{(2) unverändert} & \textit{(2) unverändert} \\
\hline
(3) Hat der junge Volljährige keinen \newline gewöhnlichen Aufenthalt, so richtet sich die \newline Zuständigkeit nach seinem tatsächlichen \newline Aufenthalt zu dem in Absatz 1 genannten \newline Zeitpunkt; Absatz 2 bleibt unberührt. & \textit{(3) unverändert} & \textit{(3) unverändert} \\
\hline
(4) Wird eine Leistung nach \S{} 13 Absatz 3 \newline oder nach \S{} 21 über die Vollendung des \newline 18. Lebensjahres hinaus weitergeführt oder \newline geht der Hilfe für junge Volljährige nach \S{} \newline 41 eine dieser Leistungen, eine Leistung \newline nach \S{} 19 oder eine Hilfe nach den \S{}\S{} 27 \newline bis 35a voraus, so bleibt der örtliche Träger \newline zuständig, der bis zu diesem Zeitpunkt \newline zuständig war. Eine Unterbrechung der \newline Hilfeleistung von bis zu drei Monaten bleibt \newline dabei außer Betracht. Die Sätze 1 und 2 \newline gelten entsprechend, wenn eine Hilfe für \newline junge Volljährige nach \S{} 41 beendet war \newline und innerhalb von drei Monaten erneut \newline Hilfe für junge Volljährige nach \S{} 41 \newline erforderlich wird. & (4) Wird eine Leistung nach \S{} 13 Absatz 3 \newline oder nach \S{} 21 über die Vollendung des \newline 18. Lebensjahres hinaus weitergeführt oder \newline geht der Hilfe für junge Volljährige nach \S{} \newline 41 eine dieser Leistungen, eine Leistung \newline nach \S{} 19 oder eine Hilfe nach den \S{}\S{} 27 \newline bis \textbf{35i} voraus, so bleibt der örtliche Träger \newline zuständig, der bis zu diesem Zeitpunkt \newline zuständig war. Eine Unterbrechung der \newline Hilfeleistung von bis zu drei Monaten bleibt \newline dabei außer Betracht. Die Sätze 1 und 2 \newline gelten entsprechend, wenn eine Hilfe für \newline junge Volljährige nach \S{} 41 beendet war \newline und innerhalb von drei Monaten erneut \newline Hilfe für junge Volljährige nach \S{} 41 \newline erforderlich wird. & (4) Wird eine Leistung nach \S{} 13 Absatz 3 \newline oder nach \S{} 21 über die Vollendung des \newline 18. Lebensjahres hinaus weitergeführt oder \newline geht der Hilfe für junge Volljährige nach \S{} \newline 41 eine dieser Leistungen, eine Leistung \newline nach \S{} 19 oder eine Hilfe nach den \S{}\S{} 27 \newline bis \textbf{35i} voraus, so bleibt der örtliche Träger \newline zuständig, der bis zu diesem Zeitpunkt \newline zuständig war. Eine Unterbrechung der \newline Hilfeleistung von bis zu drei Monaten bleibt \newline dabei außer Betracht. Die Sätze 1 und 2 \newline gelten entsprechend, wenn eine Hilfe für \newline junge Volljährige nach \S{} 41 beendet war \newline und innerhalb von drei Monaten erneut \newline Hilfe für junge Volljährige nach \S{} 41 \newline erforderlich wird. \\
\hline
\multicolumn{3}{|c|}{\cellcolor{gray!10}\textbf{\S{} 86b -- Örtliche Zuständigkeit für Leistungen in gemeinsamen Wohnformen für Mütter/Väter und Kinder}} \\
\hline
(1) Für Leistungen in gemeinsamen \newline Wohnformen für Mütter oder Väter und \newline Kinder ist der örtliche Träger zuständig, in \newline dessen Bereich der nach \S{} 19 \newline Leistungsberechtigte vor Beginn der \newline Leistung seinen gewöhnlichen Aufenthalt \newline hat. \S{} 86a Absatz 2 gilt entsprechend. & \textit{(1) unverändert} & \textit{(1) unverändert} \\
\hline
(2) Hat der Leistungsberechtigte keinen \newline gewöhnlichen Aufenthalt, so richtet sich die \newline Zuständigkeit nach seinem tatsächlichen \newline Aufenthalt zu dem in Absatz 1 genannten \newline Zeitpunkt. & \textit{(2) unverändert} & \textit{(2) unverändert} \\
\hline
(3) Geht der Leistung Hilfe nach den \S{}\S{} 27 \newline bis 35a oder eine Leistung nach \S{} 13 \newline Absatz 3, \S{} 21 oder \S{} 41 voraus, so bleibt \newline der örtliche Träger zuständig, der bisher \newline zuständig war. Eine Unterbrechung der \newline Hilfeleistung von bis zu drei Monaten bleibt \newline dabei außer Betracht. & (3) Geht der Leistung Hilfe nach den \S{}\S{} 27 \newline bis \textbf{35i }oder eine Leistung nach \S{} 13 \newline Absatz 3, \S{} 21 oder \S{} 41 voraus, so bleibt \newline der örtliche Träger zuständig, der bisher \newline zuständig war. Eine Unterbrechung der \newline Hilfeleistung von bis zu drei Monaten bleibt \newline dabei außer Betracht. & (3) Geht der Leistung Hilfe nach den \S{}\S{} 27 \newline bis \textbf{35i }oder eine Leistung nach \S{} 13 \newline Absatz 3, \S{} 21 oder \S{} 41 voraus, so bleibt \newline der örtliche Träger zuständig, der bisher \newline zuständig war. Eine Unterbrechung der \newline Hilfeleistung von bis zu drei Monaten bleibt \newline dabei außer Betracht. \\
\hline
\multicolumn{3}{|c|}{\cellcolor{gray!10}\textbf{\S{} 86c -- Fortdauernde Leistungsverpflichtung und Fallübergabe bei Zuständigkeitswechsel}} \\
\hline
(1) Wechselt die örtliche Zuständigkeit für \newline eine Leistung, so bleibt der bisher \newline zuständige örtliche Träger so lange zur \newline Gewährung der Leistung verpflichtet, bis der \newline nunmehr zuständige örtliche Träger die \newline Leistung fortsetzt. Dieser hat dafür Sorge zu \newline tragen, dass der Hilfeprozess und die im \newline Rahmen der Hilfeplanung vereinbarten \newline Hilfeziele durch den Zuständigkeitswechsel \newline nicht gefährdet werden. & (1) Wechselt die örtliche Zuständigkeit für \newline eine Leistung, so bleibt der bisher \newline zuständige örtliche Träger so lange zur \newline Gewährung der Leistung verpflichtet, bis \newline der nunmehr zuständige örtliche Träger die \newline Leistung fortsetzt. Dieser hat dafür Sorge \newline zu tragen, dass der Hilfeprozess und die im \newline Rahmen der \textbf{Hilfe- und Leistungsplanung} \newline vereinbarten Hilfeziele durch den \newline Zuständigkeitswechsel nicht gefährdet \newline werden. & (1) Wechselt die örtliche Zuständigkeit für \newline eine Leistung, so bleibt der bisher \newline zuständige örtliche Träger so lange zur \newline Gewährung der Leistung verpflichtet, bis \newline der nunmehr zuständige örtliche Träger die \newline Leistung fortsetzt. Dieser hat dafür Sorge \newline zu tragen, dass der Hilfeprozess und die im \newline Rahmen der \textbf{Hilfe- und Leistungsplanung} \newline vereinbarten Hilfeziele durch den \newline Zuständigkeitswechsel nicht gefährdet \newline werden. \\
\hline
(2) Der örtliche Träger, der von den \newline Umständen Kenntnis erhält, die den \newline Wechsel der Zuständigkeit begründen, hat \newline den anderen davon unverzüglich zu \newline unterrichten. Der bisher zuständige örtliche \newline Träger hat dem nunmehr zuständigen \newline örtlichen Träger unverzüglich die für die \newline Hilfegewährung sowie den \newline Zuständigkeitswechsel maßgeblichen \newline Sozialdaten zu übermitteln. Bei der \newline Fortsetzung von Leistungen, die der \newline Hilfeplanung nach \S{} 36 Absatz 2 \newline unterliegen, ist die Fallverantwortung im \newline Rahmen eines Gespräches zu übergeben. \newline Die Personensorgeberechtigten und das \newline Kind oder der Jugendliche sowie der junge \newline Volljährige oder der Leistungsberechtigte \newline nach \S{} 19 sind an der Übergabe \newline angemessen zu beteiligen. & (2) Der örtliche Träger, der von den \newline Umständen Kenntnis erhält, die den \newline Wechsel der Zuständigkeit begründen, hat \newline den anderen davon unverzüglich zu \newline unterrichten. Der bisher zuständige örtliche \newline Träger hat dem nunmehr zuständigen \newline örtlichen Träger unverzüglich die für die \newline Hilfegewährung sowie den \newline Zuständigkeitswechsel maßgeblichen \newline Sozialdaten zu übermitteln. Bei der \newline Fortsetzung von Leistungen, die \textbf{der Hilfe- \newline und Leistungsplanung} unterliegen, ist die \newline Fallverantwortung im Rahmen eines \newline Gespräches zu übergeben. Die \newline Personensorgeberechtigten und das Kind \newline oder der Jugendliche sowie der junge \newline Volljährige oder der Leistungsberechtigte \newline nach \S{} 19 sind an der Übergabe \newline angemessen zu beteiligen. & (2) Der örtliche Träger, der von den \newline Umständen Kenntnis erhält, die den \newline Wechsel der Zuständigkeit begründen, hat \newline den anderen davon unverzüglich zu \newline unterrichten. Der bisher zuständige örtliche \newline Träger hat dem nunmehr zuständigen \newline örtlichen Träger unverzüglich die für die \newline Hilfegewährung sowie den \newline Zuständigkeitswechsel maßgeblichen \newline Sozialdaten zu übermitteln. Bei der \newline Fortsetzung von Leistungen, die \textbf{der Hilfe- \newline und Leistungsplanung} unterliegen, ist die \newline Fallverantwortung im Rahmen eines \newline Gespräches zu übergeben. Die \newline Personensorgeberechtigten und das Kind \newline oder der Jugendliche sowie der junge \newline Volljährige oder der Leistungsberechtigte \newline nach \S{} 19 sind an der Übergabe \newline angemessen zu beteiligen. \\
\hline
\multicolumn{3}{|c|}{\cellcolor{gray!10}\textbf{\S{} 87c -- Örtliche Zuständigkeit für die Beistandschaft, die Pflegschaft, die Vormundschaft und die schriftliche Auskunft nach \S{} 58}} \\
\hline
(1) Für die Vormundschaft nach \S{} 1786 des \newline Bürgerlichen Gesetzbuchs ist das \newline Jugendamt zuständig, in dessen Bereich \newline die Mutter ihren gewöhnlichen Aufenthalt \newline hat. Wurde die Vaterschaft nach \S{} 1592 \newline Nummer 1 oder 2 des Bürgerlichen \newline Gesetzbuchs durch Anfechtung beseitigt, \newline so ist der gewöhnliche Aufenthalt der \newline Mutter zu dem Zeitpunkt maßgeblich, zu \newline dem die Entscheidung rechtskräftig wird. \newline Ist ein gewöhnlicher Aufenthalt der Mutter \newline nicht festzustellen, so richtet sich die \newline örtliche Zuständigkeit nach ihrem \newline tatsächlichen Aufenthalt. & \textit{(1) unverändert} & \textit{(1) unverändert} \\
\hline
(2) Sobald die Mutter ihren gewöhnlichen \newline Aufenthalt im Bereich eines anderen \newline Jugendamts nimmt, hat das die \newline Amtsvormundschaft führende Jugendamt \newline bei dem Jugendamt des anderen Bereichs \newline die Weiterführung der Amtsvormundschaft \newline zu beantragen; der Antrag kann auch von \newline dem anderen Jugendamt, von jedem \newline Elternteil und von jedem, der ein \newline berechtigtes Interesse des Kindes oder des \newline Jugendlichen geltend macht, bei dem die \newline Amtsvormundschaft führenden Jugendamt \newline gestellt werden. Die Vormundschaft geht \newline mit der Erklärung des anderen Jugendamts \newline auf dieses über. Das abgebende \newline Jugendamt hat den Übergang dem \newline Familiengericht und jedem Elternteil \newline unverzüglich mitzuteilen. Gegen die \newline Ablehnung des Antrags kann das \newline Familiengericht angerufen werden. & \textit{(2) unverändert} & \textit{(2) unverändert} \\
\hline
(2a) Für die Vormundschaft nach \S{} 1787 \newline des Bürgerlichen Gesetzbuchs ist das \newline Jugendamt zuständig, in dessen Bereich \newline der Geburtsort des Kindes liegt. & \textit{(2a) unverändert} & \textit{(2a) unverändert} \\
\hline
\textit{--- Fassung 2024 ---} \newline (3) Für die Pflegschaft oder \newline Vormundschaft, die durch Bestellung des \newline Familiengerichts eintritt, ist das Jugendamt \newline zuständig, in dessen Bereich das Kind oder \newline der Jugendliche zum Zeitpunkt der \newline Bestellung seinen gewöhnlichen Aufenthalt \newline hat. Hat das Kind oder der Jugendliche \newline keinen gewöhnlichen Aufenthalt, so richtet \newline sich die Zuständigkeit nach seinem \newline tatsächlichen Aufenthalt zum Zeitpunkt der \newline Bestellung. Sobald das Kind oder der \newline Jugendliche seinen gewöhnlichen \newline Aufenthalt nimmt oder wechselt, hat das \newline Jugendamt beim Familiengericht einen \newline Antrag auf Entlassung zu stellen. \newline \textit{--- Fassung 2026 ---} \newline (3) Für die Pflegschaft oder \newline Vormundschaft, die durch Bestellung des & (3) Für die Pflegschaft oder \newline Vormundschaft, die durch Bestellung des \newline Familiengerichts eintritt, ist das Jugendamt \newline zuständig, in dessen Bereich das Kind oder \newline der Jugendliche zum Zeitpunkt der \newline Bestellung seinen gewöhnlichen Aufenthalt \newline hat. Hat das Kind oder der Jugendliche \newline keinen gewöhnlichen Aufenthalt, so richtet \newline sich die Zuständigkeit nach seinem \newline tatsächlichen Aufenthalt zum Zeitpunkt der \newline Bestellung. \textbf{Sobald das Kind oder der \newline Jugendliche seinen gewöhnlichen \newline Aufenthalt nimmt oder wechselt, soll \newline das Jugendamt beim Familiengericht \newline einen Antrag auf Entlassung stellen, \newline wenn es die Voraussetzungen des \S{} \newline 1804 Absatz 3 Satz 1 und 2 des \newline Bürgerlichen Gesetzbuchs für gegeben \newline hält. Lehnt das Familiengericht den \newline Antrag auf Entlassung nach \S{} 1804 \newline Absatz 3 des Bürgerlichen Gesetzbuchs \newline ab, bleibt das zum Vormund oder \newline Pfleger bestellte Jugendamt zuständig. } & (3) Für die Pflegschaft oder \newline Vormundschaft, die durch Bestellung des \\
\hline
Familiengerichts eintritt, ist das Jugendamt \newline zuständig, in dessen Bereich das Kind oder \newline der Jugendliche zum Zeitpunkt der \newline Bestellung seinen gewöhnlichen Aufenthalt \newline hat. Hat das Kind oder der Jugendliche \newline keinen gewöhnlichen Aufenthalt, so richtet \newline sich die Zuständigkeit nach seinem \newline tatsächlichen Aufenthalt zum Zeitpunkt der \newline Bestellung. Sobald das Kind oder der \newline Jugendliche seinen gewöhnlichen \newline Aufenthalt nimmt oder wechselt, hat das \newline Jugendamt beim Familiengericht einen \newline Antrag auf Entlassung zu stellen. & \textit{unverändert} & Familiengerichts eintritt, ist das Jugendamt \newline zuständig, in dessen Bereich das Kind oder \newline der Jugendliche zum Zeitpunkt der \newline Bestellung seinen gewöhnlichen Aufenthalt \newline hat. Hat das Kind oder der Jugendliche \newline keinen gewöhnlichen Aufenthalt, so richtet \newline sich die Zuständigkeit nach seinem \newline tatsächlichen Aufenthalt zum Zeitpunkt der \newline Bestellung. \textbf{Sobald das Kind oder der \newline Jugendliche seinen gewöhnlichen \newline Aufenthalt nimmt oder wechselt, soll \newline das Jugendamt beim Familiengericht \newline einen Antrag auf Entlassung stellen, \newline wenn es die Voraussetzungen des \S{} \newline 1804 Absatz 3 Satz 1 und 2 des \newline Bürgerlichen Gesetzbuchs für gegeben \newline hält. Lehnt das Familiengericht den \newline Antrag auf Entlassung nach \S{} 1804 \newline Absatz 3 des Bürgerlichen Gesetzbuchs \newline ab, bleibt das zum Vormund oder \newline Pfleger bestellte Jugendamt zuständig. } \\
\hline
(4) Für die Vormundschaft, die im Rahmen \newline des Verfahrens zur Annahme als Kind \newline eintritt, ist das Jugendamt zuständig, in \newline dessen Bereich die annehmende Person \newline ihren gewöhnlichen Aufenthalt hat. & \textit{(4) unverändert} & \textit{(4) unverändert} \\
\hline
(5) Für die Beratung und Unterstützung \newline nach \S{} 52a sowie für die Beistandschaft gilt \newline Absatz 1 Satz 1 und 3 entsprechend. \newline Sobald der allein sorgeberechtigte \newline Elternteil seinen gewöhnlichen Aufenthalt \newline im Bereich eines anderen Jugendamts \newline nimmt, hat das die Beistandschaft führende \newline Jugendamt bei dem Jugendamt des \newline anderen Bereichs die Weiterführung der \newline Beistandschaft zu beantragen; Absatz 2 \newline Satz 2 und \S{} 86c gelten entsprechend. & \textit{(5) unverändert} & \textit{(5) unverändert} \\
\hline
(6) Für die Erteilung der schriftlichen \newline Auskunft nach \S{} 58 Absatz 2 gilt Absatz 1 \newline entsprechend. Die Mitteilungen nach \S{} \newline 1626d Absatz 2 des Bürgerlichen \newline Gesetzbuchs, die Mitteilungen nach \S{} 155a \newline Absatz 3 Satz 3 und Absatz 5 Satz 2 des \newline Gesetzes über das Verfahren in \newline Familiensachen und in den \newline Angelegenheiten der freiwilligen \newline Gerichtsbarkeit sowie die Mitteilungen nach \newline \S{} 50 Absatz 3 sind an das für den \newline Geburtsort des Kindes oder des \newline Jugendlichen zuständige Jugendamt zu \newline richten; \S{} 88 Absatz 1 Satz 2 gilt \newline entsprechend. Das nach Satz 2 zuständige \newline Jugendamt teilt dem nach Satz 1 \newline zuständigen Jugendamt auf dessen \newline Ersuchen mit, ob ihm Mitteilungen nach \S{} \newline 1626d Absatz 2 des Bürgerlichen \newline Gesetzbuchs, Mitteilungen nach \S{} 155a \newline Absatz 3 Satz 3 oder Absatz 5 Satz 2 des \newline Gesetzes über das Verfahren in \newline Familiensachen und in den \newline Angelegenheiten der freiwilligen \newline Gerichtsbarkeit oder Mitteilungen nach \S{} \newline 50 Absatz 3 vorliegen. Betrifft die \newline gerichtliche Entscheidung nur Teile der \newline elterlichen Sorge, so enthalten die \newline Mitteilungen auch die Angabe, in welchen \newline Bereichen die elterliche Sorge der Mutter \newline entzogen wurde, den Eltern gemeinsam \newline übertragen wurde oder dem Vater allein \newline übertragen wurde. & \textit{(6) unverändert} & \textit{(6) unverändert} \\
\hline
\multicolumn{3}{|c|}{\cellcolor{gray!10}\textbf{\S{} 89d}} \\
\hline
(1) Kosten, die ein örtlicher Träger \newline aufwendet, sind vom Land zu erstatten, \newline wenn \newline 1. innerhalb eines Monats nach der \newline Einreise eines jungen Menschen oder \newline eines Leistungsberechtigten nach \S{} 19 \newline Jugendhilfe gewährt wird und \newline 2. sich die örtliche Zuständigkeit nach dem \newline tatsächlichen Aufenthalt dieser Person \newline oder nach der Zuweisungsentscheidung \newline der zuständigen Landesbehörde richtet. \newline Als Tag der Einreise gilt der Tag des \newline Grenzübertritts, sofern dieser amtlich \newline festgestellt wurde, oder der Tag, an dem \newline der Aufenthalt im Inland erstmals \newline festgestellt wurde, andernfalls der Tag der \newline ersten Vorsprache bei einem Jugendamt. \newline Die Erstattungspflicht nach Satz 1 bleibt \newline unberührt, wenn die Person um Asyl \newline nachsucht oder einen Asylantrag stellt. & \textit{unverändert} & (1) Kosten, die ein örtlicher Träger \newline aufwendet, sind vom Land zu erstatten, \newline wenn \newline 1. innerhalb eines Monats \textbf{nach Kenntnis \newline des tatsächlichen Aufenthalts eines \newline eingereisten jungen Menschen in \newline dessen Bereich }oder eines \newline Leistungsberechtigten nach \S{} 19 \newline Jugendhilfe gewährt wird und \newline 2. sich die örtliche Zuständigkeit nach dem \newline tatsächlichen Aufenthalt dieser Person \newline oder nach der Zuweisungsentscheidung \newline der zuständigen Landesbehörde richtet. \newline Als Tag der Einreise gilt der Tag des \newline Grenzübertritts, sofern dieser amtlich \newline festgestellt wurde, oder der Tag, an dem \newline der Aufenthalt im Inland erstmals \newline festgestellt wurde, andernfalls der Tag der \newline ersten Vorsprache bei einem Jugendamt. \newline Die Erstattungspflicht nach Satz 1 bleibt \newline unberührt, wenn die Person um Asyl \newline nachsucht oder einen Asylantrag stellt. \\
\hline
Absatz 1 & \textit{unverändert} & \textit{unverändert} \\
\hline
(2) Kosten unter 1 000 Euro werden nur bei \newline vorläufigen Maßnahmen zum Schutz von \newline Kindern und Jugendlichen (\S{} 89b), bei \newline fortdauernder oder vorläufiger \newline Leistungsverpflichtung (\S{} 89c) und bei \newline Gewährung von Jugendhilfe nach der \newline Einreise (\S{} 89d) erstattet. Verzugszinsen \newline können nicht verlangt werden. & \textit{unverändert} & (2) Kosten unter 1 000 Euro werden nur bei \newline vorläufigen Maßnahmen zum Schutz von \newline Kindern und Jugendlichen (\S{} 89b) \textbf{und} bei \newline fortdauernder oder vorläufiger \newline Leistungsverpflichtung (\S{} 89c) und bei \newline Gewährung von Jugendhilfe nach der \newline Einreise (\S{} 89d) erstattet. \textbf{Bei Gewährung \newline von Jugendhilfe an eingereiste junge \newline Menschen (\S{} 89d) kann Landesrecht \newline regeln, dass Kosten unter 1000 Euro \newline erstattet werden. Landesrecht kann eine \newline pauschale Abgeltung aufgewendeter \newline Kosten vorsehen, soweit dies \newline zweckmäßig ist.} Verzugszinsen können \newline nicht verlangt werden. \\
\hline
\multicolumn{3}{|c|}{\cellcolor{gray!10}\textbf{\S{} 91 -- Anwendungsbereich}} \\
\hline
(1) Zu folgenden vollstationären Leistungen \newline und vorläufigen Maßnahmen werden \newline Kostenbeiträge erhoben: & \textit{(1) unverändert} & \textit{(1) unverändert} \\
\hline
1. \newline der Unterkunft junger Menschen in einer \newline sozialpädagogisch begleiteten Wohnform \newline (\S{} 13 Absatz 3), & 1. unverändert & 1. unverändert \\
\hline
2. \newline der Betreuung von Müttern oder Vätern \newline und Kindern in gemeinsamen Wohnformen \newline (\S{} 19), & 2. unverändert & 2. unverändert \\
\hline
3. \newline der Betreuung und Versorgung von \newline Kindern in Notsituationen (\S{} 20), & 3. unverändert & 3. unverändert \\
\hline
4. \newline der Unterstützung bei notwendiger \newline Unterbringung junger Menschen zur \newline Erfüllung der Schulpflicht und zum \newline Abschluss der Schulausbildung (\S{} 21), & 4. unverändert & 4. unverändert \\
\hline
5. \newline der Hilfe zur Erziehung & 5. unverändert & 5. unverändert \\
\hline
d) \newline auf der Grundlage von \S{} 27 in stationärer \newline Form, & d) unverändert & d) unverändert \\
\hline
6. \newline der Eingliederungshilfe für seelisch \newline behinderte Kinder und Jugendliche durch \newline geeignete Pflegepersonen sowie in \newline Einrichtungen über Tag und Nacht und in \newline sonstigen Wohnformen (\S{} 35a Absatz 2 \newline Nummer 3 und 4), & 6. \newline der Eingliederungshilfe für Kinder und \newline Jugendliche \textbf{mit Behinderungen} durch \newline geeignete Pflegepersonen sowie in \newline Einrichtungen über Tag und Nacht und in \newline sonstigen Wohnformen (\S{} 35a Absatz \textbf{4} \newline Nummer 3 und 4), & 6. \newline der Eingliederungshilfe für Kinder und \newline Jugendliche \textbf{mit Behinderungen} durch \newline geeignete Pflegepersonen sowie in \newline Einrichtungen über Tag und Nacht und in \newline sonstigen Wohnformen (\S{} 35a Absatz \textbf{4} \newline Nummer 3 und 4), \\
\hline
7. \newline der Inobhutnahme von Kindern und \newline Jugendlichen (\S{} 42), & 7. unverändert & 7. unverändert \\
\hline
8. \newline der Hilfe für junge Volljährige, soweit sie \newline den in den Nummern 5 und 6 genannten \newline Leistungen entspricht (\S{} 41). & 8. unverändert & 8. unverändert \\
\hline
(2) Zu folgenden teilstationären Leistungen \newline werden Kostenbeiträge erhoben: & \textit{(2) unverändert} & \textit{(2) unverändert} \\
\hline
1. \newline der Betreuung und Versorgung von \newline Kindern in Notsituationen nach \S{} 20, & 1. unverändert & 1. unverändert \\
\hline
2. \newline Hilfe zur Erziehung in einer Tagesgruppe \newline nach \S{} 32 und anderen teilstationären \newline Leistungen nach \S{} 27, & 2. unverändert & 2. unverändert \\
\hline
3. \newline Eingliederungshilfe für seelisch behinderte \newline Kinder und Jugendliche in \newline Tageseinrichtungen und anderen \newline teilstationären Einrichtungen nach \S{} 35a \newline Absatz 2 Nummer 2 und & 3. \newline Eingliederungshilfe \textbf{für Kinder und \newline Jugendliche mit Behinderungen} in \newline Tageseinrichtungen und anderen \newline teilstationären Einrichtungen nach \S{} 35a \newline Absatz 2 Nummer 2 und & 3. \newline Eingliederungshilfe \textbf{für Kinder und \newline Jugendliche mit Behinderungen} in \newline Tageseinrichtungen und anderen \newline teilstationären Einrichtungen nach \S{} 35a \newline Absatz \textbf{4} Nummer 2 und \\
\hline
4. \newline Hilfe für junge Volljährige, soweit sie den in \newline den Nummern 2 und 3 genannten \newline Leistungen entspricht (\S{} 41). & 4. unverändert & 4. unverändert \\
\hline
(3) Die Kosten umfassen auch die \newline Aufwendungen für den notwendigen \newline Unterhalt und die Krankenhilfe. & \textbf{(3) Ausgenommen von der \newline Kostenbeitragspflicht nach den \newline Absätzen 1 und 2 sind Leistungen zur \newline Beschäftigung sowie Leistungen zum \newline Erwerb und Erhalt praktischer \newline Kenntnisse und Fähigkeiten nach \S{} 35f \newline Absatz 2 Nummer 5, sowie diese der \newline Vorbereitung auf die Teilhabe am \newline Arbeitsleben nach \S{} 35e Absatz 1 \newline dienen. } & \textbf{(3) Ausgenommen von der \newline Kostenbeitragspflicht nach den \newline Absätzen 1 und 2 sind Leistungen zur \newline Teilhabe am Arbeitsleben sowie \newline Leistungen zum Erwerb und Erhalt \newline praktischer Kenntnisse und Fähigkeiten \newline nach \S{} 35f Absatz 2 Nummer 5, soweit \newline diese der Vorbereitung auf die Teilhabe \newline am Arbeitsleben nach \S{} 35e Absatz 1 \newline dienen. } \\
\hline
(4) Verwaltungskosten bleiben außer \newline Betracht. & \textbf{(4) Unter den weiteren Leistungen \newline werden Kostenbeiträge zu Leistungen \newline zur Mobilität und Leistungen für \newline Wohnraum (\S{} 35f Absatz 2 Nummern 1 \newline und 7) erhoben. } & \textbf{(4) Neben den kostenbeitragspflichtigen \newline Leistungen nach Absatz 1 und 2 werden \newline Kostenbeiträge zu Leistungen zur \newline Mobilität und Leistungen für Wohnraum \newline (\S{} 35f Absatz 2 Nummer 1 und 7) \newline erhoben. } \\
\hline
(5) Die Träger der öffentlichen Jugendhilfe \newline tragen die Kosten der in den Absätzen 1 \newline und 2 genannten Leistungen unabhängig \newline von der Erhebung eines Kostenbeitrags. & (5) Die Kosten umfassen auch die \newline Aufwendungen für den notwendigen \newline Unterhalt und die Krankenhilfe. & (5) Die Kosten umfassen auch die \newline Aufwendungen für den notwendigen \newline Unterhalt und die Krankenhilfe. \\
\hline
 & (6) Verwaltungskosten bleiben außer \newline Betracht. & (6) Verwaltungskosten bleiben außer \newline Betracht. \\
\hline
 & (7) Die Träger der öffentlichen Jugendhilfe \newline tragen die Kosten der in den Absätzen 1 \newline und 2 genannten Leistungen unabhängig \newline von der Erhebung eines Kostenbeitrags. & (7) Die Träger der öffentlichen Jugendhilfe \newline tragen die Kosten der in den Absätzen 1 \newline und 2 genannten Leistungen unabhängig \newline von der Erhebung eines Kostenbeitrags. \\
\hline
\multicolumn{3}{|c|}{\cellcolor{gray!10}\textbf{\S{} 92 -- Ausgestaltung der Heranziehung}} \\
\hline
\textit{--- Fassung 2024 ---} \newline (1) Zu den Kosten der in \S{} 91 Absatz 1 \newline genannten Leistungen und vorläufigen \newline Maßnahmen sind Elternteile aus ihrem \newline Einkommen nach Maßgabe der \S{}\S{} 93 und \newline 94 heranzuziehen; leben sie mit dem \newline jungen Menschen zusammen, so werden \newline sie auch zu den Kosten der in \S{} 91 Absatz \newline 2 genannten Leistungen herangezogen. \newline \textit{--- Fassung 2026 ---} \newline (1) Zu den Kosten der in \S{} 91 Absatz 1 \newline genannten Leistungen und vorläufigen \newline Maßnahmen sind Elternteile aus ihrem & (1) Zu den Kosten der in \S{} 91 Absatz 1 \textbf{und \newline 4} genannten Leistungen und vorläufigen \newline Maßnahmen sind Elternteile aus ihrem \newline Einkommen nach Maßgabe der \S{}\S{} 93 und \newline 94 heranzuziehen; leben sie mit dem \newline jungen Menschen zusammen, so werden \newline sie auch zu den Kosten der in \S{} 91 Absatz \newline 2 genannten Leistungen herangezogen. & (1) Zu den Kosten der in \S{} 91 Absatz 1 \textbf{und \newline 4} genannten Leistungen und vorläufigen \newline Maßnahmen sind Elternteile aus ihrem \\
\hline
Einkommen nach Maßgabe der \S{}\S{} 93 und \newline 94 heranzuziehen; leben sie mit dem \newline jungen Menschen zusammen, so werden \newline sie auch zu den Kosten der in \S{} 91 Absatz \newline 2 genannten Leistungen herangezogen. & \textit{unverändert} & Einkommen nach Maßgabe der \S{}\S{} 93 und \newline 94 heranzuziehen; leben sie mit dem \newline jungen Menschen zusammen, so werden \newline sie auch zu den Kosten der in \S{} 91 Absatz \newline 2 genannten Leistungen herangezogen. \\
\hline
(1a) Unabhängig von ihrem Einkommen \newline sind nach Maßgabe von \S{} 93 Absatz 1 \newline Satz 3 und \S{} 94 Absatz 3 heranzuziehen: & (1a) Unabhängig von ihrem Einkommen \newline sind nach Maßgabe von \S{} 93 Absatz 1 \newline Satz 3 und \S{} 94 Absatz 3 \textbf{aus ihren \newline Einnahmen} heranzuziehen: & (1a) Unabhängig von ihrem Einkommen \newline sind nach Maßgabe von \S{} 93 Absatz 1 \newline Satz 3 und \S{} 94 Absatz 3 \textbf{aus ihren \newline Einnahmen} heranzuziehen: \\
\hline
1. \newline Kinder und Jugendliche zu den Kosten der \newline in \S{} 91 Absatz 1 Nummer 1 bis 7 \newline genannten Leistungen und vorläufigen \newline Maßnahmen, & 1. \newline Kinder und Jugendliche zu den Kosten der \newline in \S{} 91 Absatz 1 Nummer 1 bis 7 \textbf{und \newline Absatz 4} genannten Leistungen und \newline vorläufigen Maßnahmen, & 1. \newline Kinder und Jugendliche zu den Kosten der \newline in \S{} 91 Absatz 1 Nummer 1 bis 7 \textbf{und \newline Absatz 4} genannten Leistungen und \newline vorläufigen Maßnahmen, \\
\hline
2. \newline junge Volljährige zu den Kosten der in \S{} 91 \newline Absatz 1 Nummer 1, 4 und 8 genannten \newline Leistungen, & 2. \newline junge Volljährige zu den Kosten der in \S{} 91 \newline Absatz 1 Nummer 1, 4 und 8 \textbf{und Absatz 4} \newline genannten Leistungen, & 2. \newline junge Volljährige zu den Kosten der in \S{} 91 \newline Absatz 1 Nummer 1, 4 und 8 \textbf{und Absatz 4} \newline genannten Leistungen, \\
\hline
3. \newline Leistungsberechtigte nach \S{} 19 zu den \newline Kosten der in \S{} 91 Absatz 1 Nummer 2 \newline genannten Leistungen, & 3. unverändert & 3. unverändert \\
\hline
4. \newline Elternteile zu den Kosten der in \S{} 91 \newline Absatz 1 genannten Leistungen und \newline vorläufigen Maßnahmen; leben sie mit dem \newline jungen Menschen zusammen, so werden \newline sie auch zu den Kosten der in \S{} 91 Absatz \newline 2 genannten Leistungen herangezogen. & 4. \newline Elternteile zu den Kosten der in \S{} 91 \newline Absatz 1 \textbf{und Absatz 4} genannten \newline Leistungen und vorläufigen Maßnahmen; \newline leben sie mit dem jungen Menschen \newline zusammen, so werden sie auch zu den \newline Kosten der in \S{} 91 Absatz 2 genannten \newline Leistungen herangezogen. & 4. \newline Elternteile zu den Kosten der in \S{} 91 \newline Absatz 1 \textbf{und Absatz 4} genannten \newline Leistungen und vorläufigen Maßnahmen; \newline leben sie mit dem jungen Menschen \newline zusammen, so werden sie auch zu den \newline Kosten der in \S{} 91 Absatz 2 genannten \newline Leistungen herangezogen. \\
\hline
(2) Die Heranziehung erfolgt durch \newline Erhebung eines Kostenbeitrags, der durch \newline Leistungsbescheid festgesetzt wird; \newline Elternteile werden getrennt herangezogen. & (2) Die Heranziehung erfolgt durch \newline Erhebung eines Kostenbeitrags, der durch \newline Leistungsbescheid festgesetzt wird; \textbf{Eltern} \newline werden \textbf{gemeinsam} herangezogen, \textbf{wenn \newline sie in einem gemeinsamen Haushalt \newline leben, in dem auch der junge Mensch \newline üblicherweise lebt, wenn er nicht im \newline Rahmen einer Leistung über Tag und \newline Nacht in einer betreuten Wohnform oder \newline bei einer Pflegeperson untergebracht \newline ist. Ist dies nicht der Fall, erfolgt eine \newline getrennte Heranziehung.} & (2) Die Heranziehung erfolgt durch \newline Erhebung eines Kostenbeitrags, der durch \newline Leistungsbescheid festgesetzt wird. \\
\hline
(3) Ein Kostenbeitrag kann bei Eltern ab \newline dem Zeitpunkt erhoben werden, ab \newline welchem dem Pflichtigen die Gewährung \newline der Leistung mitgeteilt und er über die \newline Folgen für seine Unterhaltspflicht \newline gegenüber dem jungen Menschen \newline aufgeklärt wurde. Ohne vorherige \newline Mitteilung kann ein Kostenbeitrag für den \newline Zeitraum erhoben werden, in welchem der \newline Träger der öffentlichen Jugendhilfe aus \newline rechtlichen oder tatsächlichen Gründen, die \newline in den Verantwortungsbereich des \newline Pflichtigen fallen, an der Geltendmachung \newline gehindert war. Entfallen diese Gründe, ist \newline der Pflichtige unverzüglich zu unterrichten. & \textit{(3) unverändert} & \textit{(3) unverändert} \\
\hline
(4) Ein Kostenbeitrag kann nur erhoben \newline werden, soweit Unterhaltsansprüche \newline vorrangig oder gleichrangig Berechtigter \newline nicht geschmälert werden. Von der \newline Heranziehung der Eltern ist abzusehen, \newline wenn das Kind, die Jugendliche, die junge \newline Volljährige oder die Leistungsberechtigte \newline nach \S{} 19 schwanger ist oder der junge \newline Mensch oder die nach \S{} 19 \newline leistungsberechtigte Person ein leibliches \newline Kind bis zur Vollendung des sechsten \newline Lebensjahres betreut. & \textit{(4) unverändert} & \textit{(4) unverändert} \\
\hline
(5) Von der Heranziehung soll im Einzelfall \newline ganz oder teilweise abgesehen werden, \newline wenn sonst Ziel und Zweck der Leistung \newline gefährdet würden oder sich aus der \newline Heranziehung eine besondere Härte \newline ergäbe. Von der Heranziehung kann \newline abgesehen werden, wenn anzunehmen ist, \newline dass der damit verbundene \newline Verwaltungsaufwand in keinem \newline angemessenen Verhältnis zu dem \newline Kostenbeitrag stehen wird. & \textit{(5) unverändert} & \textit{(5) unverändert} \\
\hline
\multicolumn{3}{|c|}{\cellcolor{gray!10}\textbf{\S{} 92a -- Pauschale Heranziehung}} \\
\hline
 & \textit{unverändert} & \textbf{(1) Zur Heranziehung zu den Kosten der \newline Hilfen oder Leistungen nach \S{} 91 Absatz \newline 1, bei denen Leistungen zum Unterhalt \newline nach \S{} 39 umfassend gewährt werden, \newline wird bei Elternteilen jeweils ein \newline pauschaler Kostenbeitrag in Höhe von \newline 50 Prozent der Regelbedarfsstufen im \newline Sinne der Anlage des \S{} 28 des Zwölften \newline Buches erhoben. } \\
\hline
 & \textit{unverändert} & \textbf{(2) Zur Heranziehung zu den Kosten der \newline anderen Hilfen oder Leistungen nach \S{} \newline 91 Absatz 1, bei denen keine Leistungen \newline zum Unterhalt nach \S{} 39 gewährt \newline werden, sowie zu den Kosten der Hilfen \newline oder Leistungen nach \S{} 91 Absatz 2 } \newline \textbf{wird bei den Elternteilen jeweils ein \newline pauschaler Kostenbeitrag in Höhe von \newline 20 Prozent der Regelbedarfsstufen im \newline Sinne der Anlage des \S{} 28 des Zwölften \newline Buches erhoben. } \\
\hline
 & \textit{unverändert} & \textbf{(3) Zur Heranziehung zu den Kosten der \newline Leistungen nach \S{} 91 Absatz 4 wird bei \newline den Elternteilen jeweils ein pauschaler \newline Kostenbeitrag in Höhe von 10 Prozent \newline der Regelbedarfsstufen im Sinne der \newline Anlage des \S{} 28 des Zwölften Buches \newline erhoben. } \\
\hline
\multicolumn{3}{|c|}{\cellcolor{gray!10}\textbf{\S{} 93 -- Berechnung des Einkommens}} \\
\hline
(1) Zum Einkommen gehören alle Einkünfte \newline in Geld oder Geldeswert mit Ausnahme der \newline Leistungen nach diesem Buch, der \newline Leistungen nach dem Vierzehnten Buch \newline und der Leistungen nach Gesetzen, die \newline eine entsprechende Anwendung des \newline Vierzehnten Buches vorsehen, und der \newline Renten oder Beihilfen nach dem \newline Bundesentschädigungsgesetz für Schaden \newline an Leben sowie an Körper oder \newline Gesundheit bis zur Höhe der \newline vergleichbaren Leistungen nach dem \newline Vierzehnten Buch. Eine Entschädigung, die \newline nach \S{} 253 Absatz 2 des Bürgerlichen \newline Gesetzbuchs wegen eines Schadens, der \newline nicht Vermögensschaden ist, geleistet wird, \newline ist nicht als Einkommen zu \newline berücksichtigen. Geldleistungen, die dem \newline gleichen Zwecke wie die jeweilige Leistung \newline der Jugendhilfe dienen, zählen nicht zum \newline Einkommen und sind unabhängig von \newline einem Kostenbeitrag einzusetzen; dies gilt \newline nicht für & \textit{(1) unverändert} & \textit{(1) unverändert} \\
\hline
1. \newline monatliche Leistungen nach \S{} 56 des \newline Dritten Buches bis zu einer Höhe des in \S{} \newline 61 Absatz 2 Satz 1 und \S{} 62 Absatz 3 Satz \newline 1 des Dritten Buches für sonstige \newline Bedürfnisse genannten Betrages und & 1. unverändert & 1. unverändert \\
\hline
\textit{--- Fassung 2024 ---} \newline 2. \newline monatliche Leistungen nach \S{} 122 des \newline Dritten Buches bis zu einer Höhe des in \S{} \newline 123 Satz 1 Nummer 2, \S{} 124 Nummer 2 \newline und \S{} 125 des Dritten Buches genannten \newline Betrages. Kindergeld und Leistungen, die \newline auf Grund öffentlich-rechtlicher Vorschriften \newline zu einem ausdrücklich genannten Zweck \newline erbracht werden, sind nicht als Einkommen \newline zu berücksichtigen. \newline \textit{--- Fassung 2026 ---} \newline 2. & 2. \newline monatliche Leistungen nach \S{} 122 des \newline Dritten Buches bis zu einer Höhe des in \S{} \newline 123 Satz 1 Nummer 2, \S{} 124 Nummer 2 \newline und \S{} 125 des Dritten Buches genannten \newline Betrages. \textbf{Der Einsatz von \newline Geldleistungen im Sinne des Satz 3 ist \newline auf den in der Anlage zur Verordnung im \newline Sinne des \S{} 94 Absatz 5 genannten \newline Höchstbetrag begrenzt.} Leistungen, die \newline auf Grund öffentlich-rechtlicher Vorschriften \newline zu einem ausdrücklich genannten Zweck \newline erbracht werden, sind nicht als Einkommen \newline zu berücksichtigen. & 2. \\
\hline
monatliche Leistungen nach \S{} 122 des \newline Dritten Buches bis zu einer Höhe des in \S{} \newline 123 Satz 1 Nummer 2, \S{} 124 Nummer 2 \newline und \S{} 125 des Dritten Buches genannten \newline Betrages. Kindergeld und Leistungen, die \newline auf Grund öffentlich-rechtlicher Vorschriften \newline zu einem ausdrücklich genannten Zweck \newline erbracht werden, sind nicht als Einkommen \newline zu berücksichtigen. & \textit{unverändert} & monatliche Leistungen nach \S{} 122 des \newline Dritten Buches bis zu einer Höhe des in \S{} \newline 123 Satz 1 Nummer 2, \S{} 124 Nummer 2 \newline und \S{} 125 des Dritten Buches genannten \newline Betrages. \textbf{Der Einsatz von \newline Geldleistungen im Sinne des Satz 3 ist \newline auf den in der Anlage zur Verordnung im \newline Sinne des \S{} 94 Absatz 5 genannten \newline Höchstbetrag begrenzt.} Leistungen, die \newline auf Grund öffentlich-rechtlicher Vorschriften \newline zu einem ausdrücklich genannten Zweck \newline erbracht werden, sind nicht als Einkommen \newline zu berücksichtigen. \\
\hline
(2) Von dem Einkommen sind abzusetzen & \textit{(2) unverändert} & \textit{(2) unverändert} \\
\hline
1. \newline auf das Einkommen gezahlte Steuern und & 1. unverändert & 1. unverändert \\
\hline
2. \newline Pflichtbeiträge zur Sozialversicherung \newline einschließlich der Beiträge zur \newline Arbeitsförderung sowie & 2. unverändert & 2. unverändert \\
\hline
3. \newline nach Grund und Höhe angemessene \newline Beiträge zu öffentlichen oder privaten \newline Versicherungen oder ähnlichen \newline Einrichtungen zur Absicherung der Risiken \newline Alter, Krankheit, Pflegebedürftigkeit und \newline Arbeitslosigkeit. & 3. unverändert & 3. unverändert \\
\hline
(3) Von dem nach den Absätzen 1 und 2 \newline errechneten Betrag sind Belastungen der \newline kostenbeitragspflichtigen Person \newline abzuziehen. Der Abzug erfolgt durch eine \newline Kürzung des nach den Absätzen 1 und 2 \newline errechneten Betrages um pauschal 25 vom \newline Hundert. Sind die Belastungen höher als \newline der pauschale Abzug, so können sie \newline abgezogen werden, soweit sie nach Grund \newline und Höhe angemessen sind und die \newline Grundsätze einer wirtschaftlichen \newline Lebensführung nicht verletzen. In Betracht \newline kommen insbesondere & \textit{(3) unverändert} & \textit{(3) unverändert} \\
\hline
1. \newline Beiträge zu öffentlichen oder privaten \newline Versicherungen oder ähnlichen \newline Einrichtungen, & 1. unverändert & 1. unverändert \\
\hline
2. \newline die mit der Erzielung des Einkommens \newline verbundenen notwendigen Ausgaben, & 2. unverändert & 2. unverändert \\
\hline
3. \newline Schuldverpflichtungen. Die \newline kostenbeitragspflichtige Person muss die \newline Belastungen nachweisen. & 3. unverändert & 3. unverändert \\
\hline
(4) Maßgeblich ist das durchschnittliche \newline Monatseinkommen, das die \newline kostenbeitragspflichtige Person in dem \newline Kalenderjahr erzielt hat, welches dem \newline jeweiligen Kalenderjahr der Leistung oder \newline Maßnahme vorangeht. Auf Antrag der \newline kostenbeitragspflichtigen Person wird \newline dieses Einkommen nachträglich durch das \newline durchschnittliche Monatseinkommen \newline ersetzt, welches die Person in dem \newline jeweiligen Kalenderjahr der Leistung oder \newline Maßnahme erzielt hat. Der Antrag kann \newline innerhalb eines Jahres nach Ablauf dieses \newline Kalenderjahres gestellt werden. Macht die \newline kostenbeitragspflichtige Person glaubhaft, \newline dass die Heranziehung zu den Kosten aus \newline dem Einkommen nach Satz 1 in einem \newline bestimmten Zeitraum eine besondere Härte \newline für sie ergäbe, wird vorläufig von den \newline glaubhaft gemachten, dem Zeitraum \newline entsprechenden Monatseinkommen \newline ausgegangen; endgültig ist in diesem Fall \newline das nach Ablauf des Kalenderjahres zu \newline ermittelnde durchschnittliche \newline Monatseinkommen dieses Jahres \newline maßgeblich. & \textit{(4) unverändert} & \textit{(4) unverändert} \\
\hline
 & \textbf{(5) Kindergeld, das für den jungen \newline Menschen, der die Leistung erhält, \newline geleistet wird, wird dem maßgeblichen \newline Einkommen im Sinne des Absatz 3 des \newline Elternteils hinzugerechnet, der das \newline Kindergeld erhält. Erhält der junge \newline Mensch selbst das Kindergeld, so gilt es \newline unabhängig vom Einkommen als \newline Einnahme des jungen Menschen. } & \textbf{(5) Kindergeld, das für den jungen \newline Menschen, der die Leistung erhält, \newline geleistet wird, wird dem maßgeblichen \newline Einkommen im Sinne des Absatzes 3 \newline des Elternteils, der das Kindergeld \newline erhält, hinzugerechnet. Erhält der junge \newline Mensch für sich selbst das Kindergeld \newline nach \S{} 1 Absatz 2 des \newline Bundeskindergeldgesetzes oder durch \newline Abzweigung nach \S{} 74 Absatz 1 des \newline Einkommensteuergesetzes, so gilt \S{} 94 \newline Absatz 3. } \\
\hline
\multicolumn{3}{|c|}{\cellcolor{gray!10}\textbf{\S{} 94 -- Umfang der Heranziehung}} \\
\hline
(1) Die Kostenbeitragspflichtigen sind aus \newline ihrem Einkommen in angemessenem \newline Umfang zu den Kosten heranzuziehen. Die \newline Kostenbeiträge dürfen die tatsächlichen \newline Aufwendungen nicht überschreiten. & \textit{(1) unverändert} & \textit{(1) unverändert} \\
\hline
(2) Für die Bestimmung des Umfangs sind \newline bei jedem Elternteil die Höhe des nach \S{} \newline 93 ermittelten Einkommens und die Anzahl \newline der Personen, die mindestens im gleichen \newline Range wie der untergebrachte junge \newline Mensch oder Leistungsberechtigte nach \S{} \newline 19 unterhaltsberechtigt sind, angemessen \newline zu berücksichtigen. & \textit{(2) unverändert} & \textit{(2) unverändert} \\
\hline
(3) Werden Leistungen über Tag und Nacht \newline außerhalb des Elternhauses erbracht und \newline bezieht einer der Elternteile Kindergeld für \newline den jungen Menschen, so hat dieser \newline unabhängig von einer Heranziehung nach \newline Absatz 1 Satz 1 und 2 einen Kostenbeitrag \newline in Höhe des Kindergeldes zu zahlen. Zahlt \newline der Elternteil den Kostenbeitrag nach Satz \newline 1 nicht, so sind die Träger der öffentlichen \newline Jugendhilfe insoweit berechtigt, das auf \newline dieses Kind entfallende Kindergeld durch \newline Geltendmachung eines \newline Erstattungsanspruchs nach \S{} 74 Absatz 2 \newline des Einkommensteuergesetzes in \newline Anspruch zu nehmen. Bezieht der Elternteil \newline Kindergeld nach \S{} 1 Absatz 1 des \newline Bundeskindergeldgesetzes, gilt Satz 2 \newline entsprechend. Bezieht der junge Mensch \newline das Kindergeld selbst, gelten die Sätze 1 \newline und 2 entsprechend. Die Heranziehung der \newline Elternteile erfolgt nachrangig zu der \newline Heranziehung der jungen Menschen zu \newline einem Kostenbeitrag in Höhe des \newline Kindergeldes. & (3) Werden Leistungen über Tag und Nacht \newline außerhalb des Elternhauses erbracht und \newline \textbf{erhält der junge Mensch das Kindergeld \newline selbst, hat er einen Kostenbeitrag in \newline Höhe des Kindegeldes als \newline Kostenbeitrag zu zahlen.} Zahlt der \textbf{junge \newline Mensch} den Kostenbeitrag nach Satz 1 \newline nicht, so sind die Träger der öffentlichen \newline Jugendhilfe insoweit berechtigt, das auf \newline dieses Kind entfallende Kindergeld durch \newline Geltendmachung eines \newline Erstattungsanspruchs nach \S{} 74 Absatz 2 \newline des Einkommensteuergesetzes in \newline Anspruch zu nehmen. & \textbf{(3) Werden Leistungen über Tag und \newline Nacht außerhalb des Elternhauses \newline erbracht und erhält der junge Mensch \newline das Kindergeld für sich selbst nach \S{} 1 \newline Absatz 2 des \newline Bundeskindergeldgesetzes oder durch \newline Abzweigung nach \S{} 74 Absatz 1 des \newline Einkommensteuergesetzes, hat er einen \newline Kostenbeitrag in Höhe des Kindergeldes \newline zu zahlen. Zahlt der junge Mensch den \newline Kostenbeitrag nach Satz 1 nicht, so sind \newline die Träger der öffentlichen Jugendhilfe \newline insoweit berechtigt, das nach Satz 1 auf \newline diesen jungen Menschen entfallende \newline Kindergeld durch Geltendmachung \newline eines Erstattungsanspruchs nach \S{} 74 \newline Absatz 2 des Einkommensteuergesetzes \newline in Anspruch zu nehmen. Der Einsatz \newline von Geldleistungen nach \S{} 93 Absatz 1 \newline Satz 3 geht der Heranziehung nach Satz \newline 1 vor. Kommt sowohl eine Heranziehung \newline nach Satz 1 als auch nach \S{} 93 Absatz 1 \newline Satz 3 in Betracht, darf die Summe der \newline Heranziehung den Höchstbetrag nach \S{} \newline 93 Absatz 1 Satz 4 nicht überschreiten. } \\
\hline
(4) Werden Leistungen über Tag und Nacht \newline erbracht und hält sich der junge Mensch \newline nicht nur im Rahmen von \newline Umgangskontakten bei einem \newline Kostenbeitragspflichtigen auf, so ist die \newline tatsächliche Betreuungsleistung über Tag \newline und Nacht auf den Kostenbeitrag \newline anzurechnen. & \textit{(4) unverändert} & \textit{(4) unverändert} \\
\hline
(5) Für die Festsetzung der Kostenbeiträge \newline von Eltern werden nach \newline Einkommensgruppen gestaffelte \newline Pauschalbeträge durch Rechtsverordnung \newline des zuständigen Bundesministeriums mit \newline Zustimmung des Bundesrates bestimmt. & (5) Für die Festsetzung der Kostenbeiträge \newline von Eltern \textbf{für Leistungen nach \S{} 91 \newline Absatz 1 und 2} werden nach \newline Einkommensgruppen gestaffelte \newline Pauschalbeträge durch Rechtsverordnung \newline des zuständigen Bundesministeriums mit \newline Zustimmung des Bundesrates bestimmt. \newline \textbf{Die Pauschalbeträge orientieren sich an \newline der Höhe der Regelbedarfsstufen im \newline Sinne der Anlage des \S{} 28 Zwölften \newline Buches in einer Spanne von 0 bis 150 \newline Prozent. Bei getrennt herangezogenen \newline Elternteilen darf die Summe beider \newline Kostenbeiträge den Höchstbetrag der \newline Kostenbeiträge aus der Anlage zur \newline Rechtsverordnung nicht überschreiten. \newline Bei Leistungen nach \S{} 41 steht die \newline Heranziehung unter der Bedingung, \newline dass Kindergeld für den jungen \newline Menschen bezogen wird; die \newline Heranziehung ist auf die Höhe des \newline Kindergeldes begrenzt. Für die \newline Festsetzung der Kostenbeiträge von \newline Eltern für Leistungen nach \S{} 91 Absatz 3 \newline werden in der Rechtsverordnung die \newline Anteile der Kostenbeteiligung an den \newline Kosten der Leistung bestimmt.} & (5) Für die Festsetzung der Kostenbeiträge \newline von Eltern \textbf{für Leistungen nach \S{} 91 \newline Absatz 1 und 2} werden nach \newline Einkommensgruppen gestaffelte \newline Pauschalbeträge durch Rechtsverordnung \newline des zuständigen Bundesministeriums mit \newline Zustimmung des Bundesrates bestimmt. \newline \textbf{Für Hilfen oder Leistungen nach \S{} 91 \newline Absatz 1, bei denen Leistungen zum \newline Unterhalt nach \S{} 39 umfassend gewährt \newline werden, orientieren sich die \newline Pauschalbeträge an den \newline Regelbedarfsstufen im Sinne der Anlage \newline des \S{} 28 des Zwölften Buches in einer \newline Spanne von 0 bis 100 Prozent; für \newline andere Hilfen oder Leistungen nach \S{} 91 \newline Absatz 1 sowie Hilfen oder Leistungen \newline nach \S{} 91 Absatz 2 orientieren sich die \newline Pauschalbeträge an den für den \newline häuslichen Lebensunterhalt vermuteten \newline ersparten Aufwendungen. Werden beide \newline Elternteile zu den Kosten herangezogen, \newline darf die Summe beider Kostenbeiträge \newline den Höchstbetrag der Kostenbeiträge \newline aus der Anlage zur Rechtsverordnung \newline nicht überschreiten. Bei Leistungen \newline nach \S{} 41 steht die Heranziehung der \newline Elternteile unter der Bedingung, dass \newline die Elternteile Kindergeld für den jungen \newline Menschen erhalten; die Heranziehung \newline ist auf die Höhe des Kindergeldes \newline begrenzt. Für die Festsetzung der \newline Kostenbeiträge von Eltern für \newline Leistungen nach \S{} 91 Absatz 4 werden \newline in der Rechtsverordnung die Anteile der \newline Beteiligung an den Kosten der Leistung \newline bestimmt. Die Rechtsverordnung \newline benennt für Kostenbeiträge nach \S{} 93 \newline Absatz 1 Satz 3 einen Höchstbetrag \newline nach \S{} 93 Absatz 1 Satz 4, der sich an \newline der Höhe des geleisteten \newline Lebensunterhalts orientiert}. \\
\hline
 & \textbf{(6) Elternteile werden nachrangig \newline gegenüber den jungen Menschen zu den \newline Kosten herangezogen. Die Höhe des \newline Kostenbeitrags des jungen Menschen \newline wird auf den Kostenbeitrag der \newline Elternteile angerechnet. Das Nähere \newline bestimmt die Rechtsverordnung nach \newline Absatz 5. } & \textbf{(6) Elternteile werden nachrangig \newline gegenüber den jungen Menschen zu den \newline Kosten herangezogen. Die Höhe des \newline Kostenbeitrags des jungen Menschen \newline wird auf den Kostenbeitrag der \newline Elternteile angerechnet. Das Nähere \newline bestimmt die Rechtsverordnung nach \newline Absatz 5. } \\
\hline
\multicolumn{3}{|c|}{\cellcolor{gray!10}\textbf{\S{} 97a -- Pflicht zur Auskunft}} \\
\hline
Absatz 5 & \textit{unverändert} & \textit{unverändert} \\
\hline
(4) Kommt eine der nach den Absätzen 1 \newline und 2 zur Auskunft verpflichteten Personen \newline ihrer Pflicht nicht nach oder bestehen \newline tatsächliche Anhaltspunkte für die \newline Unrichtigkeit ihrer Auskunft, so ist der \newline Arbeitgeber dieser Person verpflichtet, dem \newline örtlichen Träger über die Art des \newline Beschäftigungsverhältnisses und den \newline Arbeitsverdienst dieser Person Auskunft zu \newline geben; Absatz 3 Satz 2 gilt entsprechend. \newline Der zur Auskunft verpflichteten Person ist \newline vor einer Nachfrage beim Arbeitgeber eine \newline angemessene Frist zur Erteilung der \newline Auskunft zu setzen. Sie ist darauf \newline hinzuweisen, dass nach Fristablauf die \newline erforderlichen Auskünfte beim Arbeitgeber \newline eingeholt werden. & \textit{unverändert} & (4) \textbf{Bestehen tatsächliche Anhaltspunkte \newline für die Unrichtigkeit der Auskunft einer \newline nach den Absätzen 1 und 2 zur Auskunft \newline verpflichteten Person, so ist der \newline Arbeitgeber dieser Person verpflichtet, \newline dem örtlichen Träger über die Art des \newline Beschäftigungsverhältnisses und den \newline Arbeitsverdienst dieser Person \newline Auskunft zu geben; Absatz 3 Satz 2 gilt \newline entsprechend.} Der zur Auskunft \newline verpflichteten Person ist vor einer \newline Nachfrage beim Arbeitgeber eine \newline angemessene Frist zur Erteilung der \newline Auskunft zu setzen. Sie ist darauf \newline hinzuweisen, dass nach Fristablauf die \newline erforderlichen Auskünfte beim Arbeitgeber \newline eingeholt werden. \\
\hline
\multicolumn{3}{|c|}{\cellcolor{gray!10}\textbf{\S{} 104 -- Bußgeldvorschriften}} \\
\hline
(1) Ordnungswidrig handelt, wer \newline 1. ohne Erlaubnis nach \S{} 43 Absatz 1 oder \newline \S{} 44 Absatz 1 Satz 1 ein Kind oder einen \newline Jugendlichen betreut oder ihm Unterkunft \newline gewährt, \newline 2. entgegen \S{} 45 Absatz 1 Satz 1, auch in \newline Verbindung mit \S{} 48a Absatz 1, ohne \newline Erlaubnis eine Einrichtung oder eine \newline sonstige Wohnform betreibt oder \newline 3. entgegen \S{} 47 eine Anzeige nicht, nicht \newline richtig, nicht vollständig oder nicht \newline rechtzeitig erstattet oder eine Meldung nicht, \newline nicht richtig, nicht vollständig oder nicht \newline rechtzeitig macht oder vorsätzlich oder \newline fahrlässig seiner Verpflichtung zur \newline Dokumentation oder Aufbewahrung & \textit{unverändert} & (1) Ordnungswidrig handelt, wer \newline 1. ohne Erlaubnis nach \S{} 43 Absatz 1 oder \newline \S{} 44 Absatz 1 Satz 1 ein Kind oder einen \newline Jugendlichen betreut oder ihm Unterkunft \newline gewährt, \newline 2. entgegen \S{} 45 Absatz 1 Satz 1, auch in \newline Verbindung mit \S{} 48a Absatz 1, ohne \newline Erlaubnis eine Einrichtung oder eine \newline sonstige Wohnform betreibt oder \newline 3. entgegen \S{} 47 eine Anzeige nicht, nicht \newline richtig, nicht vollständig oder nicht \newline rechtzeitig erstattet oder eine Meldung nicht, \newline nicht richtig, nicht vollständig oder nicht \newline rechtzeitig macht oder vorsätzlich oder \newline fahrlässig seiner Verpflichtung zur \newline Dokumentation oder Aufbewahrung \\
\hline
derselben oder zum Nachweis der \newline ordnungsgemäßen Buchführung auf \newline entsprechendes Verlangen nicht \newline nachkommt oder \newline 4. entgegen \S{} 97a Absatz 4 vorsätzlich oder \newline fahrlässig als Arbeitgeber eine Auskunft \newline nicht, nicht richtig oder nicht vollständig \newline erteilt. & \textit{unverändert} & derselben oder zum Nachweis der \newline ordnungsgemäßen Buchführung auf \newline entsprechendes Verlangen nicht \newline nachkommt oder \newline 4. entgegen \S{} 97a Absatz 4 vorsätzlich oder \newline fahrlässig als Arbeitgeber eine Auskunft \newline nicht, nicht richtig oder nicht vollständig \newline erteilt \textbf{oder \newline 5. entgegen \S{} 42e vorsätzlich einen \newline anderen gewöhnlichen Aufenthalt \newline nimmt.} \\
\hline
(2) Die Ordnungswidrigkeiten nach Absatz 1 \newline Nummer 1, 3 und 4 können mit einer \newline Geldbuße bis zu fünfhundert Euro, die \newline Ordnungswidrigkeit nach Absatz 1 Nummer \newline 2 kann mit einer Geldbuße bis zu \newline fünfzehntausend Euro geahndet werden. & \textit{unverändert} & (2) Die Ordnungswidrigkeiten nach Absatz 1 \newline Nummer 1, 3, \textbf{4 und 5} können mit einer \newline Geldbuße bis zu fünfhundert Euro, die \newline Ordnungswidrigkeit nach Absatz 1 Nummer \newline 2 kann mit einer Geldbuße bis zu \newline fünfzehntausend Euro geahndet werden. \\
\hline
\multicolumn{3}{|c|}{\cellcolor{gray!10}\textbf{\S{} 108 -- Übergangsregelung}} \\
\hline
(1) Das Bundesministerium für Familie, \newline Senioren, Frauen und Jugend begleitet und \newline untersucht & (1) entfällt & \textbf{(1) Das Bundesministerium für Familie, \newline Senioren, Frauen und Jugend \newline untersucht unter Beteiligung der Länder \newline das Kinder- und \newline Jugendstärkungsgesetz vom 3. Juni \newline 2024 (BGBl. 9. Juni 2021, 29) \newline einschließlich der Regelungen des \newline Bundesgesetzes im Sinne des Artikel 1 \newline Nummer 12 \S{} 10 Absatz 3 Satz 3 des \newline Kinder- und Jugendstärkungsgesetzes \newline auf seine Wirkungen. Es wird \newline untersucht, inwiefern die Regelungen \newline das Ziel, gesellschaftliche Teilhabe und \newline Chancengleichheit für alle jungen \newline Menschen zu sichern oder herzustellen, \newline erreicht werden konnte. Zudem wird \newline untersucht, welche finanziellen \newline Auswirkungen die Regelungen auf \newline Länder und Kommunen haben. Als \newline Kriterien für die Evaluation dienen \newline insbesondere die Vollständigkeit der \newline Umsetzung der Regelungen sowie die \newline Inanspruchnahme von Leistungen auch \newline unter Berücksichtigung der Perspektive \newline der Normadressatinnen und - \newline adressaten. Als Grundlage dienen die } \newline \textbf{Daten der Kinder- und \newline Jugendhilfestatistik. Das \newline Bundesministerium für Familien, \newline Senioren, Frauen und Jugend berichtet \newline dem Deutschen Bundestag und dem \newline Bundesrat über die Ergebnisse dieser \newline Untersuchung}. \\
\hline
1. \newline bis zum Inkrafttreten von \S{} 10b am 1. \newline Januar 2024 sowie & 1. entfällt & 1. entfällt \\
\hline
2. \newline bis zum Inkrafttreten von \S{} 10 Absatz 4 \newline Satz 1 und 2 am 1. Januar 2028 die \newline Umsetzung der für die Ausführung dieser \newline Regelungen jeweils notwendigen \newline Maßnahmen in den Ländern. Bei der \newline Untersuchung nach Satz 1 Nummer 1 \newline werden insbesondere auch die \newline Erfahrungen der örtlichen Träger der \newline öffentlichen Jugendhilfe einbezogen, die \newline bereits vor dem 1. Januar 2024 \newline Verfahrenslotsen entsprechend \S{} 10b \newline einsetzen. Bei der Untersuchung nach Satz \newline 1 Nummer 2 findet das Bundesgesetz nach \newline \S{} 10 Absatz 4 Satz 3 ab dem Zeitpunkt \newline seiner Verkündung, die als Bedingung für \newline das Inkrafttreten von \S{} 10 Absatz 4 Satz 1 \newline und 2 spätestens bis zum 1. Januar 2027 \newline erfolgen muss, besondere \newline Berücksichtigung. & 2. entfällt & 2. entfällt \\
\hline
(2) Das Bundesministerium für Familie, \newline Senioren, Frauen und Jugend untersucht in \newline den Jahren 2022 bis 2024 die rechtlichen \newline Wirkungen von \S{} 10 Absatz 4 und legt dem \newline Bundestag und dem Bundesrat bis zum 31. \newline Dezember 2024 einen Bericht über das \newline Ergebnis der Untersuchung vor. Dabei \newline sollen insbesondere die gesetzlichen \newline Festlegungen des Achten und Neunten \newline Buches & (2) entfällt & \textbf{(2) Es wird ein Konzept zur künftigen \newline inhaltlichen Ausgestaltung der Kinder- \newline und Jugendhilfestatistik entwickelt, auf \newline dessen Grundlage in einem \newline Bundesgesetz die zur Beurteilung der \newline Auswirkungen der Bestimmungen \newline dieses Buches und zu seiner \newline Fortentwicklung notwendigen laufenden \newline Erhebungen auch im Hinblick auf die ab \newline dem 1. Januar 2028 vorrangige \newline Zuständigkeit der Kinder- und \newline Jugendhilfe für Leistungen der \newline Eingliederungshilfe für junge Menschen \newline mit körperlichen oder geistigen \newline Behinderungen oder hiervon bedrohte \newline junge Menschen geregelt und 2030 \newline beginnend durchgeführt werden \newline können. } \\
\hline
1. \newline zur Bestimmung des leistungsberechtigten \newline Personenkreises, & 1. entfällt & 1. entfällt \\
\hline
2. \newline zur Bestimmung von Art und Umfang der \newline Leistungen, & 2. entfällt & 2. entfällt \\
\hline
3. \newline zur Ausgestaltung der Kostenbeteiligung \newline bei diesen Leistungen und & 3. entfällt & 3. entfällt \\
\hline
4. \newline zur Ausgestaltung des Verfahrens \newline untersucht werden mit dem Ziel, den \newline leistungsberechtigten Personenkreis, Art \newline und Umfang der Leistungen sowie den \newline Umfang der Kostenbeteiligung für die \newline hierzu Verpflichteten nach dem am 1. \newline Januar 2023 für die Eingliederungshilfe \newline geltenden Recht beizubehalten, \newline insbesondere einerseits keine \newline Verschlechterungen für \newline leistungsberechtigte oder \newline kostenbeitragspflichtige Personen und \newline andererseits keine Ausweitung des Kreises \newline der Leistungsberechtigten sowie des \newline Leistungsumfangs im Vergleich zur \newline Rechtslage am 1. Januar 2023 \newline herbeizuführen, sowie Hinweise auf die zu \newline bestimmenden Inhalte des \newline Bundesgesetzes nach \S{} 10 Absatz 4 Satz \newline 3 zu geben. In die Untersuchung werden \newline auch mögliche finanzielle Auswirkungen \newline gesetzlicher Gestaltungsoptionen \newline einbezogen. & 4. entfällt & 4. entfällt \\
\hline
(3) Soweit das Bundesministerium für \newline Familie, Senioren, Frauen und Jugend \newline Dritte in die Durchführung der \newline Untersuchungen nach den Absätzen 1 und \newline 2 einbezieht, beteiligt es hierzu vorab die \newline Länder. & (3) entfällt & (3) entfällt \\
\hline
(4) Das Bundesministerium für Familie, \newline Senioren, Frauen und Jugend untersucht \newline unter Beteiligung der Länder die Wirkungen \newline dieses Gesetzes im Übrigen einschließlich \newline seiner finanziellen Auswirkungen auf \newline Länder und Kommunen und berichtet dem \newline Deutschen Bundestag und dem Bundesrat \newline über die Ergebnisse dieser Untersuchung. & Das Bundesministerium für Familie, \newline Senioren, Frauen und Jugend untersucht \newline unter Beteiligung der Länder die Wirkungen \newline \textbf{des Kinder- und \newline Jugendstärkungsgesetzes vom 3. Juni \newline 2024 (BGBl. 9. Juni 2021, 29)} im Übrigen \newline einschließlich seiner finanziellen \newline Auswirkungen auf Länder und Kommunen \newline und berichtet dem Deutschen Bundestag \newline und dem Bundesrat über die Ergebnisse \newline dieser Untersuchung. &  \\
\hline
\multicolumn{3}{|c|}{\cellcolor{gray!10}\textbf{\S{} 109 -- Übergangsregelungen}} \\
\hline
 & \textbf{(1) Die Leistungs- und \newline Vergütungsvereinbarungen nach Kapitel \newline 8 des Teils 2 des Neunten Buches \newline gelten für die in \S{} 78a benannten \newline Leistungen als Vereinbarungen nach \S{} \newline 78b und bei ambulanten Leistungen als \newline Vereinbarungen nach \S{} 77 SGB VIII bis \newline zum 31. Dezember 2032 fort. Die \newline Vereinbarungen umfassen die \newline Leistungen für minderjährige \newline Leistungsberechtigte, auf die sich die \newline Leistungs- und \newline Vergütungsvereinbarungen im Sinne \newline des Satz 1 bisher bezogen haben, sowie \newline Leistungen nach \S{} 41, die inhaltlich den \newline bisher vereinbarten Leistungen \newline entsprechen. } & \textbf{(1) Die Leistungs- und \newline Vergütungsvereinbarungen nach Kapitel \newline 8 des Teils 2 des Neunten Buches \newline gelten für die in \S{} 78a benannten \newline Leistungen als Vereinbarungen nach \S{} \newline 78b und bei ambulanten Leistungen als \newline Vereinbarungen nach \S{} 77 bis zum 31. \newline Dezember 2032 fort. Die \newline Vereinbarungen, die als Vereinbarungen \newline nach \S{} 78b oder \S{} 77 fortgelten, \newline umfassen die Leistungen für \newline minderjährige Leistungsberechtigte, auf \newline die sich die Leistungs- und \newline Vergütungsvereinbarungen im Sinne \newline des Satzes 1 bisher bezogen haben, \newline sowie Leistungen nach \S{} 41, die \newline inhaltlich den bisher vereinbarten \newline Leistungen entsprechen. } \\
\hline
 & \textbf{(2) Leistungen nach dem Neunten Buch \newline für junge Menschen, die vor dem 1. \newline Januar 2028 das 18. Lebensjahr \newline vollendet haben, gehen Leistungen nach \newline diesem Buch vor. } & \textbf{(2) Leistungen nach dem Neunten Buch \newline für junge Menschen, die vor dem 1. \newline Januar 2028 das 18. Lebensjahr \newline vollendet haben, gehen Leistungen nach \newline diesem Buch vor. } \\
\hline
 & \textbf{(3) Die Leistungsbescheide auf \newline Grundlage des \S{} 99 des Neunten \newline Buches gelten als Bescheide nach \S{} 27 \newline Absatz 3 fort. } & \textbf{(3) Die Leistungsbescheide für \newline minderjährige Leistungsberechtigte auf \newline Grundlage des \S{} 99 des Neunten \newline Buches gelten als Bescheide nach \S{} 27 \newline Absatz 3 fort. } \\
\hline
 & \textbf{(4) Jede Vertragspartei der Leistungs- \newline und Vergütungsvereinbarungen nach \newline Ab-satz 1 hat zum 1. Januar 2028 \newline unbeschadet der Laufzeit der nach \newline Absatz 1 und 2 fort-geltenden Verträge \newline und Leistungsbestandteile einen \newline Anspruch auf Neuverhandlung einer \newline Vereinbarung nach \S{} 78b SGB VIII. } & \textbf{(4) Jede Vertragspartei der Leistungs- \newline und Vergütungsvereinbarungen nach \newline Absatz 1 hat zum 1. Januar 2028 \newline unbeschadet der Laufzeit der nach \newline Absatz 1 und 2 fortgeltenden Verträge \newline und Leistungsbestandteile einen \newline Anspruch auf Neuverhandlung einer \newline Vereinbarung nach \S{} 78b SGB VIII.} \textbf{Die \newline Frist nach \S{} 78g Absatz 2 Satz1 wird im \newline Falle von Neuverhandlungen nach Satz \newline 1 um 12 Wochen verlängert. } \\
\hline
 & \textbf{(5) Die Bescheide zur Festsetzung des \newline Beitrags aus Einkommen zu den \newline Aufwendungen nach \S{} 136 des Neunten \newline Buches sowie die Festsetzung der \newline Aufbringung der Mittel für die Kosten \newline des Lebensunterhaltes nach \S{} 142 des \newline Neunten Buches gelten bis zu ihrer \newline Aufhebung fort, sofern die Leistung \newline nach \S{} 99 des Neunten Buches \newline entsprechend des Absatz 2 ab dem 1. \newline Januar 2028 als Leistung nach \S{} 27 \newline Absatz 3 fortgesetzt wird. Für den \newline Träger der Eingliederungshilfe tritt der \newline zuständige Träger der öffentlichen \newline Jugendhilfe ein. Die Aufhebung der \newline Bescheide muss rückwirkend zum 1. \newline Januar 2028 erfolgen; die Aufhebung \newline soll spätestens am 31. Dezember 2028 \newline bei den durch die Bescheide \newline Verpflichteten eingehen}. & \textbf{(5) Die Bescheide zur Festsetzung des \newline Beitrags aus Einkommen zu den } \newline \textbf{Aufwendungen nach \S{} 136 des Neunten \newline Buches sowie die Festsetzung der \newline Aufbringung der Mittel für die Kosten \newline des Lebensunterhaltes nach \S{} 142 des \newline Neunten Buches gelten bis zu ihrer \newline Aufhebung fort, sofern die Leistung \newline nach \S{} 99 des Neunten Buches \newline entsprechend des Absatzes 3 ab dem 1. \newline Januar 2028 als Leistung nach \S{} 27 \newline Absatz 3 fortgesetzt wird. An die Stelle \newline des Trägers der Eingliederungshilfe tritt \newline der zuständige Träger der öffentlichen \newline Jugendhilfe. Die Aufhebung der \newline Bescheide muss rückwirkend zum 1. \newline Januar 2028 erfolgen; die Aufhebung \newline muss spätestens am 31. Dezember 2028 \newline den durch die Bescheide Verpflichteten \newline zugehen. } \\
\hline
 & \textbf{(6) Abweichend von den \S{}\S{} 91 bis 94 gilt \newline für den Kostenbeitrag für die \newline Erbringung von Leistungen für \newline Leistungsberechtigte auf der Grundlage \newline von Leistungsbescheiden im Sinne des \newline Absatz 2 das Folgende: } & \textbf{(6) Abweichend von den \S{}\S{} 91 bis 94 gilt \newline für den Kostenbeitrag für die \newline Erbringung von Leistungen für \newline Leistungsberechtigte auf der Grundlage \newline von Leistungsbescheiden im Sinne des \newline Absatz 2 das Folgende: } \newline \textbf{den \S{}\S{} 91 bis 94. Die Sätze 1 bis 3 gelten \newline für die Kostenbeiträge von Elternteilen \newline entsprechend, die auf der Grundlage der \newline \S{}\S{} 91 bis 94 mit Gültigkeit vom 31. \newline Dezember 2027 zu den Kosten \newline herangezogen wurden. } \\
\hline
 & \textbf{1. \newline Wurde mindestens von einem Elternteil \newline des Leistungsberechtigten der Einsatz \newline des Einkommens nach \S{} 136 des \newline Neunten Buches oder die Aufbringung \newline der Mittel für die Kosten des \newline Lebensunterhalts in Höhe der für den \newline häuslichen Lebensunterhalt ersparten \newline Aufwendungen nach \S{} 142 des Neunten \newline Buches gefordert und ist der nach den \newline \S{}\S{} 91 bis 94 aufzubringende Beitrag \newline höher als der Einkommenseinsatz oder \newline die Höhe der aufzubringenden ersparten \newline Aufwendungen } & \textbf{1. \newline Wurde mindestens von einem Elternteil \newline des Leistungsberechtigten der Einsatz \newline des Einkommens nach \S{} 136 des \newline Neunten Buches oder die Aufbringung \newline der Mittel für die Kosten des \newline Lebensunterhalts in Höhe der für den \newline häuslichen Lebensunterhalt ersparten \newline Aufwendungen nach \S{} 142 des Neunten \newline Buches gefordert und ist der nach den \newline \S{}\S{} 91 bis 94 aufzubringende Betrag \newline höher als der Einkommenseinsatz oder \newline als die aufzubringenden ersparten \newline Aufwendungen nach Kapitel 9 des Teils \newline 2 des Neunten Buches mit Gültigkeit \newline vom 31. Dezember 2027, so ist der \newline Kostenbeitrag nach den \S{}\S{} 91 bis 94 auf \newline diesen Betrag begrenzt. Die Begrenzung \newline gilt auch dann, wenn bis zum 31. \newline Dezember 2027 nur ein Elternteil zu den \newline Kosten herangezogen wurde und nach \newline den \S{}\S{} 91 bis 94 beide Elternteile \newline getrennt zu den Kosten herangezogen \newline werden. Der bisher von einem Elternteil \newline aufgebrachte Betrag gilt dann als \newline Höchstbetrag für die Summe der \newline Kostenbeiträge beider Elternteile nach } \\
\hline
 & \textbf{2. \newline Wurde bisher von keinem Elternteil der \newline Einsatz des Einkommens nach \S{} 136 \newline des Neunten Buches oder die \newline Aufbringung der Mittel für die Kosten \newline des Lebensunter-halts nach \S{} 142 des \newline Neunten Buches gefordert und wird die \newline bis zum 31. Dezember 2027 erbrachte \newline Leistung nach \S{} 99 des Neunten Buchs \newline auf der Grundlage des \S{} 27 Absatz 3 \newline fortgesetzt oder neu bewilligt, so wird \newline für diese Leistung kein Kostenbeitrag \newline nach den \S{}\S{} 91 bis 94 erhoben. } & \textbf{2. \newline Wurde bisher von keinem Elternteil der \newline Einsatz des Einkommens nach \S{} 136 \newline des Neunten Buches oder die \newline Aufbringung der Mittel für die Kosten \newline des Lebensunterhalts nach \S{} 142 des \newline Neunten Buches gefordert und wird die \newline bis zum 31. Dezember 2027 erbrachte \newline Leistung nach \S{} 99 des Neunten Buchs \newline auf der Grundlage des \S{} 27 Absatz 3 \newline fortgesetzt oder neu bewilligt, so wird \newline für diese Leistung kein Kostenbeitrag \newline nach den \S{}\S{} 91 bis 94 erhoben. } \\
\hline
 & \textit{unverändert} & \textbf{(7) Für Leistungen auf der Grundlage \newline von Bescheiden nach Absatz 3 richtet \newline sich die örtliche Zuständigkeit nach \S{}\S{} \newline 86, 86c, 86d und 88. Spätestens bis zum \newline 31. Oktober 2028 ist die örtliche \newline Zuständigkeit für Leistungen nach Satz \newline 1 zu prüfen. Fand die Übergabe eines \newline Falles der Gewährung von Leistungen \newline nach Satz 1 an den nach Satz 1 \newline zuständigen Träger der Jugendhilfe \newline nicht zum 1. Januar 2028 statt, sind dem \newline Träger, der bis zum 31. Dezember 2027 \newline zuständig war, die Kosten bis zur \newline tatsächlichen Übergabe zu erstatten. } \\
\hline
\multicolumn{3}{|c|}{\cellcolor{gray!10}\textbf{\S{} 98 -- Zweck und Umfang der Erhebung}} \\
\hline
(1) Zur Beurteilung der Auswirkungen der \newline Bestimmungen dieses Buches und zu \newline seiner Fortentwicklung sind laufende \newline Erhebungen über & \textit{(1) unverändert} & \textit{unverändert} \\
\hline
1. \newline Kinder und tätige Personen in \newline Tageseinrichtungen, & 1. unverändert & \textit{unverändert} \\
\hline
1a. \newline Kinder in den Klassenstufen eins bis vier, & 1a. unverändert & \textit{unverändert} \\
\hline
2. \newline Kinder und tätige Personen in öffentlich \newline geförderter Kindertagespflege, & 2. unverändert & \textit{unverändert} \\
\hline
3. \newline Personen, die mit öffentlichen Mitteln \newline geförderte Kindertagespflege gemeinsam \newline oder auf Grund einer Erlaubnis nach \S{} 43 \newline Absatz 3 Satz 3 in Pflegestellen \newline durchführen, und die von diesen betreuten \newline Kinder, & 3. unverändert & \textit{unverändert} \\
\hline
4. \newline die Empfänger & 4. unverändert & \textit{unverändert} \\
\hline
c) \newline der Eingliederungshilfe für seelisch \newline behinderte Kinder und Jugendliche, & c) \newline der Eingliederungshilfe für Kinder und \newline Jugendliche \textbf{mit Behinderungen oder \newline drohenden Behinderungen}, & \textit{unverändert} \\
\hline
5. \newline Kinder und Jugendliche, zu deren Schutz \newline vorläufige Maßnahmen getroffen worden \newline sind, & 5. unverändert & \textit{unverändert} \\
\hline
6. \newline Kinder und Jugendliche, die als Kind \newline angenommen worden sind, & 6. unverändert & \textit{unverändert} \\
\hline
7. \newline Kinder und Jugendliche, die unter \newline Amtspflegschaft, Amtsvormundschaft oder \newline Beistandschaft des Jugendamts stehen, & 7. unverändert & \textit{unverändert} \\
\hline
8. \newline Kinder und Jugendliche, für die eine \newline Pflegeerlaubnis erteilt worden ist, & 8. unverändert & \textit{unverändert} \\
\hline
9. \newline Maßnahmen des Familiengerichts, & 9. unverändert & \textit{unverändert} \\
\hline
10. \newline Angebote der Jugendarbeit nach \S{} 11 \newline sowie Fortbildungsmaßnahmen für \newline ehrenamtliche Mitarbeiter anerkannter \newline Träger der Jugendhilfe nach \S{} 74 Absatz 6, & 10. unverändert & \textit{unverändert} \\
\hline
11. \newline die Träger der Jugendhilfe, die dort tätigen \newline Personen und deren Einrichtungen mit \newline Ausnahme der Tageseinrichtungen, & 11. unverändert & \textit{unverändert} \\
\hline
12. \newline die Ausgaben und Einnahmen der \newline öffentlichen Jugendhilfe sowie & 12. unverändert & \textit{unverändert} \\
\hline
13. \newline Gefährdungseinschätzungen nach \S{} 8a als \newline Bundesstatistik durchzuführen. & 13. unverändert & \textit{unverändert} \\
\hline
(2) Zur Verfolgung der gesellschaftlichen \newline Entwicklung im Bereich der elterlichen \newline Sorge sind im Rahmen der Kinder- und \newline Jugendhilfestatistik auch laufende \newline Erhebungen über Sorgeerklärungen \newline durchzuführen. & \textit{(2) unverändert} & \textit{unverändert} \\
\hline
\multicolumn{3}{|c|}{\cellcolor{gray!10}\textbf{\S{} 99 -- Erhebungsmerkmale}} \\
\hline
(1) Erhebungsmerkmale bei den \newline Erhebungen über Hilfe zur Erziehung nach \newline den \S{}\S{} 27 bis 35, Eingliederungshilfe für \newline seelisch behinderte Kinder und \newline Jugendliche nach \S{} 35a und Hilfe für junge \newline Volljährige nach \S{} 41 sind & (1) Erhebungsmerkmale bei den \newline Erhebungen über Hilfe zur Erziehung nach \newline den \S{}\S{} 27 \textbf{Absatz 2, 27a} bis 35, \newline Eingliederungshilfe für Kinder und \newline Jugendliche \textbf{mit Behinderungen oder \newline drohenden Behinderungen} nach \S{} 35a \newline \textbf{bis \S{} 35i} und Hilfe für junge Volljährige \newline nach \S{} 41 sind & \textit{unverändert} \\
\hline
1. im Hinblick auf die Hilfe & 1. im Hinblick auf die Hilfe \textbf{oder Leistung} & \textit{unverändert} \\
\hline
c) \newline Bundesland des überwiegenden \newline Einsatzortes. & c) unverändert & \textit{unverändert} \\
\hline
(2) Erhebungsmerkmale bei den \newline Erhebungen über vorläufige Maßnahmen \newline zum Schutz von Kindern und Jugendlichen \newline sind Kinder und Jugendliche, zu deren \newline Schutz Maßnahmen nach \S{} 42 oder \S{} 42a \newline getroffen worden sind, gegliedert nach & \textit{(2) unverändert} & \textit{unverändert} \\
\hline
1. \newline Art der Maßnahme, Art des Trägers der \newline Maßnahme, Form der Unterbringung \newline während der Maßnahme, hinweisgebender \newline Institution oder Person, Zeitpunkt des \newline Beginns und Dauer der Maßnahme, \newline Durchführung aufgrund einer \newline vorangegangenen \newline Gefährdungseinschätzung nach \S{} 8a \newline Absatz 1, Maßnahmeanlass, im \newline Kalenderjahr bereits wiederholt \newline stattfindende Inobhutnahme, Widerspruch \newline der Personensorge- oder \newline Erziehungsberechtigten gegen die \newline Maßnahme, im Fall des Widerspruchs \newline gegen die Maßnahme Herbeiführung einer \newline Entscheidung des Familiengerichts nach \S{} \newline 42 Absatz 3 Satz 2 Nummer 2, Grund für \newline die Beendigung der Maßnahme, \newline anschließendem Aufenthalt, Art der \newline anschließenden Hilfe, & 1. unverändert & \textit{unverändert} \\
\hline
2. \newline bei Kindern und Jugendlichen zusätzlich zu \newline den unter Nummer 1 genannten \newline Merkmalen nach Geschlecht, Altersgruppe \newline zu Beginn der Maßnahme, ausländischer \newline Herkunft mindestens eines Elternteils, \newline Deutsch als in der Familie vorrangig \newline gesprochene Sprache, Art des Aufenthalts \newline vor Beginn der Maßnahme. & 2. unverändert & \textit{unverändert} \\
\hline
(3) Erhebungsmerkmale bei den \newline Erhebungen über die Annahme als Kind \newline sind & \textit{(3) unverändert} & \textit{unverändert} \\
\hline
1. \newline angenommene Kinder und Jugendliche, \newline gegliedert & 1. unverändert & \textit{unverändert} \\
\hline
(4) Erhebungsmerkmal bei den Erhebungen \newline über die Amtspflegschaft und die \newline Amtsvormundschaft sowie die \newline Beistandschaft ist die Zahl der Kinder und \newline Jugendlichen unter & \textit{(4) unverändert} & \textit{unverändert} \\
\hline
1. \newline gesetzlicher Amtsvormundschaft, & 1. unverändert & \textit{unverändert} \\
\hline
2. \newline bestellter Amtsvormundschaft, & 2. unverändert & \textit{unverändert} \\
\hline
3. \newline bestellter Amtspflegschaft sowie & 3. unverändert & \textit{unverändert} \\
\hline
4. \newline Beistandschaft, & 4. unverändert & \textit{unverändert} \\
\hline
(5) Erhebungsmerkmal bei den Erhebungen \newline über & \textit{(5) unverändert} & \textit{unverändert} \\
\hline
1. \newline die Pflegeerlaubnis nach \S{} 43 ist die Zahl \newline der Tagespflegepersonen, & 1. unverändert & \textit{unverändert} \\
\hline
2. \newline die Pflegeerlaubnis nach \S{} 44 ist die Zahl \newline der Kinder und Jugendlichen, gegliedert \newline nach Geschlecht und Art der Pflege. & 2. unverändert & \textit{unverändert} \\
\hline
(6) Erhebungsmerkmale bei der Erhebung \newline zum Schutzauftrag bei \newline Kindeswohlgefährdung nach \S{} 8a sind \newline Kinder und Jugendliche, bei denen eine \newline Gefährdungseinschätzung nach Absatz 1 \newline vorgenommen worden ist, gegliedert & \textit{(6) unverändert} & \textit{unverändert} \\
\hline
1. \newline nach der hinweisgebenden Institution oder \newline Person, der Art der Kindeswohlgefährdung, \newline der Person, von der die Gefährdung \newline ausgeht, dem Ergebnis der \newline Gefährdungseinschätzung sowie \newline wiederholter Meldung zu demselben Kind \newline oder Jugendlichen im jeweiligen \newline Kalenderjahr, & 1. unverändert & \textit{unverändert} \\
\hline
2. \newline bei Kindern und Jugendlichen zusätzlich zu \newline den in Nummer 1 genannten Merkmalen \newline nach Geschlecht, Geburtsmonat, \newline Geburtsjahr, ausländischer Herkunft \newline mindestens eines Elternteils, Deutsch als in \newline der Familie vorrangig gesprochene \newline Sprache, Eingliederungshilfe und \newline Aufenthaltsort des Kindes oder \newline Jugendlichen zum Zeitpunkt der Meldung \newline sowie den Altersgruppen der Eltern und der \newline Inanspruchnahme einer Leistung gemäß \newline den \S{}\S{} 16 bis 19 sowie 27 bis 35a und der \newline Durchführung einer Maßnahme nach \S{} 42. & 2. \newline bei Kindern und Jugendlichen zusätzlich zu \newline den in Nummer 1 genannten Merkmalen \newline nach Geschlecht, Geburtsmonat, \newline Geburtsjahr, ausländischer Herkunft \newline mindestens eines Elternteils, Deutsch als in \newline der Familie vorrangig gesprochene \newline Sprache, \textbf{Leistung der} Eingliederungshilfe \newline und Aufenthaltsort des Kindes oder \newline Jugendlichen zum Zeitpunkt der Meldung \newline sowie den Altersgruppen der Eltern und der \newline Inanspruchnahme einer Leistung gemäß \newline den \S{}\S{} 16 bis 19 sowie 27, \textbf{27a} bis 35 und \newline der Durchführung einer Maßnahme nach \S{} \newline 42. & \textit{unverändert} \\
\hline
(6a) Erhebungsmerkmal bei den \newline Erhebungen über Sorgeerklärungen und \newline die gerichtliche Übertragung der \newline gemeinsamen elterlichen Sorge nach \S{} \newline 1626a Absatz 1 Nummer 3 des \newline Bürgerlichen Gesetzbuchs ist die \newline gemeinsame elterliche Sorge nicht \newline miteinander verheirateter Eltern, gegliedert \newline danach, ob Sorgeerklärungen beider Eltern \newline vorliegen oder den Eltern die elterliche \newline Sorge aufgrund einer gerichtlichen \newline Entscheidung ganz oder zum Teil \newline gemeinsam übertragen worden ist. & \textit{(6a) unverändert} & \textit{unverändert} \\
\hline
(6b) Erhebungsmerkmal bei den \newline Erhebungen über Maßnahmen des \newline Familiengerichts ist die Zahl der Kinder und \newline Jugendlichen, bei denen wegen einer \newline Gefährdung ihres Wohls das \newline familiengerichtliche Verfahren auf Grund \newline einer Anrufung durch das Jugendamt nach \newline \S{} 8a Absatz 2 Satz 1 oder \S{} 42 Absatz 3 \newline Satz 2 Nummer 2 oder auf andere Weise \newline eingeleitet worden ist und & \textit{(6b) unverändert} & \textit{unverändert} \\
\hline
1. \newline den Personensorgeberechtigten auferlegt \newline worden ist, Leistungen nach diesem Buch \newline in Anspruch zu nehmen, & 1. unverändert & \textit{unverändert} \\
\hline
2. \newline andere Gebote oder Verbote gegenüber \newline den Personensorgeberechtigten oder \newline Dritten ausgesprochen worden sind, & 2. unverändert & \textit{unverändert} \\
\hline
3. \newline Erklärungen der \newline Personensorgeberechtigten ersetzt worden \newline sind, & 3. unverändert & \textit{unverändert} \\
\hline
4. \newline die elterliche Sorge ganz oder teilweise \newline entzogen und auf das Jugendamt oder \newline einen Dritten als Vormund oder Pfleger \newline übertragen worden ist, & 4. unverändert & \textit{unverändert} \\
\hline
(7) Erhebungsmerkmale bei den \newline Erhebungen über Kinder und tätige \newline Personen in Tageseinrichtungen sind & \textit{(7) unverändert} & \textit{unverändert} \\
\hline
1. \newline die Einrichtungen, gegliedert nach & 1. unverändert & \textit{unverändert} \\
\hline
(7a) \newline Erhebungsmerkmale bei den Erhebungen \newline über Kinder in mit öffentlichen Mitteln \newline geförderter Kindertagespflege sowie die die \newline Kindertagespflege durchführenden \newline Personen sind: & \textit{(7a) unverändert} & \textit{unverändert} \\
\hline
1. \newline für jede tätige Person & 1. unverändert & \textit{unverändert} \\
\hline
(7b) Erhebungsmerkmale bei den \newline Erhebungen über Personen, die mit \newline öffentlichen Mitteln geförderte \newline Kindertagespflege gemeinsam oder auf \newline Grund einer Erlaubnis nach \S{} 43 Absatz 3 \newline Satz 3 durchführen und die von diesen \newline betreuten Kinder sind die Zahl der \newline Kindertagespflegepersonen und die Zahl \newline der von diesen betreuten Kinder jeweils \newline gegliedert nach Pflegestellen. & \textit{(7b) unverändert} & \textit{unverändert} \\
\hline
(7c) Erhebungsmerkmale bei den \newline Erhebungen über Kinder in den \newline Klassenstufen eins bis vier sind & \textit{(7c) unverändert} & \textit{unverändert} \\
\hline
1. \newline Klassenstufe, & 1. unverändert & \textit{unverändert} \\
\hline
2. \newline Anzahl der Wochenstunden, die das Kind in \newline Angeboten nach \S{} 24 Absatz 4 verbringt, & 2. unverändert & \textit{unverändert} \\
\hline
3. \newline Art der Angebote nach \S{} 24 Absatz 4. & 3. unverändert & \textit{unverändert} \\
\hline
(8) Erhebungsmerkmale bei den \newline Erhebungen über die Angebote der \newline Jugendarbeit nach \S{} 11 sowie bei den \newline Erhebungen über Fortbildungsmaßnahmen \newline für ehrenamtliche Mitarbeiter anerkannter \newline Träger der Jugendhilfe nach \S{} 74 Absatz 6 \newline sind offene und Gruppenangebote sowie \newline Veranstaltungen und Projekte der \newline Jugendarbeit, soweit diese mit öffentlichen \newline Mitteln pauschal oder maßnahmenbezogen \newline gefördert werden oder der Träger eine \newline öffentliche Förderung erhält, gegliedert \newline nach & \textit{(8) unverändert} & \textit{unverändert} \\
\hline
1. \newline Art und Rechtsform des Trägers sowie bei \newline Trägern der freien Jugendhilfe deren \newline Verbandszugehörigkeit, & 1. unverändert & \textit{unverändert} \\
\hline
2. \newline Dauer, Häufigkeit, Durchführungsort und \newline Art des Angebots; zusätzlich bei \newline schulbezogenen Angeboten die Art der \newline kooperierenden Schule, & 2. unverändert & \textit{unverändert} \\
\hline
3. \newline Art der Beschäftigung und Tätigkeit der bei \newline der Durchführung des Angebots tätigen \newline Personen sowie, mit Ausnahme der \newline sonstigen pädagogisch tätigen Personen, \newline deren Altersgruppe und Geschlecht, & 3. unverändert & \textit{unverändert} \\
\hline
4. \newline Zahl der Teilnehmenden und der Besucher \newline sowie, mit Ausnahme von Festen, Feiern, \newline Konzerten, Sportveranstaltungen und \newline sonstigen Veranstaltungen, deren \newline Geschlecht und Altersgruppe, & 4. unverändert & \textit{unverändert} \\
\hline
5. \newline Partnerländer und Veranstaltungen im In- \newline oder Ausland bei Veranstaltungen und \newline Projekten der internationalen Jugendarbeit. & 5. unverändert & \textit{unverändert} \\
\hline
(9) Erhebungsmerkmale bei den \newline Erhebungen über die Träger der \newline Jugendhilfe, die dort tätigen Personen und \newline deren Einrichtungen, soweit diese nicht in \newline Absatz 7 erfasst werden, sind & \textit{(9) unverändert} & \textit{unverändert} \\
\hline
1. \newline die Träger gegliedert nach & 1. unverändert & \textit{unverändert} \\
\hline
(10) Erhebungsmerkmale bei der Erhebung \newline der Ausgaben und Einnahmen der \newline öffentlichen Jugendhilfe sind & \textit{(10) unverändert} & \textit{unverändert} \\
\hline
1. \newline die Art des Trägers, & 1. unverändert & \textit{unverändert} \\
\hline
2. \newline die Ausgaben für Einzel- und \newline Gruppenhilfen, gegliedert nach Ausgabe- \newline und Hilfeart sowie die Einnahmen nach \newline Einnahmeart, & 2. \newline die Ausgaben für Einzel- und \newline Gruppenhilfen \textbf{oder Einzel- und \newline Gruppenleistungen}, gegliedert nach \newline Ausgabe- und Hilfe\textbf{- oder} \textbf{Leistungs}art \newline sowie die Einnahmen nach Einnahmeart, & \textit{unverändert} \\
\hline
3. \newline die Ausgaben und Einnahmen für \newline Einrichtungen nach Arten gegliedert nach \newline der Einrichtungsart, & 3.unverändert & \textit{unverändert} \\
\hline
4. \newline die Ausgaben für das Personal, das bei \newline den örtlichen und den überörtlichen \newline Trägern sowie den kreisangehörigen \newline Gemeinden und Gemeindeverbänden, die \newline nicht örtliche Träger sind, Aufgaben der \newline Jugendhilfe wahrnimmt. & 4.unverändert & \textit{unverändert} \\
\hline
\multicolumn{3}{|c|}{\cellcolor{gray!10}\textbf{\S{} 101 -- Periodizität und Berichtszeitraum}} \\
\hline
(1) Die Erhebungen nach \S{} 99 Absatz 1 bis \newline 5, 6a bis 7c und 10 sind jährlich \newline durchzuführen, die Erhebungen nach \S{} 99 \newline Absatz 3 Nummer 3 erstmalig für das Jahr \newline 2022; die Erhebungen nach \S{} 99 Absatz 1, \newline soweit sie die Eingliederungshilfe für \newline Kinder und Jugendliche mit seelischer \newline Behinderung betreffen, sind 2007 \newline beginnend jährlich durchzuführen. Die \newline Erhebung nach \S{} 99 Absatz 6 erfolgt \newline laufend. Die übrigen Erhebungen nach \S{} \newline 99 sind alle zwei Jahre durchzuführen, die \newline Erhebungen nach \S{} 99 Absatz 8 erstmalig \newline für das Jahr 2015 und die Erhebungen \newline nach \S{} 99 Absatz 9 erstmalig für das Jahr \newline 2014. & (1) Die Erhebungen nach \S{} 99 Absatz 1 bis \newline 5, 6a bis 7c und 10 sind jährlich \newline durchzuführen, die Erhebungen nach \S{} 99 \newline Absatz 3 Nummer 3 erstmalig für das Jahr \newline 2022; die Erhebungen nach \S{} 99 Absatz 1, \newline soweit sie die Eingliederungshilfe für \newline Kinder und Jugendliche mit seelischer \newline Behinderung betreffen, sind 2007 \newline beginnend jährlich durchzuführen. Die \newline Erhebung nach \S{} 99 Absatz 6 erfolgt \newline laufend. Die übrigen Erhebungen nach \S{} \newline 99 sind alle zwei Jahre durchzuführen, die \newline Erhebungen nach \S{} 99 Absatz 8 erstmalig \newline für das Jahr 2015 und die Erhebungen \newline nach \S{} 99 Absatz 9 erstmalig für das Jahr \newline 2014. \textbf{Die Erhebungen nach \S{} 99 Absatz \newline 1, soweit sie die Eingliederungshilfe für \newline Kinder und Jugendliche mit (drohenden) \newline geistigen oder körperlichen \newline Behinderungen betreffen, sind 2029 \newline beginnend jährlich durchzuführen.} & \textit{unverändert} \\
\hline
(2) Die Angaben für die Erhebung nach & \textit{(2) unverändert} & \textit{unverändert} \\
\hline
1. \newline \S{} 99 Absatz 1 sind zu dem Zeitpunkt, zu \newline dem die Hilfe endet, bei fortdauernder Hilfe \newline zum 31. Dezember, & 1.unverändert & \textit{unverändert} \\
\hline
2. \newline bis 5. (weggefallen) & 2.unverändert & \textit{unverändert} \\
\hline
6. \newline \S{} 99 Absatz 2 sind zum Zeitpunkt des \newline Endes einer vorläufigen Maßnahme, & 6.unverändert & \textit{unverändert} \\
\hline
7. \newline \S{} 99 Absatz 3 Nummer 1 sind zum \newline Zeitpunkt der rechtskräftigen gerichtlichen \newline Entscheidung über die Annahme als Kind, & 7.unverändert & \textit{unverändert} \\
\hline
8. \newline \S{} 99 Absatz 3 Nummer 2 Buchstabe a, \newline Nummer 3 und Absatz 6a, 6b und 10 sind \newline für das abgelaufene Kalenderjahr, & 8.unverändert & \textit{unverändert} \\
\hline
9. \newline \S{} 99 Absatz 3 Nummer 2 Buchstabe b und \newline Absatz 4 und 5 sind zum 31. Dezember, & 9.unverändert & \textit{unverändert} \\
\hline
10. \newline \S{} 99 Absatz 7, 7a bis 7c sind zum 1. März, & 10.unverändert & \textit{unverändert} \\
\hline
11. \newline \S{} 99 Absatz 6 sind zum Zeitpunkt des \newline Abschlusses der \newline Gefährdungseinschätzung, & 11.unverändert & \textit{unverändert} \\
\hline
12. \newline \S{} 99 Absatz 8 sind für das abgelaufene \newline Kalenderjahr, & 12.unverändert & \textit{unverändert} \\
\hline
13. \newline \S{} 99 Absatz 9 sind zum 15. Dezember. zu \newline erteilen. & 13.unverändert & \textit{unverändert} \\
\hline
\multicolumn{3}{|c|}{\cellcolor{gray!10}\textbf{\S{} 102 -- Auskunftspflicht}} \\
\hline
(1) Für die Erhebungen besteht \newline Auskunftspflicht. Die Angaben zu \S{} 100 \newline Nummer 4 sind freiwillig. & \textit{(1) unverändert} & \textit{unverändert} \\
\hline
(2) Auskunftspflichtig sind & \textit{(2) unverändert} & \textit{unverändert} \\
\hline
1. \newline die örtlichen Träger der Jugendhilfe für die \newline Erhebungen nach \S{} 99 Absatz 1 bis 10, \newline nach Absatz 8 nur, soweit eigene Angebote \newline gemacht wurden, & 1.unverändert & \textit{unverändert} \\
\hline
2. \newline die überörtlichen Träger der Jugendhilfe für \newline die Erhebungen nach \S{} 99 Absatz 3 und 7 \newline und 8 bis 10, nach Absatz 8 nur, soweit \newline eigene Angebote gemacht wurden, & 2. \newline die überörtlichen Träger der Jugendhilfe für \newline die Erhebungen nach \S{} 99 Absatz \textbf{1,} 3 und \newline 7 und 8 bis 10, nach Absatz 8 nur, soweit \newline eigene Angebote gemacht wurden, & \textit{unverändert} \\
\hline
3. \newline die obersten Landesjugendbehörden für \newline die Erhebungen nach \S{} 99 Absatz 7 und 8 \newline bis 10, & 3.unverändert & \textit{unverändert} \\
\hline
4. \newline die fachlich zuständige oberste \newline Bundesbehörde für die Erhebung nach \S{} \newline 99 Absatz 10, & 4.unverändert & \textit{unverändert} \\
\hline
5. \newline die kreisangehörigen Gemeinden und die \newline Gemeindeverbände, soweit sie Aufgaben \newline der Jugendhilfe wahrnehmen, für die \newline Erhebungen nach \S{} 99 Absatz 7 bis 10, & 5.unverändert & \textit{unverändert} \\
\hline
6. \newline die Träger der freien Jugendhilfe für \newline Erhebungen nach \S{} 99 Absatz 1, soweit sie \newline eine Beratung nach \S{} 28 oder \S{} 41 \newline betreffen, nach \S{} 99 Absatz 8, soweit sie \newline anerkannte Träger der freien Jugendhilfe \newline nach \S{} 75 Absatz 1 oder Absatz 3 sind, \newline und nach \S{} 99 Absatz 3, 7 und 9, & 6.unverändert & \textit{unverändert} \\
\hline
7. \newline Adoptionsvermittlungsstellen nach \S{} 2 \newline Absatz 3 des \newline Adoptionsvermittlungsgesetzes aufgrund \newline ihrer Tätigkeit nach \S{} 1 des \newline Adoptionsvermittlungsgesetzes sowie \newline anerkannte Auslandsvermittlungsstellen \newline nach \S{} 4 Absatz 2 Satz 3 des \newline Adoptionsvermittlungsgesetzes aufgrund \newline ihrer Tätigkeit nach \S{} 2a Absatz 4 Nummer \newline 2 des Adoptionsvermittlungsgesetzes \newline gemäß \S{} 99 Absatz 3 Nummer 1 sowie \newline gemäß \S{} 99 Absatz 3 Nummer 2a für die \newline Zahl der ausgesprochenen Annahmen und \newline gemäß \S{} 99 Absatz 3 Nummer 2b für die \newline Zahl der vorgemerkten Adoptionsbewerber, & 7.unverändert & \textit{unverändert} \\
\hline
8. \newline die Leiter der Einrichtungen, Behörden und \newline Geschäftsstellen in der Jugendhilfe für die \newline Erhebungen nach \S{} 99 Absatz 7. Die \newline Auskunftspflichtigen für Erhebungen nach \newline \S{} 99 Absatz 7c werden durch Landesrecht \newline bestimmt. & 8.unverändert & \textit{unverändert} \\
\hline
(3) Zur Durchführung der Erhebungen nach \newline \S{} 99 Absatz 1, 3, 7, 8 und 9 übermitteln die \newline Träger der öffentlichen Jugendhilfe den \newline statistischen Ämtern der Länder auf \newline Anforderung die erforderlichen Anschriften \newline der übrigen Auskunftspflichtigen. & \textit{(3) unverändert} & \textit{unverändert} \\
\hline
\end{longtable}
\newpage

\begin{longtable}{|L{8.6cm}|L{8.6cm}|L{8.6cm}|}
\hline
\multicolumn{3}{|c|}{\cellcolor{gray!20}\textbf{Artikel 2 (2026): Änderung des Neunten Buches Sozialgesetzbuch}} \\
\hline
\textbf{Geltendes Recht} & \textbf{Änderungen RefE 2024} & \textbf{Änderungen RefE 2026} \\
\hline
\endhead

\multicolumn{3}{|c|}{\cellcolor{gray!10}\textbf{\S{} 21 -- Besondere Anforderungen an das Teilhabeplanverfahren}} \\
\hline
Ist der Träger der Eingliederungshilfe der \newline für die Durchführung des \newline Teilhabeplanverfahrens verantwortliche \newline Rehabilitationsträger, gelten für ihn die \newline Vorschriften für die Gesamtplanung \newline ergänzend; dabei ist das \newline Gesamtplanverfahren ein Gegenstand des \newline Teilhabeplanverfahrens. Ist der Träger der \newline öffentlichen Jugendhilfe der für die \newline Durchführung des Teilhabeplans \newline verantwortliche Rehabilitationsträger, \newline gelten für ihn die Vorschriften für den \newline Hilfeplan nach den \S{}\S{} 36, 36b und 37c des \newline Achten Buches ergänzend. Ist der Träger \newline der Sozialen Entschädigung der für die \newline Durchführung des Teilhabeplanverfahrens \newline verantwortliche Rehabilitationsträger, \newline gelten für ihn die Vorschriften für das \newline Fallmanagement nach \S{} 30 des \newline Vierzehnten Buches ergänzend. & Ist der Träger der Eingliederungshilfe der \newline für die Durchführung des \newline Teilhabeplanverfahrens verantwortliche \newline Rehabilitationsträger, gelten für ihn die \newline Vorschriften für die Gesamtplanung \newline ergänzend; dabei ist das \newline Gesamtplanverfahren ein Gegenstand des \newline Teilhabeplanverfahrens. Ist der Träger der \newline öffentlichen Jugendhilfe der für die \newline Durchführung des Teilhabeplans \newline verantwortliche Rehabilitationsträger, \newline gelten für ihn die Vorschriften für den \newline Hilfeplan nach den \textbf{\S{}\S{} 36 bis 36b, 37a und \newline 38c} des Achten Buches ergänzend. Ist der \newline Träger der Sozialen Entschädigung der für \newline die Durchführung des \newline Teilhabeplanverfahrens verantwortliche \newline Rehabilitationsträger, gelten für ihn die \newline Vorschriften für das Fallmanagement nach \newline \S{} 30 des Vierzehnten Buches ergänzend. & Ist der Träger der Eingliederungshilfe der \newline für die Durchführung des \newline Teilhabeplanverfahrens verantwortliche \newline Rehabilitationsträger, gelten für ihn die \newline Vorschriften für die Gesamtplanung \newline ergänzend; dabei ist das \newline Gesamtplanverfahren ein Gegenstand des \newline Teilhabeplanverfahrens. Ist der Träger der \newline öffentlichen Jugendhilfe der für die \newline Durchführung des Teilhabeplans \newline verantwortliche Rehabilitationsträger, \newline gelten für ihn die Vorschriften für den \newline Hilfeplan nach den \textbf{\S{}\S{} 36 bis 36b, 37a und \newline 38c} des Achten Buches ergänzend. Ist der \newline Träger der Sozialen Entschädigung der für \newline die Durchführung des \newline Teilhabeplanverfahrens verantwortliche \newline Rehabilitationsträger, gelten für ihn die \newline Vorschriften für das Fallmanagement nach \newline \S{} 30 des Vierzehnten Buches ergänzend. \\
\hline
\multicolumn{3}{|c|}{\cellcolor{gray!10}\textbf{\S{} 63 -- Zuständigkeit nach den Leistungsgesetzen}} \\
\hline
(1) Die Leistungen im Eingangsverfahren \newline und im Berufsbildungsbereich einer \newline anerkannten Werkstatt für behinderte \newline Menschen erbringen & \textit{(1) unverändert} & \textit{(1) unverändert} \\
\hline
1. \newline die Bundesagentur für Arbeit, soweit nicht \newline einer der in den Nummern 2 bis 4 \newline genannten Träger zuständig ist, & 1. unverändert & 1. unverändert \\
\hline
2. \newline die Träger der Unfallversicherung im \newline Rahmen ihrer Zuständigkeit für durch \newline Arbeitsunfälle Verletzte und von \newline Berufskrankheiten Betroffene, & 2. unverändert & 2. unverändert \\
\hline
3. \newline die Träger der Rentenversicherung unter \newline den Voraussetzungen der \S{}\S{} 11 bis 13 des \newline Sechsten Buches und & 3. unverändert & 3. unverändert \\
\hline
4. \newline die Träger der Sozialen Entschädigung \newline unter den Voraussetzungen der \S{}\S{} 63 und \newline 64 des Vierzehnten Buches. & 4. unverändert & 4. unverändert \\
\hline
(2) Die Leistungen im Arbeitsbereich einer \newline anerkannten Werkstatt für behinderte \newline Menschen erbringen & \textit{(2) unverändert} & \textit{(2) unverändert} \\
\hline
1. \newline die Träger der Unfallversicherung im \newline Rahmen ihrer Zuständigkeit für durch \newline Arbeitsunfälle Verletzte und von \newline Berufskrankheiten Betroffene, & 1. unverändert & 1. unverändert \\
\hline
2. \newline die Träger der Sozialen Entschädigung \newline unter den Voraussetzungen des \S{} 63 des \newline Vierzehnten Buches, & 2. unverändert & 2. unverändert \\
\hline
3. \newline die Träger der öffentlichen Jugendhilfe \newline unter den Voraussetzungen des \S{} 35a des \newline Achten Buches und & 3. \newline die Träger der öffentlichen Jugendhilfe \newline unter den Voraussetzungen des \textbf{\S{}\S{} 27, 35a} \newline des Achten Buches und & 3. \newline die Träger der öffentlichen Jugendhilfe \newline unter den Voraussetzungen des \textbf{\S{}\S{} 27, 35a} \newline des Achten Buches und \\
\hline
4. \newline im Übrigen die Träger der \newline Eingliederungshilfe unter den \newline Voraussetzungen des \S{} 99. & 4. unverändert & 4. unverändert \\
\hline
(3) Absatz 1 gilt auch für die Leistungen zur \newline beruflichen Bildung bei einem anderen \newline Leistungsanbieter sowie für die Leistung \newline des Budgets für Ausbildung an Menschen \newline mit Behinderungen, die Anspruch auf \newline Leistungen nach \S{} 57 haben. Absatz 2 gilt \newline auch für die Leistungen zur Beschäftigung \newline bei einem anderen Leistungsanbieter, für \newline die Leistung des Budgets für Ausbildung an \newline Menschen mit Behinderungen, die \newline Anspruch auf Leistungen nach \S{} 58 haben \newline und die keinen Anspruch auf Leistungen \newline nach \S{} 57 haben, sowie für die Leistung \newline des Budgets für Arbeit. & \textit{(3) unverändert} & \textit{(3) unverändert} \\
\hline
\multicolumn{3}{|c|}{\cellcolor{gray!10}\textbf{\S{} 85 -- Klagerecht der Verbände}} \\
\hline
Werden Menschen mit Behinderungen in \newline ihren Rechten nach diesem Buch verletzt, \newline können an ihrer Stelle und mit ihrem \newline Einverständnis Verbände klagen, die nach \newline ihrer Satzung Menschen mit \newline Behinderungen auf Bundes- oder \newline Landesebene vertreten und nicht selbst am \newline Prozess beteiligt sind. In diesem Fall \newline müssen alle Verfahrensvoraussetzungen \newline wie bei einem Rechtsschutzersuchen durch \newline den Menschen mit Behinderungen selbst \newline vorliegen. & \textit{unverändert} & Werden Menschen mit Behinderungen in \newline ihren Rechten nach diesem Buch \textbf{oder \newline nach dem zweiten Kapitel, Vierter \newline Abschnitt, Dritter Unterabschnitt des \newline Achten Buches }verletzt, können an ihrer \newline Stelle und mit ihrem Einverständnis \newline Verbände klagen, die nach ihrer Satzung \newline Menschen mit Behinderungen auf Bundes- \newline oder Landesebene vertreten und nicht \newline selbst am Prozess beteiligt sind. In diesem \newline Fall müssen alle \newline Verfahrensvoraussetzungen wie bei einem \newline Rechtsschutzersuchen durch den \newline Menschen mit Behinderungen selbst \newline vorliegen. \\
\hline
\multicolumn{3}{|c|}{\cellcolor{gray!10}\textbf{\S{} 98 -- Örtliche Zuständigkeit}} \\
\hline
(1) Für die Eingliederungshilfe örtlich \newline zuständig ist der Träger der \newline Eingliederungshilfe, in dessen Bereich die \newline leistungsberechtigte Person ihren \newline gewöhnlichen Aufenthalt zum Zeitpunkt der \newline ersten Antragstellung nach \S{} 108 Absatz 1 \newline hat oder in den zwei Monaten vor den \newline Leistungen einer Betreuung über Tag und \newline Nacht zuletzt gehabt hatte. Bedarf es nach \newline \S{} 108 Absatz 2 keines Antrags, ist der \newline Beginn des Verfahrens nach Kapitel 7 \newline maßgeblich. Diese Zuständigkeit bleibt bis \newline zur Beendigung des Leistungsbezuges \newline bestehen. Sie ist neu festzustellen, wenn \newline für einen zusammenhängenden Zeitraum \newline von mindestens sechs Monaten keine \newline Leistungen bezogen wurden. Eine \newline Unterbrechung des Leistungsbezuges \newline wegen stationärer \newline Krankenhausbehandlung oder \newline medizinischer Rehabilitation gilt nicht als \newline Beendigung des Leistungsbezuges. & \textit{(1) unverändert} & \textit{(1) unverändert} \\
\hline
(2) Steht innerhalb von vier Wochen nicht \newline fest, ob und wo der gewöhnliche Aufenthalt \newline begründet worden ist, oder ist ein \newline gewöhnlicher Aufenthalt nicht vorhanden \newline oder nicht zu ermitteln, hat der für den \newline tatsächlichen Aufenthalt zuständige Träger \newline der Eingliederungshilfe über die Leistung \newline unverzüglich zu entscheiden und sie \newline vorläufig zu erbringen. Steht der \newline gewöhnliche Aufenthalt in den Fällen des \newline Satzes 1 fest, wird der Träger der \newline Eingliederungshilfe nach Absatz 1 örtlich \newline zuständig und hat dem nach Satz 1 \newline leistenden Träger die Kosten zu erstatten. \newline Ist ein gewöhnlicher Aufenthalt im \newline Bundesgebiet nicht vorhanden oder nicht \newline zu ermitteln, ist der Träger der \newline Eingliederungshilfe örtlich zuständig, in \newline dessen Bereich sich die \newline leistungsberechtigte Person tatsächlich \newline aufhält. & \textit{(2) unverändert} & \textit{(2) unverändert} \\
\hline
(3) Werden für ein Kind vom Zeitpunkt der \newline Geburt an Leistungen nach diesem Teil \newline des Buches über Tag und Nacht beantragt, \newline tritt an die Stelle seines gewöhnlichen \newline Aufenthalts der gewöhnliche Aufenthalt der \newline Mutter. & (3) entfällt & (3) entfällt \\
\hline
(4) Als gewöhnlicher Aufenthalt im Sinne \newline dieser Vorschrift gilt nicht der stationäre \newline Aufenthalt oder der auf richterlich \newline angeordneter Freiheitsentziehung \newline beruhende Aufenthalt in einer \newline Vollzugsanstalt. In diesen Fällen ist der \newline Träger der Eingliederungshilfe örtlich \newline zuständig, in dessen Bereich die \newline leistungsberechtigte Person ihren \newline gewöhnlichen Aufenthalt in den letzten \newline zwei Monaten vor der Aufnahme zuletzt \newline hatte. & \textit{(4) unverändert} & \textit{(4) unverändert} \\
\hline
(5) Bei Personen, die am 31. Dezember \newline 2019 Leistungen nach dem Sechsten \newline Kapitel des Zwölften Buches in der am 31. \newline Dezember 2019 geltenden Fassung \newline bezogen haben und auch ab dem 1. \newline Januar 2020 Leistungen nach Teil 2 dieses \newline Buches erhalten, ist der Träger der \newline Eingliederungshilfe örtlich zuständig, \newline dessen örtliche Zuständigkeit sich am 1. \newline Januar 2020 im Einzelfall in \newline entsprechender Anwendung von \S{} 98 \newline Absatz 1 Satz 1 oder Absatz 5 des \newline Zwölften Buches oder in entsprechender \newline Anwendung von \S{} 98 Absatz 2 Satz 1 und \newline 2 in Verbindung mit \S{} 107 des Zwölften \newline Buches ergeben würde. Absatz 1 Satz 3 \newline bis 5 gilt entsprechend. Im Übrigen bleiben \newline die Absätze 2 bis 4 unberührt. & \textit{(5) unverändert} & \textit{(5) unverändert} \\
\hline
\multicolumn{3}{|c|}{\cellcolor{gray!10}\textbf{\S{} 117 -- Gesamtplanverfahren}} \\
\hline
Absätze 1 bis 5 & \textit{unverändert} & \textit{unverändert} \\
\hline
(6) Bei minderjährigen \newline Leistungsberechtigten wird der nach \S{} 86 \newline des Achten Buches zuständige örtliche \newline Träger der öffentlichen Jugendhilfe vom \newline Träger der Eingliederungshilfe mit \newline Zustimmung des \newline Personensorgeberechtigten informiert und \newline nimmt am Gesamtplanverfahren beratend \newline teil, soweit dies zur Feststellung der \newline Leistungen der Eingliederungshilfe nach \newline den Kapiteln 3 bis 6 erforderlich ist. \newline Hiervon kann in begründeten \newline Ausnahmefällen abgesehen werden, \newline insbesondere, wenn durch die Teilnahme \newline des zuständigen örtlichen Trägers der \newline öffentlichen Jugendhilfe das \newline Gesamtplanverfahren verzögert würde. & \textit{unverändert} & (6) entfällt \\
\hline
\multicolumn{3}{|c|}{\cellcolor{gray!10}\textbf{\S{} 119 -- Gesamtplankonferenz}} \\
\hline
(1) Mit Zustimmung des \newline Leistungsberechtigten kann der Träger der \newline Eingliederungshilfe eine \newline Gesamtplankonferenz durchführen, um die \newline Leistungen für den Leistungsberechtigten \newline nach den Kapiteln 3 bis 6 sicherzustellen. \newline Die Leistungsberechtigten, die beteiligten \newline Rehabilitationsträger und bei \newline minderjährigen Leistungsberechtigten der \newline nach \S{} 86 des Achten Buches zuständige \newline örtliche Träger der öffentlichen Jugendhilfe \newline können dem nach \S{} 15 verantwortlichen \newline Träger der Eingliederungshilfe die \newline Durchführung einer Gesamtplankonferenz \newline vorschlagen. Den Vorschlag auf \newline Durchführung einer Gesamtplankonferenz \newline kann der Träger der Eingliederungshilfe \newline ablehnen, wenn der maßgebliche \newline Sachverhalt schriftlich ermittelt werden \newline kann oder der Aufwand zur Durchführung \newline nicht in einem angemessenen Verhältnis \newline zum Umfang der beantragten Leistung \newline steht. & \textit{unverändert} & (1) Mit Zustimmung des \newline Leistungsberechtigten kann der Träger der \newline Eingliederungshilfe eine \newline Gesamtplankonferenz durchführen, um die \newline Leistungen für den Leistungsberechtigten \newline nach den Kapiteln 3 bis 6 sicherzustellen. \newline Die Leistungsberechtigten \textbf{und die \newline beteiligten Rehabilitationsträger können \newline dem nach \S{} 15 verantwortlichen Träger \newline der Eingliederungshilfe die \newline Durchführung einer \newline Gesamtplankonferenz vorschlagen. }Den \newline Vorschlag auf Durchführung einer \newline Gesamtplankonferenz kann der Träger der \newline Eingliederungshilfe ablehnen, wenn der \newline maßgebliche Sachverhalt schriftlich \newline ermittelt werden kann oder der Aufwand \newline zur Durchführung nicht in einem \newline angemessenen Verhältnis zum Umfang der \newline beantragten Leistung steht. \\
\hline
Absätze 2 bis 4 & \textit{unverändert} & \textit{unverändert} \\
\hline
\multicolumn{3}{|c|}{\cellcolor{gray!10}\textbf{\S{} 134 -- Sonderregelung zum Inhalt der Vereinbarungen zur Erbringung von Leistungen für minderjährige Leistungsberechtigte und in Sonderfällen}} \\
\hline
(1) In der schriftlichen Vereinbarung zur \newline Erbringung von Leistungen für \newline minderjährige Leistungsberechtigte \newline zwischen dem Träger der \newline Eingliederungshilfe und dem \newline Leistungserbringer sind zu regeln: \newline 1. \newline Inhalt, Umfang und Qualität einschließlich \newline der Wirksamkeit der Leistungen \newline (Leistungsvereinbarung) sowie \newline 2. \newline die Vergütung der Leistung \newline (Vergütungsvereinbarung). & \textit{unverändert} & (1) entfällt \\
\hline
(2) In die Leistungsvereinbarung sind als \newline wesentliche Leistungsmerkmale \newline insbesondere aufzunehmen: \newline 1. \newline die betriebsnotwendigen Anlagen des \newline Leistungserbringers, \newline 2. \newline der zu betreuende Personenkreis, \newline 3. \newline Art, Ziel und Qualität der Leistung, \newline 4. \newline die Festlegung der personellen \newline Ausstattung, \newline 5. \newline die Qualifikation des Personals sowie \newline 6. \newline die erforderliche sächliche Ausstattung. & \textit{unverändert} & (2) entfällt \\
\hline
(3) Die Vergütungsvereinbarung besteht \newline mindestens aus \newline 1. \newline der Grundpauschale für Unterkunft und \newline Verpflegung, \newline 2. \newline der Maßnahmepauschale sowie \newline 3. \newline einem Betrag für betriebsnotwendige \newline Anlagen einschließlich ihrer Ausstattung \newline (Investitionsbetrag). \newline Förderungen aus öffentlichen Mitteln sind \newline anzurechnen. Die Maßnahmepauschale ist \newline nach Gruppen für Leistungsberechtigte mit \newline vergleichbarem Bedarf zu kalkulieren. & \textit{unverändert} & (3) entfällt \\
\hline
(4) Die Absätze 1 bis 3 finden auch \newline Anwendung, wenn volljährige \newline Leistungsberechtigte Leistungen zur \newline Schulbildung nach \S{} 112 Absatz 1 Nummer \newline 1 sowie Leistungen zur schulischen \newline Ausbildung für einen Beruf nach \S{} 112 \newline Absatz 1 Nummer 2 erhalten, soweit diese \newline Leistungen in besonderen \newline Ausbildungsstätten über Tag und Nacht für \newline Menschen mit Behinderungen erbracht \newline werden. Entsprechendes gilt bei anderen \newline volljährigen Leistungsberechtigten, wenn \newline 1. \newline das Konzept des Leistungserbringers auf \newline Minderjährige als zu betreuenden \newline Personenkreis ausgerichtet ist, \newline 2. \newline der Leistungsberechtigte von diesem \newline Leistungserbringer bereits Leistungen über \newline Tag und Nacht auf Grundlage von \newline Vereinbarungen nach den Absätzen 1 bis \newline 3, \S{} 78b des Achten Buches, \S{} 75 Absatz 3 \newline des Zwölften Buches in der am 31. \newline Dezember 2019 geltenden Fassung oder \newline nach Maßgabe des \S{} 75 Absatz 4 des \newline Zwölften Buches in der am 31. Dezember \newline 2019 geltenden Fassung erhalten hat und \newline 3. \newline der Leistungsberechtigte nach Erreichen \newline der Volljährigkeit für eine kurze Zeit, in der \newline Regel nicht länger als bis zur Vollendung \newline des 21. Lebensjahres, Leistungen von \newline diesem Leistungserbringer weitererhält, mit \newline denen insbesondere vor dem Erreichen der \newline Volljährigkeit definierte Teilhabeziele \newline erreicht werden sollen. & \textit{unverändert} & (4) entfällt \\
\hline
\multicolumn{3}{|c|}{\cellcolor{gray!10}\textbf{\S{} 136 -- Beitrag aus Einkommen zu den Aufwendungen}} \\
\hline
(1) Bei den Leistungen nach diesem Teil ist \newline ein Beitrag zu den Aufwendungen \newline aufzubringen, wenn das Einkommen im \newline Sinne des \S{} 135 der antragstellenden \newline Person sowie bei minderjährigen Personen \newline der im Haushalt lebenden Eltern oder des \newline im Haushalt lebenden Elternteils die \newline Beträge nach Absatz 2 übersteigt. & \textit{unverändert} & (1) Bei den Leistungen nach diesem Teil ist \newline ein Beitrag zu den Aufwendungen \newline aufzubringen, wenn das Einkommen im \newline Sinne des \S{} 135 der antragstellenden \newline Person sowie bei minderjährigen Personen \newline der im Haushalt lebenden Eltern oder des \newline im Haushalt lebenden Elternteils die \newline Beträge nach Absatz 2 übersteigt. \\
\hline
Absätze 2 bis 5 & \textit{unverändert} & \textit{unverändert} \\
\hline
\multicolumn{3}{|c|}{\cellcolor{gray!10}\textbf{\S{} 138 -- Besondere Höhe des Beitrages zu den Aufwendungen}} \\
\hline
(1) Ein Beitrag ist nicht aufzubringen bei \newline 1. \newline heilpädagogischen Leistungen nach \S{} 113 \newline Absatz 2 Nummer 3, \newline 2. \newline Leistungen zur medizinischen \newline Rehabilitation nach \S{} 109, \newline 3. \newline Leistungen zur Teilhabe am Arbeitsleben \newline nach \S{} 111 Absatz 1, \newline 4. \newline Leistungen zur Teilhabe an Bildung nach \S{} \newline 112 Absatz 1 Nummer 1, \newline 5. \newline Leistungen zur schulischen oder \newline hochschulischen Ausbildung oder \newline Weiterbildung für einen Beruf nach \S{} 112 \newline Absatz 1 Nummer 2, soweit diese \newline Leistungen in besonderen \newline Ausbildungsstätten über Tag und Nacht für \newline Menschen mit Behinderungen erbracht \newline werden, \newline 6. \newline Leistungen zum Erwerb und Erhalt \newline praktischer Kenntnisse und Fähigkeiten \newline nach \S{} 113 Absatz 2 Nummer 5, soweit \newline diese der Vorbereitung auf die Teilhabe am \newline Arbeitsleben nach \S{} 111 Absatz 1 dienen, \newline 7. \newline Leistungen nach \S{} 113 Absatz 1, die noch \newline nicht eingeschulten leistungsberechtigten \newline Personen die für sie erreichbare Teilnahme \newline am Leben in der Gemeinschaft \newline ermöglichen sollen, \newline 8. \newline gleichzeitiger Gewährung von Leistungen \newline zum Lebensunterhalt nach dem Zweiten \newline oder Zwölften Buch oder nach \S{} 27a des \newline Bundesversorgungsgesetzes. & \textit{unverändert} & (1) Ein Beitrag ist nicht aufzubringen bei \newline 1. \newline heilpädagogischen Leistungen nach \S{} 113 \newline Absatz 2 Nummer 3, \newline 2. \newline Leistungen zur medizinischen \newline Rehabilitation nach \S{} 109, \newline 3. \newline Leistungen zur Teilhabe am Arbeitsleben \newline nach \S{} 111 Absatz 1, \newline 4. \newline Leistungen zur Teilhabe an Bildung nach \S{} \newline 112 Absatz 1 Nummer 1, \newline 5. \newline Leistungen zur schulischen oder \newline hochschulischen Ausbildung oder \newline Weiterbildung für einen Beruf nach \S{} 112 \newline Absatz 1 Nummer 2, soweit diese \newline Leistungen in besonderen \newline Ausbildungsstätten über Tag und Nacht für \newline Menschen mit Behinderungen erbracht \newline werden, \newline 6. \newline Leistungen zum Erwerb und Erhalt \newline praktischer Kenntnisse und Fähigkeiten \newline nach \S{} 113 Absatz 2 Nummer 5, soweit \newline diese der Vorbereitung auf die Teilhabe am \newline Arbeitsleben nach \S{} 111 Absatz 1 dienen, \newline 7. \newline Leistungen nach \S{} 113 Absatz 1, die noch \newline nicht eingeschulten leistungsberechtigten \newline Personen die für sie erreichbare Teilnahme \newline am Leben in der Gemeinschaft \newline ermöglichen sollen, \newline \textbf{7}. \newline gleichzeitiger Gewährung von Leistungen \newline zum Lebensunterhalt nach dem Zweiten \newline oder Zwölften Buch oder nach \S{} 27a des \newline Bundesversorgungsgesetzes. \\
\hline
(2) Wenn ein Beitrag nach \S{} 137 \newline aufzubringen ist, ist für weitere Leistungen \newline im gleichen Zeitraum oder weitere \newline Leistungen an minderjährige Kinder im \newline gleichen Haushalt nach diesem Teil kein \newline weiterer Beitrag aufzubringen. & \textit{unverändert} & (2) Wenn ein Beitrag nach \S{} 137 \newline aufzubringen ist, ist für weitere Leistungen \newline im gleichen Zeitraum oder weitere \newline Leistungen an minderjährige Kinder im \newline gleichen Haushalt nach diesem Teil kein \newline weiterer Beitrag aufzubringen. \\
\hline
\multicolumn{3}{|c|}{\cellcolor{gray!10}\textbf{\S{} 140 -- Einsatz des Vermögens}} \\
\hline
(1) Die antragstellende Person sowie bei \newline minderjährigen Personen die im Haushalt \newline lebenden Eltern oder ein Elternteil haben \newline vor der Inanspruchnahme von Leistungen \newline nach diesem Teil die erforderlichen Mittel \newline aus ihrem Vermögen aufzubringen. & \textit{unverändert} & \textbf{(1) Die antragsstellende Person hat vor \newline der Inanspruchnahme von Leistungen \newline nach diesem Teil die erforderlichen \newline Mittel aus ihrem Vermögen \newline aufzubringen. } \\
\hline
Absätze 2 und 3 & \textit{unverändert} & \textit{unverändert} \\
\hline
\multicolumn{3}{|c|}{\cellcolor{gray!10}\textbf{\S{} 142 -- Sonderregelungen für minderjährige Leistungsberechtigte und in Sonderfällen}} \\
\hline
(1) Minderjährigen Leistungsberechtigten \newline und ihren Eltern oder einem Elternteil ist \newline bei Leistungen im Sinne des \S{} 138 Absatz \newline 1 Nummer 1, 2, 4, 5 und 7 die Aufbringung \newline der Mittel für die Kosten des \newline Lebensunterhalts nur in Höhe der für den \newline häuslichen Lebensunterhalt ersparten \newline Aufwendungen zuzumuten, soweit \newline Leistungen über Tag und Nacht oder über \newline Tag erbracht werden. & \textit{unverändert} & (1) entfällt \\
\hline
(2) Sind Leistungen von einem oder \newline mehreren Anbietern über Tag und Nacht \newline oder über Tag oder für ärztliche oder \newline ärztlich verordnete Maßnahmen \newline erforderlich, sind die Leistungen, die der \newline Vereinbarung nach \S{} 134 Absatz 3 \newline zugrunde liegen, durch den Träger der \newline Eingliederungshilfe auch dann in vollem \newline Umfang zu erbringen, wenn den \newline minderjährigen Leistungsberechtigten und \newline ihren Eltern oder einem Elternteil die \newline Aufbringung der Mittel nach Absatz 1 zu \newline einem Teil zuzumuten ist. In Höhe dieses \newline Teils haben sie zu den Kosten der \newline erbrachten Leistungen beizutragen; \newline mehrere Verpflichtete haften als \newline Gesamtschuldner. & \textit{unverändert} & (2) entfällt \\
\hline
(3) Die Absätze 1 und 2 gelten \newline entsprechend für volljährige \newline Leistungsberechtigte, wenn diese \newline Leistungen erhalten, denen \newline Vereinbarungen nach \S{} 134 Absatz 4 \newline zugrunde liegen. In diesem Fall ist den \newline volljährigen Leistungsberechtigten die \newline Aufbringung der Mittel für die Kosten des \newline Lebensunterhalts nur in Höhe der für ihren \newline häuslichen Lebensunterhalt ersparten \newline Aufwendungen zuzumuten. & \textit{unverändert} & (3) entfällt \\
\hline
\end{longtable}
\newpage

\begin{longtable}{|L{8.6cm}|L{8.6cm}|L{8.6cm}|}
\hline
\multicolumn{3}{|c|}{\cellcolor{gray!20}\textbf{Artikel 3 (2026) / Artikel 4 (2024): Änderung des Ersten Buch Sozialgesetzbuch}} \\
\hline
\textbf{Geltendes Recht} & \textbf{Änderungen RefE 2024} & \textbf{Änderungen RefE 2026} \\
\hline
\endhead

\multicolumn{3}{|c|}{\cellcolor{gray!10}\textbf{\S{} 27 -- Leistungen der Kinder- und Jugendhilfe}} \\
\hline
(1) Nach dem Recht der Kinder- und \newline Jugendhilfe können in Anspruch \newline genommen werden: & \textit{(1) unverändert} & \textit{(1) unverändert} \\
\hline
1. \newline Angebote der Jugendarbeit, der \newline Jugendsozialarbeit und des erzieherischen \newline Jugendschutzes, & 1. unverändert & 1. unverändert \\
\hline
2. \newline Angebote zur Förderung der Erziehung in \newline der Familie, & 2. unverändert & 2. unverändert \\
\hline
3. Angebote zur Förderung von Kindern in \newline Tageseinrichtungen und in Tagespflege, & 3. unverändert & 3. unverändert \\
\hline
4. Hilfe zur Erziehung, Eingliederungshilfe \newline für seelisch behinderte Kinder und \newline Jugendliche sowie Hilfe für junge \newline Volljährige. & 4. Hilfe zur Erziehung, Eingliederungshilfe \newline für seelisch behinderte Kinder und \newline Jugendliche \textbf{mit Behinderungen} sowie \newline Hilfe für junge Volljährige. & 4. \textbf{Leistungen zur Entwicklung, zur \newline Erziehung und zur Teilhabe, Hilfen für \newline junge Volljährige.} \\
\hline
(2) Zuständig sind die Kreise und die \newline kreisfreien Städte, nach Maßgabe des \newline Landesrechts auch kreisangehörige \newline Gemeinden; sie arbeiten mit der freien \newline Jugendhilfe zusammen. & \textit{(2) unverändert} & \textit{(2) unverändert} \\
\hline
\end{longtable}
\newpage

\begin{longtable}{|L{8.6cm}|L{8.6cm}|L{8.6cm}|}
\hline
\multicolumn{3}{|c|}{\cellcolor{gray!20}\textbf{Artikel 4 (2026) / Artikel 5 (2024): Änderung des Zweiten Buches Sozialgesetzbuch}} \\
\hline
\textbf{Geltendes Recht} & \textbf{Änderungen RefE 2024} & \textbf{Änderungen RefE 2026} \\
\hline
\endhead

\multicolumn{3}{|c|}{\cellcolor{gray!10}\textbf{\S{} 11a -- Nicht zu berücksichtigendes Einkommen}} \\
\hline
(1) Nicht als Einkommen zu \newline berücksichtigen sind & \textit{(1) unverändert} & \textit{(1) unverändert} \\
\hline
1. \newline Leistungen nach diesem Buch, & 1. unverändert & 1. unverändert \\
\hline
2. \newline (weggefallen) & 2. unverändert & 2. unverändert \\
\hline
3. \newline die Renten oder Beihilfen, die nach dem \newline Bundesentschädigungsgesetz für Schaden \newline an Leben sowie an Körper oder \newline Gesundheit erbracht werden, bis zur Höhe \newline der vergleichbaren \newline Entschädigungszahlungen nach Kapitel 9 \newline des Vierzehnten Buches, & 3. unverändert & 3. unverändert \\
\hline
4. \newline Aufwandspauschalen nach \S{} 1878 des \newline Bürgerlichen Gesetzbuchs kalenderjährlich \newline bis zu dem in \S{} 3 Nummer 26 Satz 1 des \newline Einkommensteuergesetzes genannten \newline Betrag, & 4. unverändert & 4. unverändert \\
\hline
5. \newline Aufwandsentschädigungen oder \newline Einnahmen aus nebenberuflichen \newline Tätigkeiten, die nach \S{} 3 Nummer 12, \newline Nummer 26 oder Nummer 26a des \newline Einkommensteuergesetzes steuerfrei sind, \newline soweit diese Einnahmen einen Betrag in \newline Höhe von 3 000 Euro im Kalenderjahr nicht \newline überschreiten, & 5. unverändert & 5. unverändert \\
\hline
6. \newline Mutterschaftsgeld nach \S{} 19 des \newline Mutterschutzgesetzes & 6. unverändert & 6. unverändert \\
\hline
7. \newline Einmalige Einnahmen aus Erbschaften, \newline Vermächtnissen und \newline Pflichtteilszuwendungen. & 7. unverändert & 7. unverändert \\
\hline
(2) Entschädigungen, die wegen eines \newline Schadens, der kein Vermögensschaden ist, \newline nach \S{} 253 Absatz 2 des Bürgerlichen \newline Gesetzbuchs geleistet werden, sind nicht \newline als Einkommen zu berücksichtigen. & \textit{(2) unverändert} & \textit{(2) unverändert} \\
\hline
(3) Leistungen, die aufgrund öffentlich- \newline rechtlicher Vorschriften zu einem \newline ausdrücklich genannten Zweck erbracht \newline werden, sind nur so weit als Einkommen zu \newline berücksichtigen, als die Leistungen nach \newline diesem Buch im Einzelfall demselben \newline Zweck dienen. Abweichend von Satz 1 sind \newline als Einkommen zu berücksichtigen \newline a) \newline für das dritte Pflegekind zu 75 Prozent, & \textit{(3) unverändert} & \textit{(3) unverändert} \\
\hline
1. \newline die Leistungen nach \S{} 39 des Achten \newline Buches, die für den erzieherischen Einsatz \newline erbracht werden, & 1. \newline die Leistungen nach \S{} \textbf{39c} des Achten \newline Buches, die für den erzieherischen Einsatz \newline erbracht werden, & 1. \newline die Leistungen nach \S{} \textbf{39c} des Achten \newline Buches, die für den erzieherischen Einsatz \newline erbracht werden, \\
\hline
b) \newline für das vierte und jedes weitere Pflegekind \newline vollständig, & b) unverändert & b) unverändert \\
\hline
2. \newline die Leistungen nach \S{} 23 des Achten \newline Buches, & 2. unverändert & 2. unverändert \\
\hline
3. \newline die Leistungen der Ausbildungsförderung \newline nach dem \newline Bundesausbildungsförderungsgesetz sowie \newline vergleichbare Leistungen der \newline Begabtenförderungswerke; \S{} 14b Absatz 2 \newline Satz 1 des \newline Bundesausbildungsförderungsgesetzes \newline bleibt unberührt, & 3. unverändert & 3. unverändert \\
\hline
4. \newline die Berufsausbildungsbeihilfe nach dem \newline Dritten Buch mit Ausnahme der Bedarfe \newline nach \S{} 64 Absatz 3 Satz 1 des Dritten \newline Buches sowie & 4. unverändert & 4. unverändert \\
\hline
5. \newline Reisekosten zur Teilhabe am Arbeitsleben \newline nach \S{} 127 Absatz 1 Satz 1 des Dritten \newline Buches in Verbindung mit \S{} 73 des \newline Neunten Buches. & 5. unverändert & 5. unverändert \\
\hline
(4) Zuwendungen der freien \newline Wohlfahrtspflege sind nicht als Einkommen \newline zu berücksichtigen, soweit sie die Lage der \newline Empfängerinnen und Empfänger nicht so \newline günstig beeinflussen, dass daneben \newline Leistungen nach diesem Buch nicht \newline gerechtfertigt wären. & \textit{(4) unverändert} & \textit{(4) unverändert} \\
\hline
(5) Zuwendungen, die ein anderer erbringt, \newline ohne hierzu eine rechtliche oder sittliche \newline Pflicht zu haben, sind nicht als Einkommen \newline zu berücksichtigen, soweit & \textit{(5) unverändert} & \textit{(5) unverändert} \\
\hline
1. \newline ihre Berücksichtigung für die \newline Leistungsberechtigten grob unbillig wäre \newline oder & 1. unverändert & 1. unverändert \\
\hline
2. \newline sie die Lage der Leistungsberechtigten \newline nicht so günstig beeinflussen, dass \newline daneben Leistungen nach diesem Buch \newline nicht gerechtfertigt wären. & 2. unverändert & 2. unverändert \\
\hline
(6) Überbrückungsgeld nach \S{} 51 des \newline Strafvollzugsgesetzes oder vergleichbare \newline Leistungen nach landesrechtlichen \newline Regelungen sind nicht als Einkommen zu \newline berücksichtigen. & \textit{(6) unverändert} & \textit{(6) unverändert} \\
\hline
(7) Nicht als Einkommen zu \newline berücksichtigen sind Einnahmen von \newline Schülerinnen und Schülern allgemein- oder \newline berufsbildender Schulen, die das 25. \newline Lebensjahr noch nicht vollendet haben, aus \newline Erwerbstätigkeiten, die in den Schulferien \newline ausgeübt werden. Satz 1 gilt nicht für eine \newline Ausbildungsvergütung, auf die eine \newline Schülerin oder ein Schüler einen Anspruch \newline hat. & \textit{(7) unverändert} & \textit{(7) unverändert} \\
\hline
\end{longtable}
\newpage

\begin{longtable}{|L{8.6cm}|L{8.6cm}|L{8.6cm}|}
\hline
\multicolumn{3}{|c|}{\cellcolor{gray!20}\textbf{Artikel 5 (2026) / Artikel 6 (2024): Änderung des Fünften Buches Sozialgesetzbuch}} \\
\hline
\textbf{Geltendes Recht} & \textbf{Änderungen RefE 2024} & \textbf{Änderungen RefE 2026} \\
\hline
\endhead

\multicolumn{3}{|c|}{\cellcolor{gray!10}\textbf{\S{} 44b -- Krankengeld für eine bei stationärer Behandlung mitaufgenommene Begleitperson aus dem engsten persönlichen Umfeld}} \\
\hline
(1) Ab dem 1. November 2022 haben \newline Versicherte Anspruch auf Krankengeld, \newline wenn sie & \textit{(1) unverändert} & \textit{(1) unverändert} \\
\hline
1. \newline zur Begleitung eines Versicherten bei einer \newline stationären Krankenhausbehandlung nach \newline \S{} 39 mitaufgenommen werden, & 1. unverändert & 1. unverändert \\
\hline
Der Anspruch besteht für die Dauer der \newline Mitaufnahme. Der Mitaufnahme steht die \newline ganztägige Begleitung gleich. & \textit{unverändert} & \textit{unverändert} \\
\hline
(2) Der Gemeinsame Bundesausschuss \newline bestimmt in einer Richtlinie nach \S{} 92 bis \newline zum 1. August 2022 Kriterien zur \newline Abgrenzung des Personenkreises, der die \newline Begleitung aus medizinischen Gründen \newline benötigt. Vor der Entscheidung ist den für \newline die Wahrnehmung der Interessen von \newline Menschen mit Behinderungen \newline maßgeblichen Organisationen, der \newline Bundesarbeitsgemeinschaft der \newline überörtlichen Träger der Sozialhilfe und der \newline Eingliederungshilfe sowie der \newline Bundesarbeitsgemeinschaft der \newline Landesjugendämter Gelegenheit zur \newline Stellungnahme zu geben. Die \newline Stellungnahmen sind in die Entscheidung \newline einzubeziehen. & \textit{(2) unverändert} & \textit{(2) unverändert} \\
\hline
(3) Der Anspruch auf Krankengeld bei \newline Erkrankung des Kindes nach \S{} 45 Absatz 1 \newline bleibt unberührt. & \textit{(3) unverändert} & \textit{(3) unverändert} \\
\hline
(4) \S{} 45 Absatz 3 gilt entsprechend. Den \newline Anspruch auf unbezahlte Freistellung von \newline der Arbeitsleistung haben auch \newline Arbeitnehmer, die nicht Versicherte mit \newline Anspruch auf Krankengeld nach Absatz 1 \newline sind. & \textit{(4) unverändert} & \textit{(4) unverändert} \\
\hline
\end{longtable}
\newpage

\begin{longtable}{|L{8.6cm}|L{8.6cm}|L{8.6cm}|}
\hline
\multicolumn{3}{|c|}{\cellcolor{gray!20}\textbf{Artikel 6 (2026) / Artikel 7 (2024): Änderung des Vierzehnten Buches Sozialgesetzbuch}} \\
\hline
\textbf{Geltendes Recht} & \textbf{Änderungen RefE 2024} & \textbf{Änderungen RefE 2026} \\
\hline
\endhead

\multicolumn{3}{|c|}{\cellcolor{gray!10}\textbf{\S{} 65 -- Leistungen zur Teilhabe an Bildung}} \\
\hline
Geschädigte, die auf Grund der \newline Schädigungsfolgen zum \newline leistungsberechtigten Personenkreis im \newline Sinne von \S{} 99 des Neunten Buches \newline gehören, erhalten Leistungen zur Teilhabe \newline an Bildung entsprechend Teil 2 Kapitel 5 \newline des Neunten Buches. & \textit{unverändert} & Geschädigte, die auf Grund der \newline Schädigungsfolgen zum \newline leistungsberechtigten Personenkreis im \newline Sinne von \S{} 99 des Neunten Buches \newline gehören, erhalten Leistungen zur Teilhabe \newline an Bildung entsprechend Teil 2 Kapitel 5 \newline des Neunten Buches. \textbf{Geschädigte Kinder \newline oder Jugendliche, die aufgrund der \newline Schädigungsfolgen zum \newline leistungsberechtigten Personenkreis im \newline Sinne von \S{} 27 Absatz 3 des Achten \newline Buches gehören, erhalten Leistungen \newline zur Teilhabe an Bildung entsprechend \S{} \newline 35d des Achten Buches.} \\
\hline
\multicolumn{3}{|c|}{\cellcolor{gray!10}\textbf{\S{} 66 -- Leistungen zur Sozialen Teilhabe}} \\
\hline
(1) Geschädigte, die auf Grund der \newline Schädigungsfolgen zum \newline leistungsberechtigten Personenkreis im \newline Sinne von \S{} 99 des Neunten Buches \newline gehören, erhalten Leistungen zur Sozialen \newline Teilhabe entsprechend Teil 2 Kapitel 6 des \newline Neunten Buches. & \textit{unverändert} & (1) Geschädigte, die auf Grund der \newline Schädigungsfolgen zum \newline leistungsberechtigten Personenkreis im \newline Sinne von \S{} 99 des Neunten Buches \newline gehören, erhalten Leistungen zur Sozialen \newline Teilhabe entsprechend Teil 2 Kapitel 6 des \newline Neunten Buches. \textbf{Geschädigte Kinder \newline oder Jugendliche, die aufgrund der \newline Schädigungsfolgen zum \newline leistungsberechtigten Personenkreis im \newline Sinne von \S{} 27 Absatz 3 des Achten \newline Buches gehören, erhalten Leistungen \newline zur Sozialen Teilhabe entsprechend \S{}\S{} \newline 35f, 35h und 35i des Achten Buches.} \\
\hline
(2) Abweichend von Absatz 1 werden \newline Leistungen zur Mobilität nach \S{} 83 des \newline Neunten Buches erbracht. Sie umfassen \newline zudem Leistungen zum Betrieb, Unterhalt, \newline Unterstellen und Abstellen eines \newline Kraftfahrzeuges. & \textit{unverändert} & \textit{(2) unverändert} \\
\hline
\multicolumn{3}{|c|}{\cellcolor{gray!10}\textbf{\S{} 93 -- Leistungen zum Lebensunterhalt}} \\
\hline
(1) Geschädigte erhalten Leistungen zum \newline Lebensunterhalt. Hinterbliebene erhalten \newline Leistungen nach Satz 1 für einen Zeitraum \newline von bis zu fünf Jahren nach dem Tod der \newline oder des Geschädigten. Die Vorschriften \newline des Dritten und Vierten Kapitels des \newline Zwölften Buches gelten entsprechend unter \newline Berücksichtigung der besonderen Lage der \newline Geschädigten und Hinterbliebenen. \newline Leistungen zum Lebensunterhalt werden \newline nur erbracht, soweit der Lebensunterhalt \newline nicht aus den übrigen Leistungen nach \newline diesem Gesetz bestritten werden kann. & \textit{(1) unverändert} & \textit{(1) unverändert} \\
\hline
(2) Sind für Geschädigte und Waisen \newline Leistungen zum Lebensunterhalt während \newline der Erbringung von Leistungen nach dem \newline Achten Buch erforderlich, erbringt diese der \newline Träger der Sozialen Entschädigung nach \newline Maßgabe des Absatzes 1, soweit nicht der \newline Träger der öffentlichen Jugendhilfe \newline Leistungen nach \S{} 39 des Achten Buches \newline erbringt. & (2) Sind für Geschädigte und Waisen \newline Leistungen zum Lebensunterhalt während \newline der Erbringung von Leistungen nach dem \newline Achten Buch erforderlich, erbringt diese der \newline Träger der Sozialen Entschädigung nach \newline Maßgabe des Absatzes 1, soweit nicht der \newline Träger der öffentlichen Jugendhilfe \newline Leistungen nach \S{} \textbf{39c} des Achten Buches \newline erbringt. & (2) Sind für Geschädigte und Waisen \newline Leistungen zum Lebensunterhalt während \newline der Erbringung von Leistungen nach dem \newline Achten Buch erforderlich, erbringt diese der \newline Träger der Sozialen Entschädigung nach \newline Maßgabe des Absatzes 1, soweit nicht der \newline Träger der öffentlichen Jugendhilfe \newline Leistungen nach \S{} \textbf{39c} des Achten Buches \newline erbringt. \\
\hline
(3) Ansprüche nach dem \newline Bundesausbildungsförderungsgesetz \newline gehen Ansprüchen nach diesem Buch vor. \newline Soweit für Geschädigte weitere Leistungen \newline zum Lebensunterhalt während der \newline Erbringung von Leistungen zur \newline Ausbildungsförderung nach dem \newline Bundesausbildungsförderungsgesetz \newline erforderlich sind, erbringt diese der Träger \newline der Sozialen Entschädigung nach \newline Maßgabe des Absatzes 1. & \textit{(3) unverändert} & \textit{(3) unverändert} \\
\hline
(4) Der Anspruch auf Leistungen zum \newline Lebensunterhalt kann nicht abgetreten, \newline übertragen, verpfändet oder gepfändet \newline werden. & \textit{(4) unverändert} & \textit{(4) unverändert} \\
\hline
\end{longtable}
\newpage

\begin{longtable}{|L{8.6cm}|L{8.6cm}|L{8.6cm}|}
\hline
\multicolumn{3}{|c|}{\cellcolor{gray!20}\textbf{Artikel 7 (2026): Änderung des Jugendschutzgesetzes}} \\
\hline
\textbf{Geltendes Recht} & \textbf{Änderungen RefE 2024} & \textbf{Änderungen RefE 2026} \\
\hline
\endhead

\multicolumn{3}{|c|}{\cellcolor{gray!10}\textbf{\S{} 9 -- Absatz 1}} \\
\hline
(2) Absatz 1 Nummer 1 gilt nicht, wenn \newline Jugendliche von einer \newline personensorgeberechtigten Person \newline begleitet werden. & \textit{unverändert} & (2) entfällt \\
\hline
Absätze 3 bis 4 & \textit{unverändert} & Absätze 2 und 3 \\
\hline
\end{longtable}
\newpage

\begin{longtable}{|L{8.6cm}|L{8.6cm}|L{8.6cm}|}
\hline
\multicolumn{3}{|c|}{\cellcolor{gray!20}\textbf{Artikel 3 (2024): Änderung des Sozialgerichtsgesetz}} \\
\hline
\textbf{Geltendes Recht} & \textbf{Änderungen RefE 2024} & \textbf{Änderungen RefE 2026} \\
\hline
\endhead

\multicolumn{3}{|c|}{\cellcolor{gray!10}\textbf{\S{} 51 -- Rechtsweg und Zuständigkeit}} \\
\hline
(1) Die Gerichte der Sozialgerichtsbarkeit \newline entscheiden über öffentlich-rechtliche \newline Streitigkeiten & \textit{(1) unverändert} & \textit{unverändert} \\
\hline
1. \newline in Angelegenheiten der gesetzlichen \newline Rentenversicherung einschließlich der \newline Alterssicherung der Landwirte, & 1. unverändert & \textit{unverändert} \\
\hline
2. \newline in Angelegenheiten der gesetzlichen \newline Krankenversicherung, der sozialen \newline Pflegeversicherung und der privaten \newline Pflegeversicherung (Elftes Buch \newline Sozialgesetzbuch), auch soweit durch \newline diese Angelegenheiten Dritte betroffen \newline werden; dies gilt nicht für Streitigkeiten in \newline Angelegenheiten nach \S{} 110 des Fünften \newline Buches Sozialgesetzbuch aufgrund einer \newline Kündigung von Versorgungsverträgen, die \newline für Hochschulkliniken oder \newline Plankrankenhäuser (\S{} 108 Nr. 1 und 2 des \newline Fünften Buches Sozialgesetzbuch) gelten, & 2. unverändert & \textit{unverändert} \\
\hline
3. \newline in Angelegenheiten der gesetzlichen \newline Unfallversicherung mit Ausnahme der \newline Streitigkeiten aufgrund der Überwachung \newline der Maßnahmen zur Prävention durch die \newline Träger der gesetzlichen \newline Unfallversicherung, & 3. unverändert & \textit{unverändert} \\
\hline
4. \newline in Angelegenheiten der Arbeitsförderung \newline einschließlich der übrigen Aufgaben der \newline Bundesagentur für Arbeit, & 4. unverändert & \textit{unverändert} \\
\hline
5. \newline in sonstigen Angelegenheiten der \newline Sozialversicherung, & 5. unverändert & \textit{unverändert} \\
\hline
6. \newline in Angelegenheiten des Sozialen \newline Entschädigungsrechts, & 6. unverändert & \textit{unverändert} \\
\hline
6a. \newline in Angelegenheiten der Sozialhilfe \newline einschließlich der Angelegenheiten nach \newline Teil 2 des Neunten Buches \newline Sozialgesetzbuch und des \newline Asylbewerberleistungsgesetzes, & 6a. unverändert & \textit{unverändert} \\
\hline
 & \textbf{6b. \newline In Angelegenheiten der Kinder- und \newline Jugendhilfe, soweit sie Leistungen der \newline Eingliederungshilfe betreffen, } & \textit{unverändert} \\
\hline
7. \newline bei der Feststellung von Behinderungen \newline und ihrem Grad sowie weiterer \newline gesundheitlicher Merkmale, ferner der \newline Ausstellung, Verlängerung, Berichtigung \newline und Einziehung von Ausweisen nach \S{} 152 \newline des Neunten Buches Sozialgesetzbuch, & 7. unverändert & \textit{unverändert} \\
\hline
8. \newline die aufgrund des \newline Aufwendungsausgleichsgesetzes \newline entstehen, & 8. unverändert & \textit{unverändert} \\
\hline
9. \newline (weggefallen) & 9. unverändert & \textit{unverändert} \\
\hline
10. \newline für die durch Gesetz der Rechtsweg vor \newline diesen Gerichten eröffnet wird. & 10. unverändert & \textit{unverändert} \\
\hline
(2) Die Gerichte der Sozialgerichtsbarkeit \newline entscheiden auch über privatrechtliche \newline Streitigkeiten in Angelegenheiten der \newline Zulassung von Trägern und Maßnahmen \newline durch fachkundige Stellen nach dem \newline Fünften Kapitel des Dritten Buches \newline Sozialgesetzbuch und in Angelegenheiten \newline der gesetzlichen Krankenversicherung, \newline auch soweit durch diese Angelegenheiten \newline Dritte betroffen werden. Satz 1 gilt für die \newline soziale Pflegeversicherung und die private \newline Pflegeversicherung (Elftes Buch \newline Sozialgesetzbuch) entsprechend. & \textit{(2) unverändert} & \textit{unverändert} \\
\hline
(3) Von der Zuständigkeit der Gerichte der \newline Sozialgerichtsbarkeit nach den Absätzen 1 \newline und 2 ausgenommen sind Streitigkeiten in \newline Verfahren nach dem Gesetz gegen \newline Wettbewerbsbeschränkungen, die \newline Rechtsbeziehungen nach \S{} 69 des Fünften \newline Buches Sozialgesetzbuch betreffen & \textit{(3) unverändert} & \textit{unverändert} \\
\hline
\end{longtable}
\newpage

\end{document}